% concatenated from arxiv.tex
\documentclass[11pt,letterpaper,reqno]{amsart}
\usepackage[foot]{amsaddr}
\usepackage[margin=1in]{geometry}
\usepackage{amsmath,amsthm,amsfonts,amssymb,mathrsfs}
\usepackage{hyperref,color,enumerate,cleveref}
\usepackage{graphicx}
%\usepackage{refcheck}
\usepackage{bm}

\makeatletter
\def\paragraph{\@startsection{paragraph}{4}%
  \z@\z@{-\fontdimen2\font}%
  {\normalfont\bfseries}}
\makeatother

\allowdisplaybreaks

\newtheorem{thm}{Theorem}[section]
\newtheorem{conj}[thm]{Conjecture}
\newtheorem{cor}[thm]{Corollary}
\newtheorem{lem}[thm]{Lemma}
\newtheorem{prop}[thm]{Proposition}

\theoremstyle{definition}
\newtheorem{asp}[thm]{Assumption}
\newtheorem{exa}[thm]{Example}
\newtheorem{defi}[thm]{Definition}
\newtheorem{remark}[thm]{Remark}
\newtheorem{cla}[thm]{Claim}

% Basic symbols/operators
\def\eps{\varepsilon}
\def\d{\mathrm{d}}
\def\D{\mathrm{D}}
\def\P{\mathbb{P}}
\def\E{\mathbb{E}}
\def\Z{\mathbb{Z}}
\def\R{\mathbb{R}}
\def\Tr{\operatorname{Tr}}
\def\Pois{\operatorname{Poisson}}
\def\Law{\operatorname{Law}}
\def\GP{\operatorname{GP}}
\def\PP{\operatorname{PP}}
\def\Id{\mathrm{Id}}
\def\op{\mathrm{op}}
\def\vec{\operatorname{vec}}
\def\diag{\operatorname{diag}}

% Fonts
\def\cR{\mathcal{R}}
\def\cS{\mathcal{S}}
\def\cD{\mathcal{D}}
\def\cF{\mathcal{F}}
\def\cN{\mathcal{N}}
\def\sF{\mathrm{F}}
\def\a{\mathbf{a}}
\def\b{\mathbf{b}}
\def\e{\mathbf{e}}
\def\v{\mathbf{v}}
\def\x{\mathbf{x}}
\def\y{\mathbf{y}}
\def\z{\mathbf{z}}
\def\A{\mathbf{A}}
\def\B{\mathbf{B}}
\def\bD{\mathbf{D}}
\def\bE{\mathbf{E}}
\def\F{\mathbf{F}}
\def\H{\mathbf{H}}
\def\I{\mathbf{I}}
\def\J{\mathbf{J}}
\def\M{\mathbf{M}}
\def\bN{\mathbf{N}}
\def\bP{\mathbf{P}}
\def\U{\mathbf{U}}
\def\V{\mathbf{V}}
\def\W{\mathbf{W}}
\def\X{\mathbf{X}}
\def\Y{\mathbf{Y}}
\def\bZ{\mathbf{Z}}
\def\1{\mathbf{1}}
\def\beps{\boldsymbol{\varepsilon}}
\def\btheta{\boldsymbol{\theta}}
\def\bxi{\boldsymbol{\xi}}
\def\bnu{\boldsymbol{\nu}}
\def\bPi{\boldsymbol{\Pi}}

% Shorthands
\def\set{\mathscr{S}_\kappa}
\def\Ct{C_\theta}
\def\Rt{R_\theta}
\def\Cf{C_f}
\def\Rf{R_f}
\def\pCt{\Phi_{\Ct}}
\def\pCf{\Phi_{\Cf}}
\def\pRt{\Phi_{\Rt}}
\def\pRf{\Phi_{\Rf}}
\def\pRfs{\Phi_{\Rf^*}}
\def\trsfrm{\mathcal{T}}
\def\trsfrmA{\mathcal{T}_{\mathcal{S}_\xi \to {\mathcal{S}}_\theta}}
\def\trsfrmB{\mathcal{T}_{{\mathcal{S}}_\theta \to \mathcal{S}_\xi}}
\def\trsfrmD{\mathcal{T}^\delta}
\def\trsfrmDA{\mathcal{T}^\delta_{\mathcal{D}_\xi \to {\mathcal{D}}_\theta}}
\def\trsfrmDB{\mathcal{T}^\delta_{{\mathcal{D}}_\theta \to \mathcal{D}_\xi}}
\def\nI{(\mathrm{I})}
\def\nII{(\mathrm{II})}
\def\nIII{(\mathrm{III})}
\def\nIV{(\mathrm{IV})}
\def\nV{(\mathrm{V})}
\def\nVI{(\mathrm{VI})}
\newcommand{\lbddst}[2]{\mathsf{dist}_{\lambda,T}(#1,#2)}
\newcommand{\flr}[1]{{\lfloor {#1} \rfloor}}
\newcommand{\cil}[1]{{\lceil {#1} \rceil}}

% Norms and brackets
\newcommand{\norm}[1]{\|#1\|}
\newcommand{\normop}[1]{\|#1\|_{\mathrm{op}}}
\newcommand{\prn}[1]{\left({#1}\right)} % parentheses
\newcommand{\brk}[1]{\left[{#1}\right]} % bracket
\newcommand{\abs}[1]{\left|{#1}\right|} % norm

% Comments
\definecolor{lw}{RGB}{11,0,249}
\newcommand{\lw}[1]{\textcolor{lw}{[LW: #1]}}
\newcommand{\zf}[1]{{\color{red}[ZF: #1]}}

\title[High-dimensional learning dynamics of multi-pass SGD]
{High-dimensional learning dynamics of multi-pass Stochastic Gradient
Descent in multi-index models}

\author{Zhou Fan}
\author{Leda Wang}
\email{zhou.fan@yale.edu, leda.wang@yale.edu}
\address{Department of Statistics and Data Science, Yale University}

\begin{document}

\begin{abstract}
We study the learning dynamics of a multi-pass, mini-batch Stochastic
Gradient Descent (SGD) procedure for empirical risk minimization in
high-dimensional multi-index models with isotropic random data. In an
asymptotic regime where the sample size $n$ and data dimension $d$ increase
proportionally, for any sub-linear batch size $\kappa \asymp n^\alpha$ where
$\alpha \in [0,1)$, and for a commensurate ``critical'' scaling of the
learning rate, we provide an asymptotically exact characterization of
the coordinate-wise dynamics of SGD. This characterization takes
the form of a system of dynamical mean-field equations, driven by a scalar
Poisson jump process that represents the asymptotic limit of SGD sampling
noise. We develop an analogous characterization of the Stochastic Modified
Equation (SME) which provides a Gaussian diffusion approximation to SGD.

Our analyses imply that the limiting dynamics for SGD are the same for
any batch size scaling $\alpha \in [0,1)$, and that under a commensurate scaling of the learning rate,
dynamics of SGD, SME, and gradient flow
are mutually distinct, with those of SGD and SME coinciding in the special 
case of a linear model. We recover a known dynamical mean-field characterization
of gradient flow in a limit of small learning rate, and of
one-pass/online SGD in a limit of increasing sample size $n/d \to \infty$.
\end{abstract}


\maketitle

%\tableofcontents
%\vspace{-5ex}

\section{Introduction}

Introduced by Robbins and Monro in \cite{robbins1951stochastic}, Stochastic Gradient
Descent (SGD) and its variants have played an important role in machine
learning, and are often the optimization methods of choice in large-scale
learning systems \cite{bottou2010large,bottou2018optimization}. Existing theory for SGD explains many aspects
of its behavior, including optimization theory on convergence rates for
convex optimization problems \cite{nemirovski2009robust,moulines2011non,rakhlin2012making,bach2013non}, classical formulations of
ODE scaling limits under small learning rate asymptotics
\cite{ljung1992stochastic,kushner2003stochastic}, and
analyses of diffusion approximations, especially near critical points and
minimizers of the loss landscape
\cite{mandt2017stochastic,li2017stochastic,feng2018semigroups,hu2019diffusion,blanc2020implicit,li2022happens}.
Such results are suggestive of a behavior that is similar to (non-stochastic)
gradient descent and gradient flow,
with additional diffusivity properties near criticality.
However, in many modern examples of learning systems --- often
characterized by a high-dimensional optimization parameter, a complex and
highly non-convex loss landscape, and a combination of small batch size and
large learning rate --- it is also empirically observed that SGD may have
markedly different properties from gradient flow, possibly leading to improved
generalization \cite{keskar2017large,jastrzkebski2017three}. Our understanding
of SGD in such contexts remains incomplete.

In this work, we will study the dynamics of SGD in a prototypical
high-dimensional application of empirical risk minimization for a multi-index
model, with isotropic random data $\x \in \R^d$ and labels
\[y=\sigma^*(\x^\top \btheta^*,\eps), \qquad \btheta^* \in \R^{d \times k^*}.\]
Given $n$ such data observations $(\x_1,y_1),\ldots,(\x_n,y_n)$, we will
consider a standard mini-batch, multi-pass SGD procedure for optimizing a
(possibly non-convex) empirical risk with coordinate-separable regularizer,
\begin{equation}\label{eq:ERMintro}
\cR(\btheta)=\sum_{i=1}^n L(\sigma(\x_i^\top \btheta),y_i)+\sum_{j=1}^d
G(\theta_j), \qquad \btheta=[\theta_1,\ldots,\theta_d]^\top \in \R^{d \times
k}.
\end{equation}
Our main result, building upon the method of \cite{celentano2021high}, will
provide an asymptotically exact characterization of the
learning dynamics of SGD in this setting, with a general sub-linear batch
size $\kappa \ll n$, and under a high-dimensional asymptotic limit
as $n,d \to \infty$ proportionally.

Our analysis is motivated and inspired by three lines of related literature:

\begin{enumerate}[1.]
\item A large body of work, starting with
\cite{saad1995dynamics,saad1995exact,riegler1995line,biehl1995learning} and the mathematical
formalizations of such findings in \cite{wang2019solvable,goldt2019dynamics},
has investigated these learning dynamics in an
analogous \emph{one-pass} or \emph{online learning} setting, where a new
sample $(\x_i,y_i)$ is used in each SGD iteration.

The resulting dynamics are characterized in the high-dimensional limit by a system of ODEs for the overlap parameters
$d^{-1}(\btheta,\btheta^*)^\top (\btheta,\btheta^*) \in \R^{(k+k^*) \times
(k+k^*)}$. Diffusive behavior near critical points was further studied in
\cite{ben2022high}; escape from initial conditions and saddle points
over long time horizons in problems with symmetry has been the subject of
intensive investigation in
\cite{ge2015escaping,fang2019sharp,arous2021online,barak2022hidden,tan2023online,abbe2023sgd,damian2023smoothing,damian2025generative,joshi2025learning};
extensions to two-layer neural networks with more general dimension scalings, and connections to dynamics in the ``mean-field'' limit of infinite width
(\cite{mei2018mean,chizat2018global,rotskoff2022trainability}), were discussed
in \cite{veiga2022phase,arnaboldi2023high}; and extensions to non-isotropic
data and adaptive dynamics were obtained
in \cite{collins2024hitting,collins2024high}.

Such results have illuminated how classical ODE and SDE
approximations of SGD may pertain to high-dimensional settings, and
how problem symmetries may affect the sample complexity of learning. They
shed light also on the dynamics of early training in multi-pass SGD, where
initial iterations are akin to online learning without sample re-use.
However, the full learning dynamics of multi-pass SGD 
differs from that of online SGD in later training, and understanding these
differences may be requisite for investigating the optimization of the empirical
--- rather than population --- risk, and associated issues of overfitting
versus generalization.\\

\item Using statistical physics techniques of dynamical mean-field theory (DMFT)
\cite{sompolinsky1981dynamic,sompolinsky1982relaxational,crisanti1993spherical,cugliandolo1993analytical,arous1995large,arous1997symmetric}, a burgeoning body of literature has developed asymptotically exact
characterizations of multi-pass gradient-based algorithms in a variety
of statistical learning models
\cite{sarao2019afraid,mannelli2019passed,sarao2020complex,sarao2020marvels,sarao2021analytical,liang2023high}, with
substantial recent interest in the training dynamics for various neural network
architectures
\cite{bordelon2022self,bordelon2024dynamical,bordelon2024infinite,montanari2025dynamical,han2025precise,martin2026high}.
Closest to our current analyses is the work of \cite{mignacco2021dynamical,mignacco2021stochasticity,mignacco2022effective} which characterized the
high-dimensional limit of discrete-time SGD dynamics with a linear  batch
size $\kappa \asymp n$ per iteration, as well as a ``persistent SGD'' dynamics
in continuous time that maintains this batch size while replacing samples
individually. Investigations of differences between such dynamics and gradient
flow in a model of phase retrieval were
carried out in \cite{mignacco2021stochasticity}. Differences between the sample complexities of
one-pass and multi-pass methods for learning with
symmetries were also highlighted in \cite{lee2024neural,dandi2024benefits}.

The pioneering work of \cite{celentano2021high} developed a new approach
to mathematically formalize such DMFT characterizations, which previously
were often derived
using non-rigorous techniques. This work \cite{celentano2021high} established on rigorous
grounds the high-dimensional limiting dynamics of gradient flow for a 
class of multi-index models. Several extensions to other
discrete-time and continuous-time first-order optimization and sampling methods have since been developed
\cite{gerbelot2024rigorous,han2025entrywise,fan2025dynamical1,fan2025dynamical2,han2025long,chen2025learning,celentano2025state,dandi2025sequential}, including analyses
of discrete-time SGD with linear batch size $\kappa \asymp n$ in \cite{gerbelot2024rigorous}
and of Glauber dynamics in the
Sherrington-Kirkpatrick model for batch updates of $\kappa \asymp n$ spins in \cite{dandi2025sequential}

We remark that a batch size of $\kappa \asymp n$ implies limited stochasticity
in the gradient approximation. Motivated by a common belief that the
stochasticity of small-batch dynamics may underlie important differences between
SGD and gradient flow \cite{keskar2017large,jastrzkebski2017three}, we will extend the
results of \cite{celentano2021high,gerbelot2024rigorous} to (rigorously) characterize the dynamics of SGD for any batch size
\[\kappa \asymp n^\alpha, \qquad \alpha \in [0,1),\]
under a commensurate scaling of the learning rate.
Our results will elucidate how the stochasticity of the gradient approximation
manifests in the limiting dynamics.\\

\item An interesting line of work
\cite{paquette2021dynamics,paquette2021sgd,lee2022trajectory,paquette2022implicit,marshall2024clip,paquette2025homogenization} obtained high-dimensional
asymptotic characterizations of single-sample (i.e.\ $\alpha=0$ and $\kappa=1$) multi-pass SGD
and several of its variants, in the setting of a linear regression
model with ridge regularization, and possibly non-isotropic data. One finding of
this work, highlighted in \cite{paquette2022implicit,paquette2025homogenization}, is that the asymptotic dynamics of
quadratic observables of $\btheta$ coincide
with those of a simplified SDE approximation (dubbed ``Homogenized
SGD''). Whereas previous analyses establish the accuracy of such
diffusion approximations in the asymptotics of vanishing learning rate
$\eta \to 0$ for fixed dimensions $d$ \cite{li2019stochastic,hu2019diffusion}, the results of \cite{paquette2025homogenization}
showed that such a diffusion approximation remains asymptotically
exact under a commonly studied (large) learning rate scaling for a
high-dimensional linear model, raising a tantalizing question of whether this
holds true in more general high-dimensional learning problems.

Restricting to the setting of isotropic data $\x \in \R^d$, we provide a
negative answer to this question, by developing an analogous characterization
of the limiting dynamics of the Stochastic Modified Equation (SME) that was introduced and studied in \cite{li2017stochastic,li2019stochastic}
as a diffusion approximation for the SGD process.
Our results illustrate that the asymptotic dynamics of SME and SGD are, in
general, different for multi-index models under commonly studied learning
rate scalings in high-dimensional settings, and they coincide only in the case of linear regression.
\end{enumerate}

\subsection{Summary of results}

Our main result provides a mathematical description of the
learning dynamics of a SGD procedure for minimizing the
empirical risk \eqref{eq:ERMintro} over $O(1)$ training epochs, in the
high-dimensional limit as $n,d \to \infty$ proportionally. This will
take the form of a characterization for the exact limiting value of any
coordinate-separable observables,
\[\lim_{n,d \to \infty} \frac{1}{d}\sum_{j=1}^d
\psi(\theta_j^{t_1},\ldots,\theta_j^{t_m},\theta_j^*),
\qquad \lim_{n,d \to \infty} \frac{1}{n}\sum_{i=1}^n
\psi(\x_i^\top \btheta^{t_1},\ldots,\x_i^\top \btheta^{t_m},\x_i^\top
\btheta^*),\]
depending on the latent parameter $\btheta^*=[\theta_1^*,\ldots,\theta_d^*]^\top
\in \R^{d \times k^*}$ and SGD iterates
$\btheta^t=[\theta_1^t,\ldots,\theta_d^t]^\top \in \R^{d \times k}$
at fixed time points $t_1,\ldots,t_m$ in units of training epochs. Such
observables can encompass, for example, the overlap parameters
$d^{-1}(\btheta,\btheta^*)^\top(\btheta,\btheta^*)$ and evaluations of the
training and test losses.

As this characterization involves dynamical mean-field theory and
is somewhat complex, we summarize here a few qualitative aspects/implications
of our results:

\begin{enumerate}[1.]
\item For any sub-linear choice of batch size
\[\kappa \asymp n^\alpha\]
with $\alpha \in [0,1)$, and for a commensurate scaling of the learning rate
$\eta \asymp n^\alpha$,\footnote{Our scaling conventions for the model are $\|\x_i\|_2 \asymp
1$ and $\|\btheta\|_\sF \asymp \sqrt{d}$. Under a common alternative scaling of
$\|\x_i\|_2 \asymp \sqrt{d}$ and $\|\btheta\|_\sF \asymp 1$, the equivalent
learning rate is $\tilde \eta=\eta/d \asymp n^{-1+\alpha}$.} there is a well-defined scaling limit for
the SGD iterates under a time rescaling by $n^{1-\alpha}$,
i.e.\ to time units of training epochs.

Furthermore, this limit is the same for any choice of $\alpha \in [0,1)$ (but
different from the setting of linear batch sizes $\alpha=1$ and $\kappa
\asymp n$), depending only
on $\bar\kappa=\lim_{n,d \to \infty} \kappa/n^\alpha$ and
$\bar \eta=\lim_{n,d \to \infty} \eta/n^\alpha$. This agrees with previous
empirical observations in more complex models
\cite{jastrzkebski2017three,he2019control} that the learning
dynamics of SGD seem to depend moreso on the ratio of learning rate to batch
size than the absolute size of the batch.\\

\item The random sampling of mini-batches in SGD leads to a dynamical
mean-field characterization in which a scalar process $\{\xi^t\}_{t \geq 0}$
that tracks the distributional dynamics of $\{\{\x_i^\top \btheta^t\}_{i \in
[n]}\}_{t \geq 0}$ is
driven by a univariate Poisson jump process $\{z^t\}_{t \geq 0}$.
The presence of this additional Poisson process constitutes the primary
distinction between the asymptotic dynamics of SGD and gradient flow.

For the SME diffusion approximation, the dynamics of $\{\xi^t\}_{t \geq 0}$ are
analogously driven by a univariate Brownian diffusion with mean and covariance 
matching those of $\{z^t\}_{t \geq 0}$.\\

\item In a further scaling limit of small learning rate $\bar \eta \to 0$ (after
taking $n,d \to \infty$ with $\eta \approx n^\alpha \bar \eta$) and
rescaling of time, the dynamics of both
SGD and SME reduce to those of gradient flow. This agrees with the classical
theory on convergence of SGD dynamics to an ODE under small learning rate
asymptotics \cite{kushner2003stochastic}, and clarifies that in this high-dimensional context,
the notion of a ``small'' learning rate under which such an ODE approximation
is accurate may be understood as $\eta \ll 
n^\alpha$. For larger learning rates represented by any
fixed value of $\bar \eta=\lim_{n,d \to \infty} \eta/n^\alpha>0$, the asymptotic dynamics of SGD, SME, and
gradient flow are, in general, all distinct.\\

\item In the special case of linear regression, the
dynamical mean-field equations simplify to deterministic integro-differential
equations for the correlation and response processes, which, in particular,
depend only on the first and second moment statistics of the Poisson
process/Brownian diffusion $\{z^t\}_{t \geq 0}$. Thus, in this case, the
dynamics of quadratic observables of SME and SGD coincide even the setting of
fixed $\bar\eta>0$, in agreement with the results of
\cite{paquette2025homogenization}.\\

\item Setting $\gamma=\lim_{n,d \to \infty} n/d$, in a further scaling limit $\gamma \to \infty$ (after taking $n,d \to \infty$ with $n \approx \gamma d$), the dynamical mean-field equations also simplify,
reducing to a Markov diffusion for a scalar process $\{\theta^t\}_{t \geq
0}$ that tracks the distributional dynamics of $\{\{\theta_j^t\}_{j \in [d]}\}_{t
\geq 0}$. From this simplification, one may recover the known ODE
characterization of the overlap parameters
for one-pass/online SGD \cite{saad1995dynamics,goldt2019dynamics}.\\

\item Since the driving Poisson process $\{z^t\}_{t \geq 0}$ in the
dynamical mean-field limit for SGD is discrete, with $O_\P(1)$ jumps over finite
time horizons (in units of training epochs), this enables rapid numerical simulation of
the preceding process $\{\xi^t\}_{t \geq 0}$. In applications
where $G(\cdot)$ is also a ridge regularizer so that the statistics of the
process $\{\theta^t\}_{t \geq 0}$ are computable analytically, this leads to a
dynamical mean-field limit for SGD that is more amenable to numerical simulation than its
counterparts for gradient flow or SME.\\
\end{enumerate}

We will elaborate upon these discussions further in Section \ref{sec:results},
after formally stating our main results.

\subsection*{Concurrent work}
Concurrent and independent work of \cite{nishiyama2026high} also establishes a DMFT characterization of the high-dimensional limit dynamics for SME, referred to therein as Stochastic Gradient Flow (SGF). Their DMFT equations are equivalent to ours for SME, although formulated in a different manner. \cite{nishiyama2026high} obtains convergence in probability under a weaker assumption of a pseudo-Lipschitz loss, over times $t \in [0,T^*]$ for some $(L,G,\sigma)$-dependent time horizon $T^*<\infty$. Our results establish a stronger notion of a.s.\ convergence over any bounded time horizon $[0,T]$ and pertain also to the SGD dynamics in addition to SME, but we impose a stronger Lipschitz continuity condition for the loss.

\subsection*{Notational conventions} For vectors $\x \in \R^k$ and $\y \in \R^l$,
we write $\x \otimes \y \in \R^{k \times l}$ for the outer product. For a
function $g:\R^k \to \R^k$ and input $\btheta \in \R^{d \times k}$, we 
write $g(\btheta) \in \R^{d \times k}$ for the application of $g(\cdot)$
to each row. We write $\frac{1}{n}\sum_{i=1}^n \delta_{\v_i}$ for the empirical
distribution on $\R^k$ of the rows of
$\V=[\v_1,\ldots,\v_n]^\top \in \R^{n \times k}$.

The notation
$\|\cdot\|$ without subscript will refer to the Euclidean norm $\|\cdot\|_2$ for
vectors and Frobenius norm $\|\cdot\|_\sF$ for matrices. We write
$\|\cdot\|_\op$ for the matrix $\ell_2$-to-$\ell_2$ operator norm, and
$\|\cdot\|_\infty$ for the element-wise $\ell_\infty$ norm. $\Id_k$ is the $k
\times k$ identity matrix, and we omit the subscript when the dimension is
clear. We write $x \lesssim
y$ if $x \leq Cy$ for a constant $C>0$, which does not on the dimensions $n,d$ but will usually depend on the fixed
time horizon $T>0$. We write $x \asymp y$ if
both $x \lesssim y$ and $y \lesssim x$.

\section{Main results}\label{sec:results}

\subsection{Model and dynamics}\label{sec:model}

We consider a standard setting of supervised learning, where
\[(\x_1,y_1),\ldots,(\x_n,y_n) \in \R^d \times \R\]
are $n$ independent training observations. It will be assumed that the data
$\x_i \in \R^d$ has independent coordinates with isotropic covariance, and
that the label $y_i \in \R$ follows a multi-index model, possibly with label
noise: For some parameter $\btheta^* \in \R^{d \times k^*}$ where $k^* \geq 1$,
noise variables
\[\eps_1,\ldots,\eps_n \in \R,\]
and label map $\sigma^*:\R^{k^*} \times \R \to \R$, each label $y_i$
takes the form
\begin{equation}\label{eq:multiindex}
y_i=\sigma^*(\x_i^\top \btheta^*,\eps_i).
\end{equation}

Our main result will characterize the training dynamics of a multi-pass
Stochastic Gradient Descent (SGD) algorithm and its Stochastic Modified Equation
(SME) approximant for minimizing a regularized
empirical risk over $\btheta \in \R^{d \times k}$,
\begin{equation}\label{eq:ERM}
\cR(\btheta)=\sum_{i=1}^n L\big(\sigma(\x_i^\top \btheta),y_i\big)
+\sum_{j=1}^d G(\theta_j).
\end{equation}
Here $\sigma:\R^k \to \R$ is a non-linear activation for some $k \geq 1$,
and $L:\R \times \R \to \R$ and $G:\R^k \to \R$ are a smooth loss function and regularizer. The gradient of $\cR(\btheta)$ takes a form
\begin{equation}\label{eq:riskgradient}
\nabla \cR(\btheta)=\sum_{i=1}^n
\x_i \otimes f(\x_i^\top \btheta,\x_i^\top \btheta^*,\eps_i)
+g(\btheta) \in \R^{d \times k},
\end{equation}
where the functions $f:\R^k \times \R^{k^*} \times \R \to \R^k$ and
$g:\R^k \to \R^k$ are defined by
\begin{equation}\label{eq:fgdef}
f(\xi,w^*,\eps)=L'(\sigma(\xi),\sigma^*(w^*,\eps))
\,\nabla \sigma(\xi),
\qquad g(\theta)=\nabla G(\theta),
\end{equation}
$L'(\cdot,\cdot)$ is the derivative of $L(\cdot,\cdot)$ in its first argument,
and $g(\cdot)$ is applied row-wise to $\btheta \in \R^{d \times k}$.\\

\paragraph{SGD} The SGD dynamics, initialized at
\[\bar\btheta^0 \equiv \btheta^0\in \R^{d \times k},\]
are given by
\begin{equation}\label{eq:SGD}
\bar\btheta^{t+1}=\bar\btheta^t-\eta^t \underbrace{\left(\frac{1}{\kappa}
\sum_{i \in S^t} \x_i \otimes
f(\x_i^\top \bar\btheta^t,\x_i^\top \btheta^*,\eps_i)
+\frac{1}{n}\,g(\bar\btheta^t)\right)}_{:=\v^t},
\qquad S^t \sim \operatorname{Unif}(\set),
\end{equation}
where
\[\set=\{S \subset [n]:|S|=\kappa\}\]
denotes the collection of all subsets of $\{1,\ldots,n\}$ of size $\kappa$,
and $\kappa \geq 1$ is the batch size.
In each iteration, $\eta^t>0$ is a learning rate, and
$S^t \in \set$ is a batch of $\kappa$ samples chosen independently and
uniformly at random. The update $\v^t$ in
\eqref{eq:SGD} has expectation $n^{-1}\nabla_{\btheta} \cR(\bar\btheta^t)$
over this randomness of $S^t$, and constitutes a stochastic approximation for 
this rescaled gradient.

We will study a regime of general sub-linear batch size
\[\kappa \asymp n^\alpha, \qquad \alpha \in [0,1),\]
in a high-dimensional limit where $n,d \to \infty$. (The setting of $\alpha=0$
encompasses a constant batch size independent of $n$, including
single-sample SGD with $\kappa=1$.) Then
$\frac{n}{\kappa} \asymp n^{1-\alpha}$ iterations of the SGD dynamics
\eqref{eq:SGD} constitute one ``epoch'' of training in which, on expectation,
each training sample is used once. We define a continuous-time embedding
$\{\btheta^t\}_{t \geq 0}$ of the SGD dynamics \eqref{eq:SGD}, with time
re-indexed on the scale of epochs, by
\begin{equation}\label{eq:SGDrescaled}
\btheta^t=\bar\btheta^{\lfloor tn^{1-\alpha} \rfloor}.
\end{equation}
Our results will characterize the training dynamics of \eqref{eq:SGDrescaled}
over a fixed number of epochs not depending on $n,d$.

Under our model scaling conventions (to be specified in Assumption
\ref{asp:model}) where
$\|\bar\btheta^t\|_\sF \asymp \sqrt{d}$ and $\|\x_i\|_2 \asymp 1$,
we consider a learning rate that scales also as
\[\eta^t \asymp n^\alpha.\]
One may check that this is the scaling under which both the
mean and variance per coordinate of $\eta^t\v^t$ accumulated over one training
epoch has $O(1)$ size, and thus non-trivial learning may occur over $O(1)$
epochs (in problems without symmetry at initialization). For single-sample SGD where $\alpha=0$ and
$\kappa=1$, this coincides also with learning rate scalings studied previously
in e.g.\ \cite{goldt2019dynamics,tan2023online,paquette2025homogenization}
(in both online and multi-pass settings).

We formalize the above conditions as the following assumption.

\begin{asp}[Batch size and learning rate]\label{asp:SGD}
There exist constants $\alpha \in [0,1)$, $\bar \kappa>0$, and a
Lipschitz-continuous learning rate schedule $\{\bar\eta^t\}_{t \geq 0}$
not depending on $n,d$, such that over each fixed time horizon $T>0$,
\[\lim_{n,d \to \infty} |\kappa/n^\alpha-\bar \kappa|=0, \qquad
\lim_{n,d \to \infty} \sup_{t \in [0,T]}
|\eta^{\lfloor tn^{1-\alpha}\rfloor}/n^\alpha-\bar \eta^t|=0.\]
\end{asp}

\paragraph{SME}

Under Assumption \ref{asp:SGD}, we define the SME approximation
\cite{li2017stochastic} to the above SGD algorithm as the continuous-time
diffusion process
\begin{equation}\label{eq:SME}
\d\btheta^t={-}\bar\eta^t\left(\sum_{i=1}^n \x_i \otimes
f(\x_i^\top \btheta^t,\x_i^\top \btheta^*,\eps_i)
+g(\btheta^t)\right)\d t
+\frac{\bar\eta^t}{\sqrt{\bar \kappa}} \sum_{i=1}^n \Big(\x_i \otimes
f(\x_i^\top \btheta^t,\x_i^\top \btheta^*,\eps_i)\Big) \d b_i^t
\end{equation}
where $\{b_i^t\}_{t \geq 0}$ are independent standard Brownian motions for
$i=1,\ldots,n$. More precisely, this may be understood as a diffusion
approximation to the time-rescaled SGD process $\{\btheta^t\}_{t \geq 0}$ in
\eqref{eq:SGDrescaled}, where the drift
and anisotropic diffusion terms of \eqref{eq:SME} are defined to match the mean and
covariance of the stochastic gradient updates defining \eqref{eq:SGDrescaled}.

(In \cite{paquette2025homogenization}, the authors study a further ``Homogenized
SGD'' approximation where each factor 
$f(\x_i^\top \btheta^t,\x_i^\top \btheta^*,\eps_i)$ in the diffusion term is
further replaced by
$n^{-1}\sum_{i=1}^n f(\x_i^\top \btheta^t,\x_i^\top \btheta^*,\eps_i)$. In this
work, we will study instead the diffusion \eqref{eq:SME} which provides a closer
approximation to SGD.)

\subsection{Model assumptions}

In addition to Assumption \ref{asp:SGD} on the SGD process, we impose the
following two assumptions on the data model and empirical risk.

\begin{asp}[Data model and asymptotic scaling]\label{asp:model}
\phantom{ }
\begin{enumerate}[(a)]
\item As $n,d \to \infty$, $n/d \to \gamma$ for a constant $\gamma>0$, and
$k,k^* \geq 1$ are constants not depending on $n,d$.
\item $(\x_1,y_1),\ldots,(\x_n,y_n) \in \R^d$ are independent samples.
Each data vector $\x_i=(x_{i1},\ldots,x_{id}) \in \R^d$ has independent entries
with
\[\E[x_{ij}]=0, \quad \E[x_{ij}^2]=d^{-1}, \quad \E[|x_{ij}|^p] \leq
C_pd^{-p/2} \text{ for each $p \geq 3$ and a constant $C_p>0$},\]
and $y_i$ is given by the multi-index model \eqref{eq:multiindex}.
\item The initialization $\btheta^0=[\theta_1^0,\ldots,\theta_d^0]^\top
\in \R^{d \times k}$, parameter
$\btheta^*=[\theta_1^*,\ldots,\theta_d^*]^\top \in \R^{d \times k^*}$,
and noise variables $\beps=(\eps_1,\ldots,\eps_n)$ defining
$\y=(y_1,\ldots,y_n)$ via \eqref{eq:multiindex} are non-random and satisfy,
weakly and in Wasserstein-$p$ for each fixed order $p \geq 1$ as $n,d \to
\infty$,
\[\frac{1}{d}\sum_{j=1}^d \delta_{(\theta_j^0,\theta_j^*)} \Rightarrow
\Law(\theta^0,\theta^*),
\qquad \frac{1}{n}\sum_{i=1}^n \delta_{\eps_i} \Rightarrow \Law(\eps).\]
Here $\Law(\theta^0,\theta^*)$ and $\Law(\eps)$ are any probability
distributions on $\R^k \times \R^{k^*}$ and $\R$, respectively, that have
finite moment generating functions in a neighborhood of 0.
\end{enumerate}
\end{asp}

Assumption \ref{asp:model}(c) allows a correlation between $\theta^0$ and
$\theta^*$ in the limiting law, and thus encompasses settings where the SGD/SME
initialization $\btheta^0$ has a non-zero initial overlap with $\btheta^*$.

The assumption that $\btheta^0,\btheta^*,\beps$ are non-random is only for
convenience of analysis, so that randomness of the data arises entirely through
the data matrix
\[\X=[\x_1,\ldots,\x_n]^\top \in \R^{n \times d}.\]
We note that our results will then apply equally in settings where
$(\btheta^0,\btheta^*)$ and/or $\beps$ are random and independent of $\X$, upon
applying these results conditioned on $\btheta^0,\btheta^*,\beps$.

\begin{asp}[Loss and regularizer]\label{asp:lipschitz}
Let $f:\R^k \times \R^{k^*} \times \R \to \R^k$ and $g:\R^k \to \R^k$ be the
functions defined by \eqref{eq:fgdef}. For a constant $C>0$, these satisfy
\[\|f(\xi,w^*,\eps)\|_2 \leq C,
\qquad \|g(\theta)\|_2 \leq C(1+\|\theta\|_2).\]
Furthermore, $f$ and $g$ are twice continuously-differentiable,
with all first and second order partial derivatives uniformly bounded by $C$.
\end{asp}

These conditions hold, for example, when the activation
$\sigma(\cdot)$ and loss $L(\cdot,\cdot)$ defining the empirical risk
\eqref{eq:ERM} are both Lipschitz with
bounded derivatives up to order 3, and the regularizer $G(\cdot)$ is
pseudo-Lipschitz with bounded derivatives of orders 2 and 3
(e.g.\ $G(\theta)=\frac{\lambda}{2}\|\theta\|_2^2$).

We expect that our results may be extendable also to pseudo-Lipschitz loss
functions including the squared loss $L(\hat y,y)=\frac{1}{2}(\hat y-y)^2$,
under a suitable bound for the learning rate $\{\bar\eta^t\}_{t \geq 0}$
and additional technical arguments. To keep the technicalities simpler,
we will not pursue this extension in our current work.

\subsection{Definitions of the DMFT limit processes}\label{sec:DMFTlimits}

The high-dimensional limit of the dynamics $\{\btheta^t\}_{t \geq 0}$ in
\eqref{eq:SGDrescaled} or \eqref{eq:SME} is described via a
$\R^k$-valued process $\{\theta^t\}_{t \geq 0}$ (not depending on $n,d$),
constructed in a probability space of random variables $(\theta^0,\theta^*)$
that are distributed according to the limit law in Assumption
\ref{asp:model}(c). Our results will show a convergence of the empirical
distribution of coordinates of any finite-time marginals of
$(\{\btheta^t\}_{t \geq 0},\btheta^*)$ to the corresponding marginals of the
limit process
$(\{\theta^t\}_{t \geq 0},\theta^*)$,
\[\frac{1}{d}\sum_{j=1}^d
\delta_{(\theta_j^{t_1},\ldots,\theta_j^{t_m},\theta_j^*)}
\Rightarrow \Law(\theta^{t_1},\ldots,\theta^{t_m},\theta^*).\]
Likewise, we will show an analogous convergence of the marginals of $(\{\X\btheta^t\}_{t \geq 0},\X\btheta^*,\beps)$ to those of a limit process $(\{\xi^t\}_{t \geq 0},w^*,\eps)$.

In this section, we define these limit processes $\{\theta^t,\xi^t\}_{t \geq 0}$ for both
SGD and SME, which may be understood as an analogue of the processes described
in \cite{celentano2021high} for gradient flow.

\subsubsection{Limit processes for SGD}\label{sec:DMFT-SGD}
Fix a time horizon $T>0$. We define the high-dimensional limit
for the time-rescaled SGD process $\{\btheta^t\}_{t \in [0,T]}$
of (\ref{eq:SGDrescaled}) via a fixed-point relation
for a system of deterministic correlation and response kernels/operators
$(C_\theta,R_\theta,C_f,R_f,R_f^*,\Gamma)$:
\begin{itemize}
\item $C_\theta$ is the joint covariance kernel of a $\R^k$-valued Gaussian
process $\{w^t\}_{t \in [0,T]}$ and a Gaussian vector $w^* \in \R^{k^*}$.
We denote $C_\theta^{t,s}=\E[w^t \otimes w^s] \in \R^{k \times k}$
for $t,s \in [0,T]$,
$C_\theta^{t,*}=\E[w^t \otimes w^*] \in \R^{k \times k^*}$,
and $C_\theta^{*,*}=\E[w^* \otimes w^*] \in \R^{k^* \times k^*}$.
\item $C_f$ is the covariance kernel of a $\R^k$-valued Gaussian process
$\{u^t\}_{t \in [0,T]}$. We denote likewise
$C_f^{t,s}=\E[u^t \otimes u^s] \in \R^{k \times k}$.
\item $R_\theta \equiv \{R_\theta^{t,s}\}_{t,s \in [0,T]}$ is a $\R^{k \times
k}$-valued process on $[0,T] \times [0,T]$,
where $R_\theta^{t,s}=0$ if $t<s$.
\item $R_f \equiv \{R_f^t\}_{t \in [0,T]}$ is a process of linear operators,
where $R_f^t:L^4([0,t],\R^k) \to \R^k$ for each $t \in [0,T]$.
Given such an operator $R_f^t$ and a matrix-valued process
$x \in L^4([0,t],\R^{k \times k})$, we will write also
$R_f^t(x) \in \R^{k \times k}$ for $R_f^t$ applied column-wise, i.e.\
\begin{equation}\label{eq:matrixRf}
R_f^t(x)=[R_f^t(x_{:,1}),\ldots,R_f^t(x_{:,k})]
\end{equation}
where $x_{:,1},\ldots,x_{:,k} \in L^4([0,t],\R^k)$ are the columns of $x$.
\item $R_f^* \equiv \{R_f^{t,*}\}_{t \in [0,T]}$ is a $\R^{k \times k^*}$-valued
process on $[0,T]$.
\item $\Gamma \equiv \{\Gamma^t\}_{t \in [0,T]}$ is a $\R^{k \times k}$-valued
process on $[0,T]$.
\end{itemize}
We defer a specification of further technical conditions for
$(C_\theta,R_\theta,C_f,R_f,R_f^*,\Gamma)$ to Section
\ref{sec:existenceuniqueness}.

Given the above objects $(C_\theta,R_\theta,C_f,R_f,R_f^*,\Gamma)$, let
\begin{equation}\label{eq:thetaeps}
(\theta^0,\theta^*) \in \R^k \times \R^{k^*}, \qquad \eps \in \R
\end{equation}
be distributed according to the limit laws of Assumption \ref{asp:model}(c),
let
\begin{equation}\label{eq:GPs}
\{u^t\}_{t \in [0,T]} \sim \GP(0,C_f),
\quad (\{w^t\}_{t \in [0,T]},w^*) \sim \GP(0,C_\theta)
\end{equation}
be two centered Gaussian processes with covariance kernels $C_f,C_\theta$, and let
\begin{equation}\label{eq:poisson-z}
\{z^t\}_{t \in [0,T]} \sim \PP(\bar\kappa)
\end{equation}
be a $\{0,1,2,\ldots\}$-valued homogeneous Poisson jump process
with rate $\bar\kappa$. (Thus $z^t \sim \Pois(\bar\kappa t)$ marginally for each
$t \geq 0$.) We take
$(\theta^0,\theta^*)$, $\eps$, $\{u^t\}$, $(\{w^t\},w^*)$,
and $\{z^t\}$ to be mutually independent, and
constructed on a filtered probability space $(\Omega,\cF,\{\cF_t\}_{t \in
[0,T]},\P)$ where $\{\cF_t\}_{t \in [0,T]}$ is a complete and right-continuous
filtration, $\theta^0,\theta^*,\eps,w^*$ are $\cF_0$-measurable, and
$\{u^t\},\{w^t\},\{z^t\}$ are $\cF_t$-adapted and c\`adl\`ag (right-continuous
with left limits).

Recalling the limit
$\gamma=\lim_{n,d \to \infty} n/d$ of Assumption \ref{asp:model}(a),
we define the primary DMFT processes $\{\theta^t\}_{t \geq 0}$
and $\{\xi^t\}_{t \geq 0}$ on $\R^k$ as
\begin{align}
\theta^t&=\theta^0-\int_0^t \bar\eta^r \Big({\gamma}\Gamma^r
\theta^r+g(\theta^r)+{\gamma}R_f^r(\theta^{[r]})+
{\gamma}R_f^{r,*}\theta^*\Big)\d r+\sqrt{\gamma}\,u^t,
\label{eq:def-cont-theta}\\
\xi^t&={-}\int_0^t \frac{\bar\eta^r}{\bar\kappa} R_\theta^{t,r}
f(\xi^{r-},w^*,\eps)\d z^r+w^t,\label{eq:def-cont-xi}
\end{align}
where we write $x^{[s]}=\{x^t\}_{t \in [0,s]}$ for the restriction of a process
$\{x^t\}$ to times $t \in [0,s]$, and
$x^{t-}=\lim_{s \uparrow t} x^s$ for any c\`adl\`ag process $\{x^t\}$.
We define auxiliary response
processes $\{r_\theta^{t,s}\}_{t,s \geq 0}$ on $\R^{k \times k}$
and $\{r_f^{t,*}\}_{t \geq 0}$ on $\R^{k \times k^*}$ by
\begin{align}
r_\theta^{t,s}&=\Id_k-
\int_s^t \bar\eta^r\bigg[\Big({\gamma}\Gamma^r+\D
g(\theta^r)\Big)r_\theta^{r,s}+{\gamma}
R_f^r(r_\theta^{[r],s})\bigg]\d r \text{ for } t \geq s,
\label{eq:def-cont-deriv-t}\\
&\hspace{1in}r_\theta^{t,s}=0 \text{ for } t<s,\notag\\
r_f^{t,*}&={-}\D_\xi f(\xi^t,w^*,\eps) \int_0^t
\frac{\bar\eta^r}{\bar\kappa}
R_\theta^{t,r}r_f^{r-,*}\d z^r
+\D_{w^*} f(\xi^t,w^*,\eps),\label{eq:def-cont-deriv-xi-star}
\end{align}
and a process of random linear operators $\{r_f^t\}_{t \geq 0}$
with $r_f^t:L^4([0,t],\R^k) \to \R^k$ by
\begin{align}
r_f^t(x^{[t]})&={-}\D_\xi
f(\xi^t,w^*,\eps) \int_0^t \frac{\bar\eta^r}{\bar\kappa}
R_\theta^{t,r}\bigg(r_f^{r-}(x^{[r]})+\D_\xi
f(\xi^{r-},w^*,\eps)x^r\bigg)\d z^r.\label{eq:def-cont-resp-xi}
\end{align}
Here, $\D g(\theta)$ denotes the derivative (i.e.\ Jacobian matrix in $\R^{k
\times k}$) of $g(\cdot)$, and similarly
$\D_\xi f(\xi,w^*,\eps)$ and $\D_{w^*} f(\xi,w^*,\eps)$
denote the derivatives of $f(\cdot)$ in $\xi$ and $w^*$. In
(\ref{eq:def-cont-deriv-t}--\ref{eq:def-cont-resp-xi}),
$r_\theta^{[r],s}=\{r_\theta^{t,s}\}_{t \in [0,r]}$,
$R_f^r(r_\theta^{[r],s})$ is understood via its application to matrix-valued
processes in \eqref{eq:matrixRf}, and $r_f^{t-,*}=\lim_{s \uparrow t} r_f^{s,*}$
and $r_f^{t-}(x^{[t]})=\lim_{s \uparrow t} r_f^s(x^{[s]})$.

The fixed-point relations for
$(C_\theta,R_\theta,C_f,R_f,R_f^*,\Gamma)$ are then given by
\begin{align}
C_\theta^{t,s}&=\E[\theta^t \otimes \theta^s] \text{ for } t,s \in [0,\infty)
\cup \{*\}
\label{eq:def-cont-C-t}\\
R_\theta^{t,s}&=\E[r_\theta^{t,s}]
\label{eq:def-cont-R-t}\\
C_f^{t,s}&=\E\Big[\int_0^t \frac{\bar\eta^r}{\bar\kappa}f(\xi^{r-},w^*,\eps)\d z^r
\otimes \int_0^s \frac{\bar\eta^r}{\bar\kappa}f(\xi^{r-},w^*,\eps)\d z^r\Big]\label{eq:def-cont-C-h}\\
R_f^t(x^{[t]})&=\E\big[r_f^t(x^{[t]})\big]\label{eq:def-cont-R-h}\\
R_f^{t,*}&=\E[r_f^{t,*}]
\label{eq:def-cont-R-h-star}\\
\Gamma^t&=\E[\D_\xi f(\xi^t,w^*,\eps)]\label{eq:def-cont-Gamma}
\end{align}
We clarify that in (\ref{eq:def-cont-R-h}), the input process $x^{[t]}=\{x^s\}_{s \in [0,t]}$ is
understood as deterministic, and the expectation is taken with respect to
$\{\xi^t\}_{t \geq 0}$, $\{z^t\}_{t \geq 0}$, $w^*$, and $\eps$ defining
(\ref{eq:def-cont-resp-xi}). This yields a deterministic linear operator
$R_f^t:L^4([0,t],\R^k)\to \R^k$, which is then applied to the stochastic
inputs in (\ref{eq:def-cont-theta}) and (\ref{eq:def-cont-deriv-t}).

\subsubsection{Limit processes for SME}

For the SME diffusion process $\{\btheta^t\}_{t \geq 0}$
of (\ref{eq:SME}), the description of its high-dimensional limit is identical
to the above
equations (\ref{eq:thetaeps}--\ref{eq:def-cont-Gamma}) for SGD, except
with the Poisson jump process $\{z^t\}_{t \geq 0}$ in (\ref{eq:poisson-z})
replaced by the scalar Gaussian diffusion process
\begin{equation}\label{eq:gaussian-z}
z^t=\int_0^t \bar\kappa\,\d s+\int_0^t \sqrt{\bar\kappa}\,\d b^s,
\end{equation}
where $\{b^t\}_{t \geq 0}$ is a $\cF_t$-adapted standard univariate Brownian
motion. I.e., the Poisson process integrals in (\ref{eq:def-cont-xi}),
(\ref{eq:def-cont-deriv-xi-star}), (\ref{eq:def-cont-resp-xi}), and
(\ref{eq:def-cont-C-h}) with respect to $\d z^r$
are replaced by Wiener integrals with respect to
$\bar\kappa\,\d r+\sqrt{\bar\kappa}\,\d b^r$. Note
that the two processes $\{z^t\}_{t \geq 0}$ defined by \eqref{eq:poisson-z}
and \eqref{eq:gaussian-z} have the same mean and covariance, but differ
otherwise in law.

\begin{remark}[Forms of the response processes]\label{rmk:response}
In the above DMFT systems for both SGD and SME,
one may understand $r_f^{t,*}$ as a linear response
\begin{equation}\label{eq:rstarinterp}
r_f^{t,*}=\frac{\partial f(\xi^t,w^*,\eps)}{\partial w^*},
\end{equation}
where (\ref{eq:def-cont-deriv-xi-star}) is computed by 
substituting \eqref{eq:def-cont-xi} in \eqref{eq:rstarinterp} and
differentiating in $w^*$.

If $C_f^{t,s}$ were twice continuously-differentiable in $(t,s)$, then
setting $\bar C_f^{t,s}=\gamma \partial_t \partial_s C_f^{t,s}$, we may write
(\ref{eq:def-cont-theta}) as
\begin{equation}\label{eq:def-cont-theta-smooth}
\theta^t=\theta^0+\int_0^t \Big[{-}\bar\eta^r\Big({\gamma}
\Gamma^r\theta^r+g(\theta^r)+{\gamma}R_f^r(\theta^{[r]})
+{\gamma}R_f^{r,*}\theta^*\Big)+\bar u^r\Big]\d r
\end{equation}
where $\sqrt{\gamma}\,u^t=\int_0^t \bar u^r \d r$ and $\{\bar u^t\}_{t \geq 0}
\sim \GP(0,\bar C_f)$. Then $r_\theta^{t,s}$ may likewise be understood as the
linear response
\begin{equation}\label{eq:Rthetainterp}
r_\theta^{t,s}=\frac{\partial \theta^t}{\partial \bar u^s}
\end{equation}
computed by formally differentiating (\ref{eq:def-cont-theta-smooth})
in $\bar u^s$, and the equations \eqref{eq:def-cont-theta-smooth} and
\eqref{eq:Rthetainterp} would be analogous to those describing gradient flow in
\cite{celentano2021high}. We note, however, that in our
settings of interest where $\{z^t\}_{t \geq 0}$ is a Poisson process
\eqref{eq:poisson-z} or Gaussian diffusion process \eqref{eq:gaussian-z},
$C_f^{t,s}$ as defined by (\ref{eq:def-cont-C-h})
is not twice differentiable at the diagonal $s=t$, and $\{u^t\}_{t \geq 0}
\sim \GP(0,C_f)$ does not admit a representation $\sqrt{\gamma}\,u^t=\int_0^t
\bar u^r \d r$ in the sense of a usual Lebesgue-Stieltjes integral, as it
does not have differentiable sample paths. Thus these
representations \eqref{eq:def-cont-theta-smooth} and \eqref{eq:Rthetainterp} 
should be understood only in a formal sense.

One may also understand the linear operator $r_f^t(x^{[t]})$
in \eqref{eq:def-cont-resp-xi} as an integrated linear response
\begin{equation}\label{eq:responseinterp}
r_f^t(x^{[t]})
=\int_0^t \frac{\partial f(\xi^t,w^*,\eps)}{\partial w^q}\,x^q \d q
\end{equation}
against the test process $x^{[t]}=\{x^s\}_{s \in [0,t]}$.
The evolution equation (\ref{eq:def-cont-resp-xi}) is computed by formally
differentiating $f(\xi^t,w^*,\eps)$ in $w^q$ using (\ref{eq:def-cont-xi}),
\[\frac{\partial f(\xi^t,w^*,\eps)}{\partial w^q}
={-}\D_\xi f(\xi^t,w^*,\eps) \int_q^t \frac{\bar\eta^r}{\bar\kappa}
R_\theta^{t,r}\bigg(\frac{\partial f(\xi^{r-},w^*,\eps)}{\partial
w^q}+\D_\xi f(\xi^{q-},w^*,\eps)\1_{r=q}\bigg)\d z^r.\]
To give a simple and mathematically rigorous meaning to the quantity
$\int_q^t \1_{r=q} \d z^r$, we further integrate against $x^q\,\d q$ and
exchange the orders of integration in $r$ and $q$, yielding
\begin{align*}
&\int_0^t \frac{\partial f(\xi^t,w^*,\eps)}{\partial w^q}\,x^q\d q\\
&={-}\D_\xi f(\xi^t,w^*,\eps)\int_0^t\frac{\bar\eta^r}{\bar\kappa}
R_\theta^{t,r}\Bigg[\int_0^r \bigg(\frac{\partial f(\xi^{r-},w^*,\eps)}{\partial
w^q}+\D_\xi f(\xi^{q-},w^*,\eps)\1_{r=q}\bigg)x^q\d q\Bigg]\d z^r\\
&={-}\D_\xi f(\xi^t,w^*,\eps)\int_0^t\frac{\bar\eta^r}{\bar\kappa}
R_\theta^{t,r}\Bigg[\int_0^r \frac{\partial f(\xi^r,w^*,\eps)}{\partial
w^q}x^q\d q+\D_\xi f(\xi^{r-},w^*,\eps)x^r\Bigg]\d z^r.
\end{align*}
This gives the equation (\ref{eq:def-cont-resp-xi}) under the identification
(\ref{eq:responseinterp}).

We provide these remarks only to clarify a more intuitive interpretation of
these response processes; our mathematical proofs will operate directly on the
definitions (\ref{eq:def-cont-theta}--\ref{eq:def-cont-Gamma}),
rather than these formal interpretations.
\end{remark}

\subsubsection{Existence and uniqueness of the DMFT fixed point}

The following theorem establishes existence and
uniqueness of the above fixed points $(C_\theta,R_\theta,C_f,R_f,R_f^*,\Gamma)$ 
in a suitable domain $\cS^\text{cont} \equiv \cS^\text{cont}(T,C_0)$.
We defer a precise definition of this
domain to Section \ref{sec:existenceuniqueness}.

\begin{thm}\label{thm:uniqueness-existence}
Fix any large enough constant $C_0 \equiv C_0(T)>0$, and let
$\cS \equiv \cS(T,C_0)$ and
$\cS^\text{cont} \equiv \cS^\text{cont}(T,C_0)$ be the spaces for
$(C_\theta,R_\theta,C_f,R_f,R_f^*,\Gamma)$ given in
Definition \ref{def:S}. In both settings of
$\{z^t\}_{t \geq 0}$ in (\ref{eq:poisson-z}) for SGD and 
in (\ref{eq:gaussian-z}) for SME, the following hold:
\begin{enumerate}[(a)]
\item Given any $(C_\theta,R_\theta,C_f,R_f,R_f^*,\Gamma) \in \cS$, and given
the filtered probability space $(\Omega,\cF,\{\cF_t\}_{t \geq 0},\P)$
containing $(\theta^0,\theta^*),\eps,\{u^t\}_{t \geq 0},(\{w^t\}_{t \geq
0},w^*),\{z^t\}_{t \geq 0}$ defined by \eqref{eq:thetaeps}, \eqref{eq:GPs},
and \eqref{eq:poisson-z}/\eqref{eq:gaussian-z},
there exist unique $\cF_t$-adapted and c\`adl\`ag solutions to
\eqref{eq:def-cont-theta} and \eqref{eq:def-cont-deriv-t} for each fixed $s \in [0,T]$,
and unique $\cF_t$-adapted and c\`adl\`ag solutions to
\eqref{eq:def-cont-xi}, \eqref{eq:def-cont-deriv-xi-star}, and
\eqref{eq:def-cont-resp-xi} for each fixed deterministic input process
$x \in L^4([0,T],\R^k)$.
\item There exists a unique point $(C_\theta,R_\theta,C_f,R_f,R_f^*,\Gamma) \in
\cS^\text{cont} \subset \cS$ that satisfies the fixed-point conditions
(\ref{eq:def-cont-C-t}--\ref{eq:def-cont-Gamma}) defined by the solutions of
part (a).
\end{enumerate}
\end{thm}

This theorem guarantees, in particular, that the processes
$\{\theta^t\}_{t \in [0,T]}$ and $\{\xi^t\}_{t \in [0,T]}$ which
will characterize the high-dimensional limits of SGD and SME are well-defined.

\subsection{Convergence to the limit processes}

The following theorem is the main result of this paper. Recalling the data
matrix
\[\X=[\x_1,\ldots,\x_n]^\top \in \R^{n \times d},\]
we show that over fixed time horizons $[0,T]$, finite-time marginals
of the empirical
distributions of coordinates of $(\{\btheta^t\}_{t \in [0,T]},\btheta^*)$
and $(\{\X\btheta^t\}_{t \in [0,T]},\X\btheta^*,\beps)$ for both the rescaled SGD
process \eqref{eq:SGDrescaled} and the SME process \eqref{eq:SME}
converge in law to the preceding DMFT processes $\{\theta^t\}_{t \geq 0}$ and
$\{\xi^t\}_{t \geq 0}$,
defined via the fixed point $(C_\theta,R_\theta,C_f,R_f,R_f^*,\Gamma) \in
\cS^\text{cont} \equiv \cS^\text{cont}(T,C_0)$.

\begin{thm}\label{thm:DMFT-limit}
Fix a time horizon $T>0$ and $0 \leq t_1 \leq \ldots \leq t_m \leq T$
not depending on $n,d$. For any sufficiently large constant $C_0>0$, let
$\cS^\text{cont} \equiv \cS^\text{cont}(T,C_0)$ be the domain of Theorem
\ref{thm:uniqueness-existence}.
\begin{enumerate}[(a)]
\item Let $\{\btheta^t\}_{t \in [0,T]}$
be the time-rescaled SGD process (\ref{eq:SGDrescaled})
where $\btheta^t=[\theta_1^t,\ldots,\theta_d^t]^\top \in \R^{d \times k}$,
and let $\Law(\cdot)$ denote the joint law of
$(\{\theta^t\}_{t \in [0,T]},\theta^*)$ or
$(\{\xi^t\}_{t \in [0,T]},w^*,\eps)$ in the DMFT
system (\ref{eq:thetaeps}--\ref{eq:def-cont-Gamma}) when $\{z^t\}_{t \geq 0}$ is the Poisson process
in (\ref{eq:poisson-z}),
defined by the unique fixed point $(C_\theta,R_\theta,C_f,R_f,R_f^*,\Gamma) \in
\cS^\text{cont}$. Then almost surely as $n,d \to \infty$,
weakly and in Wasserstein-2 over $(\R^k)^m \times \R^{k^*}$ and $(\R^k)^m \times \R^{k^*} \times \R$,
\begin{align}
\frac{1}{d}\sum_{j=1}^d
\delta_{(\theta_j^{t_1},\ldots,\theta_j^{t_m},\theta_j^*)}
&\Rightarrow
\Law(\theta^{t_1},\ldots,\theta^{t_m},\theta^*),\label{eq:DMFT-limit-theta}\\
\frac{1}{n}\sum_{i=1}^n \delta_{(\x_i^\top \btheta^{t_1},
\ldots,\x_i^\top \btheta^{t_m},\x_i^\top \btheta^*,\eps_i)}
&\Rightarrow \Law(\xi^{t_1},\ldots,\xi^{t_m},w^*,\eps).\label{eq:DMFT-limit-xi}
\end{align}
\item Let $\{\btheta^t\}_{t \in [0,T]}$ be the SME process
(\ref{eq:SME}). Then the same result of part (a) holds, upon replacing
the DMFT system (\ref{eq:thetaeps}--\ref{eq:def-cont-Gamma}) by that
in the Gaussian setting of $\{z^t\}_{t \geq 0}$ given in (\ref{eq:gaussian-z}).
\end{enumerate}
\end{thm}

\begin{remark}
Under Assumption \ref{asp:SGD}, the characterization of the limiting law for SGD
in Section \ref{sec:DMFT-SGD} depends only on
$(\bar\kappa,\{\bar\eta^t\}_{t \geq 0})$ and not on $\alpha$. Thus the
high-dimensional limit of the SGD dynamics in Theorem \ref{thm:DMFT-limit}(a)
is the same for all $\alpha \in [0,1)$.
\end{remark}

\begin{remark}
Theorem \ref{thm:DMFT-limit} implies, as a consequence, that for any
pseudo-Lipschitz test functions
$\psi:(\R^k)^m \times \R^{k^*} \to \R$ 
and $\phi:(\R^k)^m \times \R^{k^*} \times \R \to \R$
(i.e.\ satisfying
$|\psi(x)-\psi(y)| \leq C\|x-y\|_2(1+\|x\|_2+\|y\|_2)$ for all inputs $x,y \in (\R^k)^m \times \R^{k^*}$), almost surely
\begin{align*}
\lim_{n,d \to \infty} \frac{1}{d}\sum_{j=1}^d
\psi(\theta_j^{t_1},\ldots,\theta_j^{t_m},\theta_j^*)
&=\E\,\psi(\theta^{t_1},\ldots,\theta^{t_m},\theta^*),\\
\lim_{n,d \to \infty} \frac{1}{n}\sum_{i=1}^n
\phi(\x_i^\top \btheta^{t_1},\ldots,\x_i^\top \btheta^{t_m},\x_i^\top
\btheta^*,\eps_i)
&=\E\,\phi(\xi^{t_1},\ldots,\xi^{t_m},w^*,\eps),
\end{align*}
where the expectations on the right side are evaluated over the joint laws of
the DMFT variables.

For example, for any fixed training time $t>0$ (in units of training epochs for
SGD, c.f.\ the time rescaling \eqref{eq:SGDrescaled}), almost surely as $n,d \to \infty$:
\begin{enumerate}
\item The pairwise overlaps between columns of the learned parameter
$\btheta^t \in \R^{d \times k}$ and true parameter
$\btheta^* \in \R^{d \times k^*}$ satisfy
\[\frac{1}{d}\,\btheta^{t\top}\btheta^*
\to \E[\theta^t \otimes \theta^*] \in \R^{k \times k^*}.\]
\item The average training loss satisfies
\[\frac{1}{n}\sum_{i=1}^n L\big(\sigma(\x_i^\top \btheta^t,y_i)\big)
\to \E L\big(\sigma(\xi^t),\sigma^*(w^*,\eps))\big).\]
\item Suppose further that the data $\x_i \sim \cN(0,\frac{1}{d}\,\Id)$ is Gaussian. Then for an
independent test
sample $\x_\mathrm{test} \sim \cN(0,\frac{1}{d}\,\Id)$ equal in law to $\x_1,\ldots,\x_n$, conditional on $(\btheta^t,\btheta^*)$ we have
\[(\x_\text{test}^\top\btheta^t,\x_\text{test}^\top\btheta^*)
\sim \cN\left(0,\frac{1}{d}(\btheta^t,\btheta^*)^\top(\btheta^t,\btheta^*)\right).\]
Thus for any test loss $L_\text{test}:\R^k \times \R^{k^*} \to \R$, the expected
test error admits a representation
\[\E_{\x_\mathrm{test}} L_\text{test}(\x_{\text{test}}^\top \btheta^t,\x_{\text{test}}^\top \btheta^*)
=h\left(\frac{1}{d}(\btheta^t,\btheta^*)^\top(\btheta^t,\btheta^*)\right)\]
for some smooth function $h:\R^{(k+k^*) \times (k+k^*)} \to \R$. Then also
\[\E_{\x_\mathrm{test}} L_\text{test}(\x_{\text{test}}^\top \btheta^t,\x_{\text{test}}^\top \btheta^*)
\to h\big(\E[(\theta,\theta^*) \otimes (\theta,\theta^*)]\big).\]
\end{enumerate}
\end{remark}


\subsection{Discussion}

In this section, to provide further interpretation of our results and some
connections to related literature, we give a (heuristic) discussion of some
simplifications of the preceding DMFT
equations in the limit of small learning rate, limit of large sample size,
and specialization to a setting of quadratic optimization for squared loss and
ridge regularizer.

\subsubsection{Small learning rate limit and gradient flow}\label{sec:comp-gf}

Consider, for simplicity, a time-independent learning rate
$\bar\eta^r=\bar\eta$.
We check that in the small learning rate limit $\bar\eta \to 0$,
the DMFT system (\ref{eq:def-cont-theta}--\ref{eq:def-cont-Gamma}) converges to
that obtained for gradient flow in \cite{celentano2021high}.

Introduce the rescaled time $\tau=\bar\eta t$,
and define the time-rescaled processes
\begin{equation}\label{eq:rescaled1}
\begin{gathered}
\tilde\theta^\tau=\theta^{\tau/\bar\eta},\quad
\tilde \xi^\tau=\xi^{\tau/\bar\eta},\quad
\tilde u^\tau = u^{\tau/\bar\eta},\quad
\tilde w^\tau=w^{\tau/\bar\eta},\quad
\tilde z^\tau=z^{\tau/\bar\eta},\\
\tilde r_\theta^{\tau,\sigma}=r_\theta^{\tau/\bar\eta,\sigma/\bar\eta},\quad
\tilde r_f^{\tau,*}=r_f^{\tau/\bar\eta,*},\quad
\tilde r_f^\tau(\tilde x^{[\tau]})=r_f^{\tau/\bar\eta}(x^{[\tau/\bar\eta]})
\text{ where } \tilde x^\tau=x^{\tau/\bar\eta},\end{gathered}
\end{equation}
together with the corresponding deterministic response and covariance processes
\begin{equation}\label{eq:rescaled2}
\begin{gathered}
\tilde C_\theta^{\tau,\sigma}=C_\theta^{\tau/\bar\eta,\sigma/\bar\eta}, \quad
\tilde R_\theta^{\tau,\sigma}=R_\theta^{\tau/\bar\eta,\sigma/\bar\eta}, \quad
\tilde C_f^{\tau,\sigma}=C_f^{\tau/\bar\eta,\sigma/\bar\eta},\\
\tilde R_f^{\tau}(x^{[\tau]})=R_f^{\tau/\bar\eta}(x^{[\tau/\bar\eta]}), \quad
\tilde R_f^{\tau,*}=R_f^{\tau/\bar\eta,*}, \quad
\tilde \Gamma^\tau=\Gamma^{\tau/\bar\eta}.\end{gathered}
\end{equation}
Note that $\{\tilde u^\tau\} \sim \GP(0,\tilde C_f)$,
$(\{\tilde w^\tau\},w^*) \sim \GP(0,\tilde C_\theta)$, and $\{\tilde z^\tau\}$
is a Poisson process with rate $\bar\kappa/\bar\eta$ in the SGD setting or a
diffusion process $\tilde z^\tau=\int_0^\tau (\bar\kappa/\bar\eta)\d \rho
+\int_0^\tau \sqrt{\bar\kappa/\bar\eta}\,\d b^\rho$ in the SME setting.

Applying this change of variables to
(\ref{eq:def-cont-theta}--\ref{eq:def-cont-resp-xi}), we have
\begin{align}
\tilde\theta^{\tau}&= \theta^0 - \int_0^{\tau} \Big({\gamma}\tilde\Gamma^\rho
\tilde\theta^\rho+g(\tilde\theta^\rho)+{\gamma}\tilde R_f^\rho(\tilde\theta^{[\rho]})+
{\gamma}\tilde R_f^{\rho,*}\theta^*\Big)\d \rho+\sqrt{\gamma}\,\tilde
u^{\tau},\label{eq:thetaGF}\\
\tilde r_\theta^{\tau,\sigma}&=\Id_k-\int_\sigma^\tau
\Big[\Big({\gamma}\tilde\Gamma^\rho+\D g(\tilde
\theta^\rho)\Big)\tilde r_\theta^{\rho,\sigma}
+\gamma\tilde R_f^\rho(\tilde
r_\theta^{[\rho],\sigma})\Big]\d \rho \text{ if } \tau \geq \sigma,
\quad \tilde r_\theta^{\tau,\sigma}=0 \text{ if }
\tau<\sigma,\label{eq:rthetaGF}\\
\tilde\xi^\tau&={-}\frac{\bar\eta}{\bar\kappa}
\int_0^\tau \tilde R_\theta^{\tau,\rho}
f(\tilde\xi^{\rho-},w^*,\eps)\d \tilde z^\rho+\tilde w^\tau,\\
\tilde r_f^{\tau,*}&={-}\D_\xi f(\tilde \xi^\tau,w^*,\eps)
\frac{\bar\eta}{\bar\kappa}
\int_0^\tau \tilde R_\theta^{\tau,\rho}\tilde r_f^{\rho-,*}\d \tilde z^\rho
+\D_{w^*} f(\tilde \xi^\tau,w^*,\eps),\\
\tilde r_f^\tau(\tilde x^{[\tau]})
&={-}\D_\xi f(\tilde \xi^\tau,w^*,\eps)
\frac{\bar\eta}{\bar\kappa}\int_0^\tau \tilde R_\theta^{\tau,\rho}
\Big(\tilde r_f^{\rho-}(\tilde x^{[\rho]})
+\D_\xi f(\tilde \xi^{\rho-},w^*,\eps)\tilde x^\rho\Big)\d \tilde z^\rho.
\end{align}
From this, it may be verified that $\tilde R_\theta,\tilde R_f,\tilde R_f^*$
remain uniformly bounded over finite intervals of rescaled time
$\tau,\sigma \in [0,\tilde T]$, as $\bar\eta \to 0$. Then, decomposing
the evolution of $\tilde \xi^\tau$ as
\[\tilde\xi^\tau= {-}\int_0^\tau \tilde R_\theta^{\tau,\rho}
f(\tilde\xi^\rho,w^*,\eps)\d \rho - \underbrace{\frac{\bar\eta}{\bar\kappa}
\int_0^\tau \tilde R_\theta^{\tau,\rho}
f(\tilde\xi^{\rho-},w^*,\eps)(\d \tilde
z^\rho-(\bar\kappa/\bar\eta)\d\rho)}_{:=M^\tau}+\,\tilde w^\tau,\]
we note that the martingale term $M^\tau$ satisfies
$\E\|M^\tau\|^2 \lesssim (\bar\eta/\bar\kappa)^2
\cdot (\bar\kappa/\bar\eta)=O(\bar\eta)$
over finite time horizons $\tau \in [0,\tilde T]$,
by the It\^o isometry. Then a
high-moment bound may be applied to show
$\lim_{\bar\eta \to 0} \sup_{\tau \in [0,\tilde T]} \|M^\tau\|=0$. The
same argument establishes that the martingale terms for the evolutions of
$\tilde r_f^{\tau,*}$ and $\tilde r_f^\tau(\tilde x^{[\tau]})$ vanish in the
limit $\bar\eta \to 0$, leading to the simplified equations
\begin{align}
\tilde\xi^\tau&={-}
\int_0^\tau \tilde R_\theta^{\tau,\rho}
f(\tilde\xi^{\rho},w^*,\eps)\d \rho+\tilde w^\tau,\label{eq:xiGF}\\
\tilde r_f^{\tau,*}&={-}\D_\xi f(\tilde \xi^\tau,w^*,\eps)
\int_0^\tau \tilde R_\theta^{\tau,\rho}\tilde r_f^{\rho,*}\d \rho
+\D_{w^*} f(\tilde \xi^\tau,w^*,\eps),\label{eq:rfstarGF}\\
\tilde r_f^\tau(\tilde x^{[\tau]})
&={-}\D_\xi f(\tilde \xi^\tau,w^*,\eps)
\int_0^\tau \tilde R_\theta^{\tau,\rho}
\Big(\tilde r_f^\rho(\tilde x^{[\rho]})
+\D_\xi f(\tilde \xi^\rho,w^*,\eps)\tilde x^\rho\Big)\d\rho.
\end{align}
Positing the form
\[\tilde r_f^\tau (\tilde x^{[\tau]})
=\int_0^\tau \frac{\partial f(\tilde\xi^\tau,w^*,\eps)}{\partial \tilde w^\sigma}
\,\tilde x^\sigma \d \sigma\]
as in Remark \ref{rmk:response}, this response process 
$\frac{\partial f(\tilde\xi^\tau,w^*,\eps)}{\partial \tilde w^\sigma}$ must then
satisfy
\begin{equation}
\frac{\partial f(\tilde\xi^\tau,w^*,\eps)}{\partial \tilde w^\sigma}
={-}\D_\xi f(\tilde\xi^\tau,w^*,\eps) \left(\int_\sigma^\tau
\tilde R_\theta^{\tau,\rho}\frac{\partial f(\tilde\xi^\rho,w^*,\eps)}{\partial
\tilde w^\sigma} \d \rho +
\tilde R_\theta^{\tau,\sigma}\D_\xi f(\tilde \xi^{\sigma},w^*,\eps)\right).
\label{eq:rfGF}
\end{equation}
These equations \eqref{eq:thetaGF}, \eqref{eq:rthetaGF}, \eqref{eq:xiGF},
\eqref{eq:rfstarGF}, and \eqref{eq:rfGF} (specialized to the setting of a
ridge-type regularizer $G(\theta)=\frac{1}{2}\theta^\top \Lambda \theta$)
are equivalent to the DMFT system established in \cite{celentano2021high}
for gradient flow.

\subsubsection{Large-sample limit and one-pass SGD}\label{sec:comp-onepass}

Recalling the parameter $\gamma=\lim_{n,d \to \infty} n/d$, we next consider a
data-rich regime where $\gamma \to \infty$. For simplicity, 
consider a constant learning rate $\eta^t=\bar\eta^t=\bar \eta$
and single-sample batch size $\kappa=\bar\kappa=1$.
To keep the stochastic gradient dynamics \eqref{eq:SGD} the same as $n$ grows
(i.e.\ depending on $d$ rather than $n$), we define
\[\tilde g(\theta)=\gamma^{-1} g(\theta)\]
and fix $\tilde g(\cdot)$ as $\gamma \to \infty$.
Rescaling time as $\tau=\gamma t$, let us show that over any fixed time horizon
$\tau \in [0,\tilde T]$ (corresponding to a number of SGD iterations $n(\tilde
T/\gamma) \approx d\tilde T$ that 
does not depend on $n$), the preceding DMFT equations
reduce to a simple diffusion process that characterizes the one-pass or online
SGD procedure
\[\bar\btheta^{k+1}=\bar\btheta^k-\bar \eta
\Big(\x^k \otimes f(\x^{k\top} \bar\btheta^k,\x^{k\top}\btheta^*,\eps^k)
+\frac{1}{d}\,\tilde g(\bar\btheta^k)\Big)\]
with a fresh sample $(\x^k,y^k)=(\x^k,\sigma^*(\x^{k\top}\btheta^*,\eps^k))$
per iteration. We note that this is the standard one-pass SGD procedure for 
minimizing the population risk
\[\E[L\big(\sigma(\x^\top \btheta),y\big)]
+\frac{1}{d}\sum_{j=1}^d \tilde G(\theta_j), \qquad \tilde G(\theta)
=\gamma^{-1} G(\theta).\]

We define the time-rescaled processes $\tilde\theta^\tau=\theta^{\tau/\gamma}$,
$\tilde \xi^\tau=\xi^{\tau/\gamma}$, $\tilde z^\tau=z^{\tau/\gamma}$,
etc.\ analogously to
(\ref{eq:rescaled1}--\ref{eq:rescaled2}). Then applying this change of variables
to (\ref{eq:def-cont-theta}--\ref{eq:def-cont-resp-xi}),
\begin{align*}
\tilde\theta^{\tau}&=\theta^0 - \bar\eta \int_0^{\tau} \Big(\tilde\Gamma^\rho
\tilde\theta^\rho+\tilde g(\tilde\theta^\rho)+\tilde R_f^\rho(\tilde\theta^{[\rho]})+\tilde R_f^{\rho,*}\theta^*\Big)\d \rho+\sqrt{\gamma}\,\tilde
u^{\tau},\\
\tilde r_\theta^{\tau,\sigma}&=\Id_k-\bar\eta\int_\sigma^\tau
\Big[\Big(\tilde\Gamma^\rho+\D g(\tilde \theta^\rho)\Big)
\tilde r_\theta^{\rho,\sigma}+\tilde R_f^\rho(\tilde
r_\theta^{[\rho],\sigma})\Big]\d \rho \text{ if } \tau \geq \sigma,
\quad \tilde r_\theta^{\tau,\sigma}=0 \text{ if }
\tau<\sigma,\\
\tilde\xi^\tau&={-}\bar\eta
\int_0^\tau \tilde R_\theta^{\tau,\rho}
f(\tilde\xi^{\rho-},w^*,\eps)\d \tilde z^\rho+\tilde w^\tau,\\
\tilde r_f^{\tau,*}&={-}\D_\xi f(\tilde \xi^\tau,w^*,\eps)
\cdot \bar\eta
\int_0^\tau \tilde R_\theta^{\tau,\rho}\tilde r_f^{\rho-,*}\d \tilde z^\rho
+\D_{w^*} f(\tilde \xi^\tau,w^*,\eps),\\
\tilde r_f^\tau(\tilde x^{[\tau]})
&={-}\D_\xi f(\tilde \xi^\tau,w^*,\eps)
\cdot \bar\eta \int_0^\tau \tilde R_\theta^{\tau,\rho}
\Big(\tilde r_f^{\rho-}(\tilde x^{[\rho]})
+\D_\xi f(\tilde \xi^{\rho-},w^*,\eps)\tilde x^\rho\Big)\d \tilde z^\rho.
\end{align*}
Here, $\{\tilde z^\tau\}_{\tau \geq 0}$ is a Poisson process with rate
$1/\gamma$ in the setting of SGD, or the diffusion process $\tilde
z^\tau=\int_0^\tau (1/\gamma)\d \rho+\int_0^\tau (1/\sqrt{\gamma})\d b^\rho$
in the setting of SME. Then one may again verify that $\tilde
R_\theta,\tilde R_f,\tilde R_f^*$ remain uniformly bounded as $\gamma \to
\infty$, and that in this limit, all stochastic integrals against $\d \tilde
z^\rho$ vanish to give the simplified equations
\begin{equation}\label{eq:theta-timeres-1}
\tilde\theta^\tau=\theta^0-\bar\eta \int_0^\tau \Big(\tilde \Gamma^\rho
\tilde\theta^\rho+\tilde g(\tilde \theta^\rho)+ \tilde
R_f^{\rho,*}\theta^*\Big)\d \rho+\sqrt{\gamma}\,\tilde u^\tau,
\end{equation}
\[\tilde \xi^\tau=\tilde w^\tau,
\quad \tilde r_f^{\tau,*}=\D_{w^*} f(\tilde \xi^\tau,w^*,\eps)
=\D_{w^*} f(\tilde w^\tau,w^*,\eps),
\quad \tilde r_f^\tau \equiv 0.\]
The covariance kernel of the rescaled Gaussian process $\{\sqrt{\gamma}\,\tilde
u^\tau\}_{\tau \geq 0}$ is given by
\begin{align*}
\gamma\,\tilde C_f^{\tau,\sigma} &= \gamma\,\E\Big[\bar\eta \int_0^\tau
f(\tilde\xi^{\rho-},w^*,\eps)\d \tilde z^\rho
\otimes \bar\eta\int_0^\sigma f(\tilde\xi^{\rho-},w^*,\eps)\d \tilde
z^\rho\Big].
\end{align*}
Applying the It\^o isometry and above simplification
$\tilde \xi^\tau=\tilde w^\tau$ in the $\gamma \to \infty$ limit,
we then obtain
\begin{align*}
\lim_{\gamma \to \infty} \gamma\,\tilde C_f^{\tau,\sigma}
&=\lim_{\gamma \to \infty}\left(O\Big(\frac{1}{\sqrt{\gamma}}\Big) +
\gamma\bar\eta^2\,\E\Big[\Big(\int_0^{\tau\wedge\sigma} f(\tilde
\xi^{\rho-},w^*,\eps)(\d \tilde z^r - (1/\gamma)\d r)\Big)^{\otimes
2}\Big]\right)\\
&=\bar \eta^2\int_0^{\tau\wedge\sigma} \E\Big[
f(\tilde w^\rho,w^*,\eps)^{\otimes 2}\Big]\d \rho.
\end{align*}
This is precisely the covariance kernel of a Brownian diffusion
\[\d \tilde u^\tau=\bar\eta \Sigma^\tau\d B^\tau,
\qquad \Sigma^\tau=\E[f(\tilde w^\tau,w^*,\eps)^{\otimes 2}]^{1/2},\]
where $\{B^\tau\}_{\tau \geq 0}$ is a standard Brownian motion on $\R^k$.

Thus, to summarize, in the limit $\gamma \to \infty$,
the law of $\{\tilde \theta^\tau\}_{\tau \geq 0}$ in \eqref{eq:theta-timeres-1}
is given by a distribution-dependent SDE
\begin{equation}\label{eq:onlineSDE}
\d \tilde\theta^\tau={-}\bar\eta \Big(\tilde \Gamma^\tau
\tilde\theta^\tau+\tilde g(\tilde \theta^\tau)+\tilde R_f^{\tau,*}
\theta^*\Big)\d \tau+\bar \eta\,\Sigma^\tau \d B^\tau,
\end{equation}
where the deterministic drift and diffusion
coefficients $\tilde \Gamma^\tau,\tilde R_f^{\tau,*},\Sigma^\tau$ are
defined self-consistently from the law of $(\tilde\theta^\tau,\theta^*)$ by
\begin{equation}\label{eq:SDEcoefs}
\begin{gathered}
\tilde C_\theta=\E[(\tilde \theta^\tau,\theta^*) \otimes (\tilde
\theta^\tau,\theta^*)] \in \R^{(k+k^*) \times (k+k^*)},
\quad (\tilde w^\tau,w^*) \sim \cN(0,\tilde C_\theta),\\
\tilde \Gamma^\tau=\E[\D_\xi f(\tilde w^\tau,w^*,\eps)],
\quad \tilde R_f^{\tau,*}=\E[\D_{w^*} f(\tilde w^\tau,w^*,\eps)],
\quad \Sigma^\tau=\E[f(\tilde w^\tau,w^*,\eps)^{\otimes 2}]^{1/2}.
\end{gathered}
\end{equation}
This coincides with the high-dimensional limit of the
learning dynamics of one-pass SGD. For example, specializing to the case of a
ridge penalty $\tilde G(\theta)=\frac{\lambda}{2}\|\theta\|^2$ and
$\tilde g(\theta)=\lambda \theta$, we may apply It\^o's formula to
\eqref{eq:onlineSDE} to obtain a closed system of ODEs for the overlap
parameters $\tilde C_\theta=\E[(\tilde \theta^\tau,\theta^*) \otimes (\tilde
\theta^\tau,\theta^*)]$, which take the form
\begin{align*}
\frac{\d}{\d \tau} \E [\tilde\theta^\tau \otimes \theta^*] &= -\bar\eta
\left((\tilde \Gamma^\tau+\lambda)\E[\tilde\theta^\tau \otimes \theta^*] +
\tilde R_f^{\tau,*}\E[\theta^* \otimes \theta^*] \right),\\
\frac{\d}{\d \tau} \mathbb{E} [\tilde\theta^\tau \otimes \tilde\theta^\tau] &= -\bar\eta \bigg((\tilde \Gamma^\tau+\lambda)
\E[\tilde\theta^\tau \otimes \tilde\theta^\tau]  + 
\tilde R_f^{\tau,*}\E[\theta^* \otimes \tilde\theta^\tau]\\
&\hspace{1in}+\E[\tilde\theta^\tau \otimes \tilde\theta^\tau]
(\tilde \Gamma^\tau+\lambda)^\top + \E[\tilde\theta^\tau \otimes \theta^*]
\tilde R_f^{\tau,*\top} \bigg) + \bar\eta^2(\Sigma^\tau)^2
\end{align*}
and $\tilde \Gamma^\tau,\tilde R_f^{\tau,*},\Sigma^\tau$ are the functions of
$\tilde C_\theta$ defined by \eqref{eq:SDEcoefs}. This recovers the known ODE
for overlap parameters in one-pass SGD (in a setting of isotropic data), see
e.g.\ \cite[Corollary 1.1]{collins2024hitting}.

\subsubsection{Deterministic integro-differential equations for squared loss and
ridge regularizer}\label{sec:linear}

We next consider a special setting of quadratic optimization, given by a
squared loss, linear activation, and ridge regularizer.\footnote{The result we establish does not
technically encompass this setting, as Assumption \ref{asp:lipschitz}
requires both $L(\cdot,\cdot)$ and $\sigma(\cdot)$ to be Lipschitz. However, we
include this discussion here, as we believe it is illuminating to understand
the simplifications that arise in this linear example.} In this setting,
the DMFT equations become linear, and the relevant mean and covariance
statistics of
$\{\theta^t,\xi^t,r_\theta^{t,s},r_f^t,r_f^{t,*}\}$ may be calculated
analytically to yield closed integro-differential equations for the two mappings
\[(C_\theta,R_\theta) \mapsto (C_f, R_f, R_f^*,\Gamma),
\qquad (C_f, R_f, R_f^*,\Gamma) \mapsto (C_\theta,R_\theta).\]
A corollary of this calculation is that, in this linear setting, the fixed point
$(C_\theta,R_\theta,C_f, R_f, R_f^*,\Gamma)$ depends only on the mean and
covariance of $\{z^t\}_{t \geq 0}$, and hence coincides for SGD and SME. (More generally, this fixed point coincides for the squared loss, linear activation, and any regularizer, as the mapping $(C_\theta,R_\theta) \mapsto (C_f, R_f, R_f^*,\Gamma)$ is always the same for both SGD and SME.)

Specifically, let $k=k^*$ and consider the setting
\[L(\hat y,y) = \frac{1}{2} \norm{\hat y - y}^2, \qquad \sigma(\xi) = \xi,
\qquad f(\xi, w^*, \eps) = \xi - \sigma^*(w^*, \eps).\]
Let us denote
\[\bar \xi^t=\xi^t-\sigma^*(w^*,\eps),
\qquad \bar w^t=w^t-\sigma^*(w^*,\eps),
\qquad \bar z^t=z^t-\bar\kappa\,t.\]
Then \eqref{eq:def-cont-xi}, \eqref{eq:def-cont-deriv-xi-star}, and
\eqref{eq:def-cont-resp-xi} simplify to
\begin{align*}
\bar \xi^t&={-}\int_0^t \frac{\bar\eta^r}{\bar\kappa} R_\theta^{t,r} \bar
\xi^{r-}\d z^r+\bar w^t
={-}\int_0^t {\bar\eta^r} R_\theta^{t,r} \bar \xi^r \d r
\underbrace{-\int_0^t \frac{\bar\eta^r}{\bar\kappa} R_\theta^{t,r}
\bar \xi^{r-}\d \bar z^r}_{:=M^t} +\,\bar w^t,\\
r_f^{t,*}&=-\int_0^t \frac{\bar\eta^r}{\bar\kappa}
R_\theta^{t,r}r_f^{r-,*}\d z^r-\D_{w^*}\sigma^*(w^*,\eps),\\
r_f^t(x^{[t]})&={-} \int_0^t \frac{\bar\eta^r}{\bar\kappa}
R_\theta^{t,r}\big(r_f^{r-}(x^{[r]})+x^r\big)\d z^r.
\end{align*}
Noting that $r_f^{r-}(x^{[r]})$ and $r_f^{r-,*}$ are $\cF_t$-predictable,
taking expectations of the last two equations yields
\begin{align*}
R_f^{t,*}&={-}\int_0^t \bar \eta^r R_\theta^{t,r}
R_f^{r,*} \d r-\E[\D_{w^*} \sigma^*(w^*,\eps)],\\
R_f^t(x^{[t]})&={-} \int_0^t {\bar\eta^r}
R_\theta^{t,r}\big(R_f^{r}(x^{[r]})+x^{r}\big)\d r.
\end{align*}
These are Volterra integral equations of the second kind, and explicitly
solvable: Let $K_\theta^{t,s} \in \R^{k \times k}$
denote the unique resolvent kernel (see e.g.~\cite[Chap 2]{brunner2004collocation}) satisfying
\begin{equation}\label{eq:Volterra-Rt-resolvent}
K_\theta^{t,s} = -{\bar\eta^s} R_\theta^{t,s} - \int_s^t {\bar\eta^r} R_\theta^{t,r} K_\theta^{r,s} \d r
={-}\bar\eta^s R_\theta^{t,s}-\int_s^t \bar\eta^s K_\theta^{t,r}
R_\theta^{r,s}\d r.
\end{equation}
Then the explicit solutions for $R_f,R_f^*,\Gamma$ are
\begin{equation}\label{eq:Rf-linear}
R_f^t(x^{[t]}) = \int_0^t K_\theta^{t,s} x^s \d s,
\quad
R_f^{t,*}=-\bigg(1 + \int_0^t K_\theta^{t,s} \d
s\bigg)\E[\D_{w^*}\sigma^*(w^*,\eps)],
\quad \Gamma^t=\Id_k.
\end{equation}
To derive the form of $C_f$, note that the above Volterra integral equation
for $\bar\xi^t$ may also be solved to yield
\begin{equation}\label{eq:xiexplicit}
\bar\xi^t = M^t + \bar w^t + \int_0^t K_\theta^{t,s} (M^s + \bar w^s) \d s.
\end{equation}
We observe that since $\{\bar z^t\}_{t \geq 0}$ is a mean-zero process
independent of $(\{w^t\}_{t \geq 0},w^*,\eps)$, for any $t,s \geq 0$ this
implies
$\E[M^t \otimes \bar w^s] = \E[\E[M^t\mid \{w^t\}_{t\in [0,T]}, w^*, \eps]
\otimes \bar w^s ] = 0$. Thus
\begin{align*}
\bar C_\xi^{t,s}:=\E[\bar\xi^t \otimes \bar\xi^s]
&=\E\left[\left(M^t+\int_0^t K_\theta^{t,r}M^r \d r\right)
\otimes \left(M^s+\int_0^s K_\theta^{s,r}M^r \d r\right)\right]\\
&\hspace{1in}
+\E\left[\left(\bar w^t+\int_0^t K_\theta^{t,r}\bar w^r \d r\right)
\otimes \left(\bar w^s+\int_0^s K_\theta^{s,r}\bar w^r \d r\right)\right].
\end{align*}
Denote
\begin{equation}\label{eq:barCtheta}
\bar C_\theta^{t,s}=\E[\bar w^t \otimes \bar w^s]=\E[
(w^t-\sigma^*(w^*,\eps)) \otimes (w^s-\sigma^*(w^*,\eps))],
\quad (\{w^t\}_{t \geq 0},w^*) \sim \GP(0,C_\theta).
\end{equation}
Then the contribution from $\{\bar w^t\}$ is
\begin{align*}
&\E\brk{\prn{\bar w^t + \int_0^t K_\theta^{t,r}\bar w^r \d r} \otimes \prn{\bar w^s + \int_0^s K_\theta^{s,r'}\bar w^{r'} \d r'}} \\
&= \bar C_\theta^{t,s} + \int_0^t K_\theta^{t,r} \bar C_\theta^{r,s}\d r +
\int_0^s \bar C_\theta^{t,r'} K_\theta^{s,r'\top} \d r' + \int_0^t \int_0^s
K_\theta^{t,r}\bar C_\theta^{r,r'}K_\theta^{s,r'\top} \d r\,\d r'.
\end{align*}
For the contribution from $\{M^t\}$, observe that
\begin{align}
M^t + \int_0^t K_\theta^{t,r}M^r \d r &=  - \int_0^t
\frac{\bar\eta^s}{\bar\kappa} R_\theta^{t,s} \bar \xi^{s-}\d \bar z^s
- \int_0^t K_\theta^{t,s}\int_0^s \frac{\bar\eta^r}{\bar\kappa}
R_\theta^{s,r} \bar \xi^{r-}\d \bar z^r \d s\notag\\
&=  - \int_0^t  \prn{\frac{\bar\eta^r}{\bar\kappa} R_\theta^{t,r} + \int_r^t \frac{\bar\eta^r}{\bar\kappa} K_\theta^{t,s}   R_\theta^{s,r}\d s} \bar \xi^{r-}
\d \bar z^r
=\frac{1}{\bar\kappa}\int_0^t K_\theta^{t,r} \bar \xi^{r-}
\d \bar z^r.\label{eq:Midentity}
\end{align}
Then by the It\^o isometry,
\begin{align*}
\E\brk{\prn{M^t + \int_0^t K_\theta^{t,r}M^r \d r} \otimes \prn{M^s + \int_0^s K_\theta^{s,r}M^r \d r}}
&=\frac{1}{\bar\kappa}\int_0^{t\wedge s} K_\theta^{t,r} \E[\bar
\xi^r \otimes \bar \xi^r] K_\theta^{s,r\top} \d r.
\end{align*}
Combining these arguments, the covariance of $\{\bar \xi^t\}$ satisfies
\begin{equation}\label{eq:cov-barxi}
\begin{aligned}
\bar C_\xi^{t,s}&= \frac{1}{\bar\kappa}\int_0^{t\wedge s}
K_\theta^{t,r}\bar C_\xi^{r,r}
K_\theta^{s,r\top} \d r\\
&\hspace{0.2in}+\bar C_\theta^{t,s} + \int_0^t K_\theta^{t,r} \bar C_\theta^{r,s}\d r +
\int_0^s \bar C_\theta^{t,r'} K_\theta^{s,r'\top} \d r' + \int_0^t \int_0^s
K_\theta^{t,r}\bar C_\theta^{r,r'}K_\theta^{s,r'\top} \d r\,\d r'.
\end{aligned}
\end{equation}
This may be understood as a closed linear integral equation for the diagonal
$\{\bar C_\xi^{t,t}\}_{t \in [0,T]}$, whose solution may then be
substituted back to determine $\{\bar C_\xi^{t,s}\}_{s \neq t}$.
Finally, by definition
\begin{align*}
C_f^{t,s} &= \E\brk{\prn{\int_0^t {\bar\eta^r} \bar \xi^{r} \d r + \int_0^t
\frac{\bar\eta^r}{\bar\kappa} \bar\xi^{r-} \d \bar z^r}
\otimes \prn{\int_0^s {\bar\eta^r} \bar \xi^r \d r + \int_0^s
\frac{\bar\eta^r}{\bar\kappa} \bar\xi^{r-} \d \bar z^r}}\nonumber\\
&=(\mathrm{I}) + (\mathrm{II}) + (\mathrm{III}) + (\mathrm{IV}),
\end{align*}
where the terms are given by
\begin{align*}
(\mathrm{I}) & = 
\int_0^t \int_0^s \bar\eta^r\bar\eta^{r'}\E[\bar\xi^r \otimes \bar\xi^{r'}]\d r
\d r'=\int_0^t \int_0^s \bar\eta^r\bar\eta^{r'}\bar C_\xi^{r,r'}\d r
\d r',\\
(\mathrm{II}) & = \E\brk{\int_0^t \frac{\bar\eta^r}{\bar\kappa} \bar\xi^{r-} \d \bar
z^r \otimes \int_0^s \frac{\bar\eta^{r'}}{\bar\kappa} \bar\xi^{r'-} \d \bar
z^{r'}} = \int_0^{t\wedge s} \frac{(\bar\eta^{r})^2}{\bar\kappa}\bar
C_\xi^{r,r}\d r,\\
(\mathrm{III}) &= \int_0^t {\bar\eta^r}\E\brk{\bar \xi^r  \otimes \int_0^s
\frac{\bar\eta^{r'}}{\bar\kappa} \bar \xi^{r'-} \d \bar z^{r'}}\d r\\
&\overset{(*)}{=}\int_0^t {\bar\eta^r}\E\brk{  \prn{M^r+\int_0^r K_\theta^{r,r'}
M^{r'} \d r'}  \otimes \int_0^s \frac{\bar\eta^{r'}}{\bar\kappa} \bar \xi^{r'-} \d \bar z^{r'}}\d r\\
&\overset{(**)}{=}\int_0^t \frac{\bar\eta^r}{\bar\kappa}\E\brk{\int_0^r K_\theta^{r,r'}\bar
\xi^{r'-} \d \bar z^{r'}  \otimes \int_0^s \frac{\bar\eta^{r'}}{\bar\kappa}
\bar\xi^{r'-} \d \bar z^{r'}}\d r = \int_0^t \int_0^{r\wedge s}
\frac{\bar\eta^r\bar\eta^{r'}}{\bar\kappa}
K_\theta^{r,r'}\bar C_\xi^{r',r'} \d {r'}\d r,\\
(\mathrm{IV})&=\int_0^s \bar\eta^{r'} \E\brk{\int_0^t
\frac{\bar\eta^r}{\bar\kappa} \bar\xi^{r-}\d \bar z^r
\otimes \bar \xi^{r'}} \d s
=\int_0^s \int_0^{t \wedge r} 
\frac{\bar\eta^r\bar \eta^{r'}}{\bar\kappa} 
\bar C_\xi^{r',r'}K_\theta^{r,r'\top}\d r'\d r.
\end{align*}
In $(*)$, we have applied the form of $\{x^t\}_{t \geq 0}$ from
\eqref{eq:xiexplicit}, and the independence of $(\{\bar w^t\}_{t \geq
0},w^*,\eps)$ and $\{\bar z^t\}_{t \geq 0}$ as above. In $(**)$, we have applied
again the identity \eqref{eq:Midentity}, and the form for $(\mathrm{IV})$
follows from symmetric arguments as $(\mathrm{III})$. Thus, to summarize
\begin{equation}\label{eq:Cf-linear}
\begin{aligned}
C_f^{t,s}&=\int_0^t \int_0^s \bar\eta^r\bar\eta^{r'}\bar C_\xi^{r,r'}\d r \d r'
+\int_0^{t\wedge s} \frac{(\bar\eta^{r})^2}{\bar\kappa}\bar C_\xi^{r,r}\d r\\
&\hspace{1in}
+\int_0^t \int_0^{r\wedge s} \frac{\bar\eta^r\bar\eta^{r'}}{\bar\kappa}
K_\theta^{r,r'}\bar C_\xi^{r',r'} \d {r'}\d r
+\int_0^s \int_0^{t \wedge r} \frac{\bar\eta^r\bar \eta^{r'}}{\bar\kappa} 
\bar C_\xi^{r',r'}K_\theta^{r,r'\top}\d r'\d r.
\end{aligned}
\end{equation}
The mapping $(C_\theta,R_\theta) \mapsto (C_f, R_f, R_f^*, \Gamma)$ is thus
given by \eqref{eq:Volterra-Rt-resolvent},
\eqref{eq:Rf-linear}, \eqref{eq:barCtheta}, \eqref{eq:cov-barxi},
and \eqref{eq:Cf-linear}, and this mapping is the same for both SGD and SME.

Consider next the setting of a ridge regularizer,
\[G(\theta)=\frac{\lambda}{2}\|\theta\|^2,
\qquad g(\theta)=\lambda \theta,\]
and suppose that $R_f$ takes the form
$R_f^t(\theta^{[t]})=\int_0^t R_f^{t,s} \theta^s ds$.
Then \eqref{eq:def-cont-theta} and \eqref{eq:def-cont-deriv-t} also simplify
to linear Volterra integral equations,
\[\theta^t=\int_0^t A_f^{t,r} \theta^r \d r + b_f^t,
\qquad r_\theta^{t,s}=\int_s^t A_f^{t,r} r_\theta^{r,s} \d r + \Id_k,\]
where
\begin{align*}
        A_f^{t,r} &= -\bar{\eta}^r \left({\gamma} \Gamma^r + \lambda\,\Id_k\right) - \gamma \int_r^t {\bar{\eta}^s} R_f^{r,s} \d s, \qquad
    b_f^t = \theta^0 - \gamma \left(\int_0^t
{\bar{\eta}^r}R_f^{r,*} \d r\right)  \theta^* + \sqrt{\gamma}\,u^t.
\end{align*}
Let $K_f^{t,s} \in \R^{k \times k}$ be the unique resolvent kernel satisfying
\begin{equation}\label{eq:Volterra-Rf-resolvent}
K_f^{t,s}=A_f^{t,s} + \int_s^t A_f^{t,r} K_f^{r,s} \d r
=A_f^{t,s}+\int_s^t K_f^{t,r}A_f^{r,s}\d s.
\end{equation}
Then
\begin{equation}\label{eq:Rt-ridge}
R_\theta^{t,s}=\E r_\theta^{t,s}=\left(1+\int_s^t K_f^{t,r}\d r \right) \Id_k,
\end{equation} 
and
\[\theta^t = b_f^t + \int_0^t K_f^{t,s} b_f^s \,\d s\]
which implies that $C_\theta^{t,s}=\E[\theta^t \otimes \theta^s]$ is given by
\begin{equation}\label{eq:Ct-ridge}
\begin{aligned}
C_\theta^{t,s}&=C_b^{t,s} 
+ \int_0^t K_f^{t,r} C_b^{r,s} \,\mathrm{d}r 
+ \int_0^s C_b^{t,r} (K_f^{s,r})^\top \,\mathrm{d}r 
+ \int_0^t \int_0^s K_f^{t,r} C_b^{r,r'} (K_f^{s,r'})^\top \,\mathrm{d}r'
\,\mathrm{d}r,\\
C_\theta^{t,*}&=C_b^{t,*}+\int_0^t K_f^{t,s}
C_b^{s,*} \d s,\\
C_b^{t,s}&=\E[b_f^t\otimes b_f^s]=\E\left[\left(\theta^0 
- \gamma \left(\int_0^t
{\bar{\eta}^r}R_f^{r,*} \d r\right)\theta^*\right)\otimes
\left(\theta^0 - \gamma \left(\int_0^s
{\bar{\eta}^r}R_f^{r,*} \d r\right)\theta^*\right)\right] + \gamma\,C_f^{t,s},\\
C_b^{t,*}&=\E[b_f^t\otimes \theta^*]=\E\left[\left(\theta^0 
- \gamma \left(\int_0^t
{\bar{\eta}^r}R_f^{r,*} \d r\right)\theta^*\right) \otimes \theta^* \right].
\end{aligned}
\end{equation}
The mapping $(C_f, R_f, R^*_f, \Gamma) \mapsto (C_\theta, R_\theta)$ is given
by \eqref{eq:Volterra-Rf-resolvent}, \eqref{eq:Rt-ridge}, and
\eqref{eq:Ct-ridge}.

\subsubsection{Numerical simulation of the DMFT system}

\begin{figure}[b]
    \centering
    \includegraphics[width=0.75\linewidth]{Gaussian_n8000_d10000_lr0.8_lam0.1_eps0.0_Square_Ridge_Lin_1e-05_Summary_Avg.png}
    \caption{Comparison of SGD and SME dynamics in a linear model.
Parameters: $n=8000$, $d=10000$, $\eta=0.8$, $\lambda=0.1$. The dynamics of the
squared norm $\|\btheta_t\|^2/d$, overlap $\langle \btheta_t, \btheta^*
\rangle/d$, and training loss agree closely for SGD and SME, whereas those of
the empirical CDF $n^{-1}\sum_{i=1}^n \1\{|\x_i^\top \btheta| \le 1\}$ show a
discrepancy.}\label{fig:linear}
    \end{figure}

    \begin{figure}[t]
    \centering
    \includegraphics[width=0.75\linewidth]{n8000_d10000_lr3.0_lam0.1_eps0.0_Huber_Ridge_Tanh_0.0001_Summary_Avg.png}
    \caption{Divergence of SGD and SME dynamics in a model with Huber loss and
$\tanh$ activation.
Parameters: $n=8000$, $d=10000$, $\eta=3.0$, $\lambda=0.1$. Solid/dotted curves
depict the dynamics of SGD/SME and their DMFT predictions over 4
training epochs, exhibiting a divergence that is predicted by DMFT.
Dashed curves for training times $t>4$ depict the subsequent dynamics of
gradient flow initialized from the last SGD/SME iterate, in which the norm
and overlap statistics re-converge.}
    \label{fig:tanh}
\end{figure}

\begin{figure}[t]
    \centering
    \includegraphics[width=0.75\linewidth]{Gaussian_n8000_d10000_lr3.0_lam0.1_eps0.0_Huber_Ridge_Sin_0.0001_Summary_Avg.png}
    \caption{Divergence of SGD and SME dynamics in a model with Huber loss and
$\sin$ activation. Parameters are the same as Figure \ref{fig:tanh}. In this
example, the difference between SGD and SME persists through the later gradient
flow training, suggesting that the dynamics reach local optima
with distinct statistical properties.}
    \label{fig:sine}
    \end{figure}


We simulate the solution to the DMFT
system (\ref{eq:thetaeps}--\ref{eq:def-cont-Gamma}) by iterating the fixed-point
mappings $(C_\theta,R_\theta) \mapsto (C_f,R_f,R^*_f,\Gamma)$
and $(C_f,R_f,R^*_f,\Gamma) \mapsto (C_\theta,R_\theta)$ until convergence, and
then simulating the processes $\{\theta^t,\xi^t\}_{t \geq 0}$ from the
parameters of the fixed point. More concretely:
\begin{itemize}
\item In the setting of a ridge regularizer, the mapping
$(C_f,R_f,R^*_f,\Gamma) \mapsto (C_\theta,R_\theta)$ may be evaluated
analytically via \eqref{eq:Volterra-Rf-resolvent}, \eqref{eq:Rt-ridge}, and
\eqref{eq:Ct-ridge}.
\item In the setting of squared loss and linear activation, the mapping
$(C_\theta,R_\theta) \mapsto (C_f,R_f,R^*_f,\Gamma)$ may likewise be evaluated
analytically via \eqref{eq:Volterra-Rt-resolvent},
\eqref{eq:Rf-linear}, \eqref{eq:barCtheta}, \eqref{eq:cov-barxi},
and \eqref{eq:Cf-linear}.
\item Outside of these settings, the expectations
(\ref{eq:def-cont-C-t}--\ref{eq:def-cont-Gamma}) may be approximated
via Monte Carlo simulation of the processes
(\ref{eq:thetaeps}--\ref{eq:def-cont-resp-xi}).
\end{itemize}
All time integrals are discretized in increments of $\delta$, as described in
Section \ref{sec:discreteDMFT}. We remark that simulation of the processes
$\{\xi^t,r_f^{t,*},r_f^{t,s}\}$ for SGD is substantially faster
than for gradient flow or SME, as integrals against the Poisson
process $\{z^t\}_{t \geq 0}$ are rapidly computable due to the sparsity of its
jumps. A software implementation of these methods is available on the authors'
web page.

We close this section by presenting some numerical results that corroborate our
preceding discussion/findings. All simulations use a single-index model
$k=k^*=1$, single-sample batch size $\kappa=1$, time discretization $\delta=0.05$, 
and ridge regularizer $G(\theta)=\frac{\lambda}{2}\theta^2$ with
analytical computation of the DMFT mapping
$(C_f,R_f,R^*_f,\Gamma) \mapsto (C_\theta,R_\theta)$. The reverse DMFT mapping $(C_f,R_f,R^*_f,\Gamma) \mapsto (C_\theta,R_\theta)$ is
computed analytically for the linear model, and via 10,000 Monte Carlo samples
otherwise. DMFT predictions for $d^{-1}\|\btheta^t\|^2$ and $d^{-1}\btheta^{t\top}\btheta^*$ are contained in $C_\theta$, and the predictions for other observables are also computed via 10,000 Monte Carlo samples. Statistics of the high-dimensional SGD and SME processes
$\{\btheta^t\}_{t \geq 0}$ are averaged over 10 independent trials.\\

\paragraph{Coincidence of quadratic observables of SGD and SME in a linear model.}

For a linear model $y=\x^\top \btheta^*$ with loss $L(\hat
y,y)=\frac{1}{2}(\hat y-y)^2$ and activation $\sigma(\xi)=\xi$,
Figure~\ref{fig:linear} verifies the claim of Section \ref{sec:linear}
that the dynamics of the squared norm
$d^{-1}\|\btheta^t\|^2$, overlap $d^{-1}\btheta^{t\top}\btheta^*$,
and training loss $(2n)^{-1}\sum_{i=1}^n (\x_i^\top \btheta^t-\x_i^\top
\btheta^*)^2$ coincide for SGD and SME. The empirical distribution of $\{\x_i^\top
\btheta^t\}_{i \in [n]}$ is predicted (via the law of $\xi^t$) to be different
for SGD and SME, and we indeed observe a discrepancy at the tested
learning rate that is also predicted by the DMFT theory.\\


\paragraph{Difference between SGD and SME dynamics in non-convex landscapes.}

Figures \ref{fig:tanh} and \ref{fig:sine} display two different non-linear
and non-convex settings, in which the dynamics of SGD and SME
with relatively large learning rate $\eta=3.0$ are markedly different:

\begin{enumerate}
    \item Figure \ref{fig:tanh} uses a (Lipschitz-continuous) Huber loss
$L(\hat y,y)=\begin{cases} \frac{1}{2}(\hat y - y)^2 & \text{if } |\hat y - y| < 1 \\ (|\hat y - y| - \frac{1}{2}) & \text{otherwise} \end{cases}$ and activation $\sigma(x)=\tanh(x)$, with noisy true labels $y=\tanh(\x^\top \btheta^*+\eps)$ where $\eps\sim\mathcal{N}(0,0.1)$.
\item Figure \ref{fig:sine} uses the same Huber loss, activation
$\sigma(x)=\sin(x)$, and noisy true labels $y=\sin(\x^\top \btheta^*)+\eps$ where $\eps \sim\mathcal{N}(0,0.1)$.
\end{enumerate}

We highlight that the mechanism driving this difference may be different in
these two examples: Both examples train SGD/SME for $T=4$ epochs, and
then continue to apply gradient flow from the last SGD/SME iterate until
convergence. The dynamics over the first $T=4$ epochs in Figure \ref{fig:tanh}
diverge between SGD and SME, but converge again upon
the subsequent application of gradient flow. This suggests that SGD and SME
reach basins of attraction of (local) minimizers of the empirical risk with
similar statistical properties, and that their differences may be driven by
different oscillatory behaviors around these minimizers.
In contrast, the difference between SGD and SME in
Figure \ref{fig:sine} persists even through the later gradient flow training,
with SGD achieving a higher overlap $d^{-1}\btheta^{t\top}\btheta^*$
and smaller squared norm $d^{-1}\|\btheta^t\|^2$. This suggests that SGD and
SME reach local minimizers of the empirical risk that have different
generalization properties. This difference is reflected (over
the first 4 training epochs) by the DMFT theory.\\



\paragraph{The small learning rate limit.}

Figure \ref{fig:comp-lr} compares the overlap dynamics of SGD and SME in units
of rescaled time $\tau=\eta t$ across various learning rates $\eta$ for the Tanh
activation setting of Figure \ref{fig:tanh}, verifying the claim in Section
\ref{sec:comp-gf} that these dynamics
converge to a common (gradient flow) limit as $\eta \to 0$.\\

 %we explore this by varying the learning rate $\eta$ in a non-linear setting. At a small $\eta=0.5$, the overlap trajectories of SGD and SME are nearly identical, showing that the small learning rate leads to a behavior closer to the gradient flow in both SGD and SME settings, which matches our remark in Section \ref{sec:comp-gf}. As $\eta$ increases to $1.25$ and $2.5$, the trajectories diverge significantly, with SGD consistently achieving higher overlap. This highlights that in the high-dimensional regime with non-trivial learning rates ($\eta \sim O(1)$), the Gaussian approximation of SME fails to capture the beneficial effects of SGD's specific noise structure.

\begin{figure}[t]
    \centering
    \includegraphics[width=1\linewidth]{Horizontal_Overlap_Etas_20260128_124050.png}
    \caption{Effect of learning rate on SGD vs.\ SME divergence.
Overlap dynamics are plotted in units of rescaled time $\tau=\eta t$, for
various learning rates $\eta \in \{0.5, 1.25, 2.5\}$. At small learning rates
($\eta=0.5$), the SME diffusion provides an accurate approximation of SGD, and
their gap widens as $\eta$ increases.}
    \label{fig:comp-lr}
\end{figure}

\begin{figure}[t]
    \centering
    \includegraphics[width=0.5\linewidth]{SGD_DMFT_Compare_GammaExp_Gaussian_d4000_lr2.0_lam0.1_Huber_Tanh.png}
    \caption{Convergence to one-pass SGD in the large sample limit.
Overlap dynamics are plotted in units of rescaled time $\tau=\gamma t$, fixing
$d=4000$ and varying $\gamma=n/d \in \{0.5, 1.0, 5.0, 20.0\}$.
As the training sample size $\gamma$ increases, multi-pass SGD dynamics (solid
lines) and their DMFT predictions (dotted lines) both converge to the simulated
overlap statistics of one-pass SGD (black dashed line).}
    \label{fig:comp-onepass}
\end{figure}

\paragraph{Transition to one-pass SGD dynamics.}

Figure \ref{fig:comp-onepass} compares the overlap dynamics of SGD in units of
rescaled time $\tau=\gamma t$ for various dimension ratios $\gamma=\lim_{n,d \to
\infty} n/d$, verifying the claim in Section \ref{sec:comp-onepass}
that as $\gamma \to \infty$, these overlap curves converge to a limit that
represents one-pass SGD, and the training time at which these curves plateau
is delayed.

\section{Existence and uniqueness of the DMFT fixed
point}\label{sec:existenceuniqueness}

In the remainder of the paper, we prove Theorems \ref{thm:uniqueness-existence}
and \ref{thm:DMFT-limit}.

This section defines the spaces $\cS \equiv \cS(T,C_0)$ and $\cS^\text{cont}
\equiv \cS^\text{cont}(T,C_0)$, and proves \Cref{thm:uniqueness-existence}. The
proof is an adaptation of \cite{celentano2021high} for gradient flow (see also \cite{fan2025dynamical1} in a
setting of Langevin dynamics), and will consist of:
\begin{enumerate}
    \item Defining spaces $\mathcal{S}_\xi$
    for $(\Cf,\Rf,\Rf^*,\Gamma)$ and ${\mathcal{S}}_\theta$ for $(\Ct,\Rt)$,
and showing that on these spaces, the primary processes
$\theta^t,\xi^t$ and auxiliary response processes $r_\theta^{t,s},r_f^{t,*},r_f^t(x^{[t]})$ are well-defined.
    \item Showing that the transformations which represent the DMFT fixed point
equations for $(\Ct,\Rt,\Cf,\Rf,\Rf^*,\Gamma)$ are contractive in a suitable
metric.
\end{enumerate}

The main differences between our argument and those of \cite{celentano2021high} involve the
treatment of the stochastic Volterra-type integrals in \eqref{eq:def-cont-xi},
\eqref{eq:def-cont-deriv-xi-star}, and \eqref{eq:def-cont-resp-xi} that arise
due to the stochastic gradient noise in both SGD and SME; this necessitates 
modified definitions and treatments of the response processes, as discussed in
Remark \ref{rmk:response}.

\subsection{Definition of the space $\cS$}

We remind the reader that we abbreviate $\norm{\cdot}$ for the $\ell_2$-norm of
a vector and Frobenius norm of a matrix, and write $\normop{\cdot}$ for the
$\ell_2$-to-$\ell_2$ matrix operator norm.

For any Euclidean space $E$, we denote by
$L^4([0,t],E)$ the space of (Borel-measurable) functions
$s \in [0,t] \mapsto x^s \in E$ for which
\[\|x\|_{L^4}:=\left(\int_0^t \|x^s\|_2^4 \d s\right)^{1/4}<\infty.\]
For a linear operator $R:L^4([0,t],\R^k) \to \R^k$, we define also
its operator norm
\[\normop{R}=\sup_{x \in L^4([0,t],\R^k):\|x\|_{L^4} \leq 1} \norm{R(x)}_2.\]
We say that $R$ is a bounded linear operator if $\normop{R}<\infty$.

Fix any $T>0$. Under Assumptions \ref{asp:SGD} and \ref{asp:lipschitz}, there
exist constants $M_f,M_g,M_\eta>0$ such that
\begin{equation}\label{eq:defMbounds}
\begin{gathered}
\sup_{(\xi,w^*,\eps) \in \R^k \times \R^{k^*} \times \R}
\{\|f(\xi,w^*,\eps)\|,\|\D_\xi f(\xi,w^*,\eps)\|_\op,
\|\D_{w^*} f(\xi,w^*,\eps)\|_\op\} \leq M_f,\\
\sup_{\theta \in \R^k} \|\D g(\theta)\|_\op \leq M_g,
\qquad \sup_{t \in [0,T]} \bar \eta^t \cdot \max\{1,1/\sqrt{\bar\kappa}\}
\leq M_\eta.
\end{gathered}
\end{equation}
For a constant $C_0>0$ (which we will take large enough, depending on $T$ and
the above constants $M_f,M_g,M_\eta$),
we define envelope functions $\pRt,\pRf,\pRfs:[0,T] \to [0,\infty)$ by
\begin{align}
\frac{\d}{\d t}\pRt(t)&=C_0\prn{\pRt(t)+\int_0^t \pRf(t-r) \pRt(r) \d r},
\quad \text{ with } \pRt(0)=C_0,\label{eq:def-pRt}\\
\pRf(t) &= C_0\prn{\int_0^t \pRt(t-r) \pRf(r) \d r +
\pRt(t)},\label{eq:def-pRf}\\
\pRfs(t)&=C_0\left(1+\int_0^t \pRt(t-r)\pRfs(r)\d r\right).\label{eq:def-pRfs}
\end{align}
We then define an envelope function $\pCt:[0,T] \to [0,\infty)$ by
\begin{align}
\frac{\d}{\d t}\pCt(t)&=C_0\left(\pCt(t)
+\int_0^t \Phi_{R_f}(t-r)\pCt(r)\d r+\Phi_{R_f^*}(t)\right),
\text{ with } \pCt(0)=C_0(1+6\gamma).
\label{eq:def-pCt}
\end{align}
One may check using a Laplace transform argument that
these systems admit a unique solution $\pRt,\pRf,\pRfs,\pCt:[0,T] \to
[0,\infty)$, where all four envelope functions are nonnegative and increasing;
we refer to \cite[Lemma 3.1]{fan2025dynamical1} for details of such an argument.

\begin{defi}\label{def:S}
Fix $T>0$ and a constant $C_0>0$ (which also defines $\pRt,\pRf,\pRfs,\pCt$).
Given a finite subset $D=\{d_1,\ldots,d_m\} \subset [0,T]$, we call
\[[0,d_1),[d_1,d_2),\ldots,[d_{m-1},d_m),[d_m,T]\]
the maximal intervals of $[0,T] \setminus D$.

Let $\mathcal{S}_\xi \equiv \mathcal{S}_\xi(T,C_0)$ be the space of all
tuples $(C_f,R_f,R_f^*,\Gamma)$ such that there exists a finite set
$D \subset [0,T]$ for which:
    \begin{enumerate}
        \item $\Cf \equiv \{\Cf^{t,s}\}_{t,s \in [0,T]}$ is a symmetric covariance kernel
on $[0,T] \otimes \R^k$, i.e.\ $\Cf^{t,s}=(\Cf^{s,t})^\top
\in \R^{k \times k}$ for all $t,s \in [0,T]$,
and for any finite subset $T_0 \subset [0,T]$,
the matrix $(\Cf^{t,s})_{t,s \in T_0} \in \R^{k|T_0| \times k|T_0|}$ is
symmetric positive-semidefinite. Furthermore, we have,
\[\Cf^{0,0}=0, \qquad
\Tr \Cf^{t,t} \leq C_0 \text{ for all } 0\leq t\leq T,\]
and for each maximal interval $I$ of $[0,T]\setminus D$,
        \begin{equation}\label{eq:Cfcontinuity}
\Tr[\Cf^{t,t}-\Cf^{t,s}-\Cf^{s,t}+\Cf^{s,s}] \leq C_0|t-s|
\text{ for all } t,s \in I.
        \end{equation}
        \item $\Rf \equiv \{\Rf^t\}_{t \in [0,T]}$ is a family of bounded linear operators
$\Rf^t:L^4([0,t],\R^k) \to \R^k$.
Fixing any process $x \in L^4([0,T],\R^k)$, the map
$t \mapsto R_f^t(x^{[t]})$ is Borel-measurable, and
        \begin{equation}\label{eq:pRf-volterra}
\norm{\Rf^{t}(x^{[t]})}^2 \leq \int_0^t \pRf(t-r) \norm{x^r}^2\d r \text{ for all }
t\in [0,T].
\end{equation}
\item $\Rf^* \equiv \{\Rf^{t,*}\}_{t \in [0,T]}$ is a Borel-measurable
process of matrices in $\R^{k \times k^*}$, where
\[\|\Rf^{t,*}\|^2 \leq \pRfs(t) \text{ for all } t \in [0,T].\]
\item $\Gamma \equiv \{\Gamma^t\}_{t \in [0,T]}$ is a Borel-measurable process
of matrices in $\R^{k \times k}$, where
\[\normop{\Gamma_t} \leq M_f \text{ for all } t \in [0,T].\]
    \end{enumerate}

Let ${\mathcal{S}}_\theta \equiv {\mathcal{S}}_\theta(T,C_0)$ be the space of
all tuples $(C_\theta,R_\theta)$ such that there exists a finite set $D \subset [0,T]$
for which:
    \begin{enumerate}
        \item $\Ct \equiv \{\Ct^{t,s}\}_{t,s \in [0,T] \cup \{*\}}$ is a symmetric
covariance kernel on $([0,T] \otimes \R^k) \times \R^{k^*}$, i.e.\
$\Ct^{t,s}=(\Ct^{s,t})^\top \in \R^{k \times k}$,
$\Ct^{t,*}=(\Ct^{*,t})^\top \in \R^{k \times k^*}$, and
$\Ct^{*,*} \in \R^{k^* \times k^*}$ for all $t,s \in [0,T]$, and
for any finite subset $T_0 \subset [0,T]$,
the matrix $(\Ct^{t,s})_{t,s \in T_0 \cup \{*\}} \in \R^{(k|T_0|+k^*) \times
(k|T_0|+k^*)}$ is symmetric positive-semidefinite. Furthermore,
we have
\[\Ct^{*,*}=\E[\theta^*\theta^{*\top}],
\quad \Ct^{0,0}=\E[\theta^0\theta^{0\top}],
\quad \Ct^{0,*}=\E[\theta^0\theta^{*\top}],\]
\[\Tr \Ct^{t,t} \leq \pCt(t) \text{ for all } 0\leq t\leq T,\]
and for each maximal interval $I$ of $[0,T]\setminus D$,
        \begin{equation}\label{eq:Cthetacontinuity}
            \Tr[\Ct^{t,t}-\Ct^{t,s}-\Ct^{s,t}+\Ct^{s,s}] \leq 
(\Phi_{C_\theta}(T)+C_0^2)|t-s|
\text{ for all } t,s \in I.
        \end{equation}
        \item $\Rt \equiv \{\Rt^{t,s}\}_{t,s \in [0,T]}$ is a Borel-measurable
process of matrices in $\R^{k \times k}$, satisfying
\[\Rt^{t,s}=0 \text{ for all $s>t$},
\qquad
\|\Rt^{t,s}\|^2 \leq \pRt(t-s) \text{ for all } s \leq t,\]
and for each maximal interval $I$ of $[0,T] \setminus D$,
\[\|\Rt^{t',s}-\Rt^{t,s}\|^2 \leq \Phi_{R_\theta}(T)|t'-t|
\text{ for all } s \in [0,T] \text{ and }
t,t' \in I \text{ with } t,t' \geq s.\]
\end{enumerate}
We set
\[\cS \equiv \cS(T,C_0)=\{(C_\theta,R_\theta,C_f,R_f,R_f^*,\Gamma):
(C_\theta,R_\theta) \in \cS_\theta,\,(C_f,R_f,R_f^*,\Gamma) \in \cS_\xi\}.\]
We denote by
$\cS^{\mathrm{cont}},\cS_\xi^{\mathrm{cont}},\cS_\theta^{\mathrm{cont}}$
these spaces where the above conditions hold with $D=\emptyset$,
i.e.\ where the continuity conditions for
$C_f,C_\theta,R_\theta$ hold on the whole interval $[0,T]$. 
\end{defi}

Note that the continuity conditions \eqref{eq:Cfcontinuity} and
\eqref{eq:Cthetacontinuity} for $C_f$ and $C_\theta$ 
ensure (via Kolmogorov's continuity theorem) 
the existence of Gaussian processes $\{u^t\}_{t \in [0,T]} \sim \GP(0,C_f)$ 
and $(\{w^t\}_{t \in [0,T]},w^*) \sim \GP(0,C_\theta)$ with c\`adl\`ag sample
paths (having discontinuities only in $D$),
as claimed in the constructions of Section \ref{sec:DMFTlimits}.

\subsection{Existence and uniqueness of stochastic processes}
We first establish \Cref{thm:uniqueness-existence}(a), showing
the existence and uniqueness of solutions to the
stochastic integro-differential equations defined in
(\ref{eq:def-cont-theta}--\ref{eq:def-cont-resp-xi}), given any
$(C_\theta,R_\theta,C_f,R_f,R_f^*,\Gamma) \in \cS$. (Here
and throughout, a solution $\{Y^t\}_{t \in [0,T]}$ refers to a process that
satisfies the desired equation a.s.\ at each $t \in [0,T]$.
The solution is unique if any other solution $\widetilde Y$ is a modification of $Y$, i.e.\ $Y^t=\widetilde Y^t$ a.s.\ for each $t \in [0,T]$, implying in the case where $Y$ and $\widetilde Y$ are both c\`adl\`ag that they have equal sample paths a.s.)

The following general lemma shows the well-posedness of a class of
Volterra-type SDEs, which we will use to show existence and uniqueness
of the specific processes $\xi^t$, $r_f^{t,*}$, and $r_f^t(x^{[t]})$. 

\begin{lem}\label{lem:exist-unique-poisson-volterra}
Let $(\Omega,\cF,\{\cF\}_{t \geq 0},\P)$ be the filtered probability space
of Theorem \ref{thm:uniqueness-existence}(a), where
$\{z^t\}_{t \in [0,T]}$ is either the Poisson process
\eqref{eq:poisson-z} or the Gaussian diffusion process \eqref{eq:gaussian-z}.
Fix any $m \geq 1$, let $\{F^t\}_{t \in [0,T]}$ and $\{G^t\}_{t \in
[0,T]}$ be c\`adl\`ag $\cF_t$-adapted
processes in $\R^{m \times k}$ and $\R^m$, and let 
$H:\R^k \times \R^{k^*} \times \R \to \R^k$ be a function that is
uniformly Lipschitz-continuous in its first argument. Suppose,
for some deterministic constants $M>0$ and $\iota>0$ that
\[\sup_{t \in [0,T]} \norm{F^t} \leq M \text{ a.s.}, \qquad \sup_{t \in [0,T]} \E \norm{G^t}^{2+\iota}<\infty,
\qquad \E\|H(0,w^*,\eps)\|^{2+\iota}<\infty.\]
Then for any $\Rt \equiv \{\Rt^{t,s}\}_{t,s \in [0,T]}$
satisfying the conditions of Definition \ref{def:S},
there exists a unique $\cF_t$-adapted and c\`adl\`ag solution
$\{Y^t\}_{t \in [0,T]}$ to
\begin{equation}\label{eq:Picardconstruction}
Y^t=F^t \int_0^t \frac{\bar\eta^r}{\bar\kappa} R_\theta^{t,r}
H(Y^{r-},w^*,\eps)\d z^r+G^t.
\end{equation}
\end{lem}
\begin{proof}
Let $L^{2+\iota,\infty}(\Omega \times [0,T],\R^k)$ be the
space of jointly measurable and $\cF_t$-adapted processes $\{Y^t\}_{t \in [0,T]}$ (identified up to modifications)
for which
$\|Y\|_{L^{2+\iota,\infty}}:=\sup_{t \in [0,T]} (\E \|Y^t\|^{2+\iota})^{\frac{1}{2+\iota}}<\infty$.
Initialize $Y^{(0)}=G \in L^{2+\iota}(\Omega \times [0,T],\R^k)$, and
suppose inductively we have constructed a c\`adl\`ag process
$Y^{(n)} \in L^{2+\iota}(\Omega \times [0,T],\R^k)$.
Consider the Volterra-type It\^o integral
\[I^{(n),t}=\int_0^t \frac{\bar \eta^r}{\bar\kappa}
R_\theta^{t,r}H(Y^{(n),r-},w^*,\eps) \d z^r,\]
which defines a jointly measurable and adapted process $\{I^{(n),t}\}_{t \in [0,T]}$
by \cite[Lemma 2.A]{berger1980volterra} in the setting of Gaussian $\{z^t\}$ or
by a pathwise construction for each $\omega \in \Omega$ in the setting of
Poisson $\{z^t\}$. Recall the maximal intervals of $[0,T] \setminus D$ 
in Definition \ref{def:S} on which $t \in [r,T] \mapsto R_\theta^{t,r}$
is uniformly continuous. For Poisson
$\{z^t\}_{t \in [0,T]}$, it is clear that this continuity implies $I^{(n),t}$
is c\`adl\`ag over each such maximal interval, and hence also over
$[0,T]$, with discontinuities at the jumps of $\{z^t\}_{t \in [0,T]}$ and the
points of $D$. For Gaussian $\{z^t\}$, note that
for any $s<t$ belonging to the same maximal interval of $[0,T] \setminus D$,
we have by the Burkholder-Davis-Gundy inequality
(c.f.\ Lemma \ref{lem:burkholder}) and H\"older's inequality,
\begin{align*}
&\E\|I^{(n),t}-I^{(n),s}\|^{2+\iota}\\
&\lesssim \E\left\|\int_0^s \frac{\bar \eta^r}{\bar\kappa}
(R_\theta^{t,r}-R_\theta^{s,r})H(Y^{(n),r-},w^*,\eps) \d z^r\right\|^{2+\iota}
+\E\left\|\int_s^t \frac{\bar \eta^r}{\bar\kappa}
R_\theta^{t,r}H(Y^{(n),r-},w^*,\eps) \d z^r\right\|^{2+\iota}\\
&\lesssim \E\left(\int_0^s \|(R_\theta^{t,r}
-R_\theta^{s,r})H(Y^{(n),r},w^*,\eps)\|^2 \d r\right)^{1+\iota/2}
+\E\left(\int_s^t\|R_\theta^{t,r}H(Y^{(n),r},w^*,\eps)\|^2 \d
r\right)^{1+\iota/2}\\
&\lesssim \int_0^s \E \|(R_\theta^{t,r}
-R_\theta^{s,r})H(Y^{(n),r},w^*,\eps)\|^{2+\iota} \d r
+|t-s|^{\iota/2}\int_s^t \E\|R_\theta^{t,r}H(Y^{(n),r},w^*,\eps)\|^{2+\iota}
\d r.
\end{align*}
Then, by the boundedness and continuity conditions for $R_\theta$ in Definition
\ref{def:S} and the Lipschitz continuity of $H(\cdot)$,
\begin{align}
\E\|I^{(n),t}-I^{(n),s}\|^{2+\iota} &\lesssim
|t-s|^{2+\iota}\left(\E\|H(0,w^*,\eps)\|^{2+\iota}+\int_0^T
\E \|Y^{(n),r}\|^{2+\iota} \d r\right)\notag\\
&\qquad+|t-s|^{\iota/2}\left(|t-s|\,\E\|H(0,w^*,\eps)\|^{2+\iota}
+\int_s^t \|Y^{(n),r}\|^{2+\iota} \d r\right)\notag\\
&\lesssim |t-s|^{2+\iota}+|t-s|^{1+\iota/2}.\label{eq:Icontinuous}
\end{align}
Kolmogorov's continuity theorem then implies that $\{I^{(n),t}\}_{t \in [0,T]}$ 
has a modification which is uniformly continuous on each maximal interval of
$[0,T] \setminus D$ and hence c\`adl\`ag on $[0,T]$.

Then defining
\begin{equation}\label{eq:PicardStepXi}
Y^{(n+1),t}=F^tI^{(n),t}+G^t
=F^t\int_0^t \frac{\bar\eta^r}{\bar\kappa} R_\theta^{t,r}
H(Y^{(n),r-},w^*,\eps)\d z^r+G^t,
\end{equation}
$Y^{(n+1)}$ is adapted and c\`adl\`ag.
Applying the boundedness of $\{F^t\}$, the Burkholder-Davis-Gundy type
inequality of Lemma \ref{lem:burkholder} in both the Poisson and Gaussian
settings, and H\"older's inequality as above,
\begin{align*}
\E \norm{Y^{(n+1),t}}^{2+\iota}
&\lesssim \E \left\|\int_0^t \frac{\bar\eta^r}{\bar\kappa} R_\theta^{t,r}
H(Y^{(n),r-},w^*,\eps)\d z^r\right\|^{2+\iota}
+\E \|G^t\|^{2+\iota}\\
&\lesssim \E\|H(0,w^*,\eps)\|^{2+\iota}
+\int_0^T \E\|Y^{(n),r}\|^{2+\iota}\d r
+\E\|G^t\|^{2+\iota}.
\end{align*}
Thus $Y^{(n+1)} \in L^{2+\iota,\infty}(\Omega \times [0,T],\R^k)$, so
$Y^{(n+1)}$ is inductively well-defined for each $n \geq 0$.

Now consider the difference $Y^{(n),t}-Y^{(n-1),t}$,
and set $\phi_n(t)=\E\|Y^{(n),t}-Y^{(n-1),t}\|^{2+\iota}$.
Then again by Lemma \ref{lem:burkholder} and similar arguments as above,
\begin{align}
\phi_1(t)&\lesssim \E \left\|\int_0^t \frac{\bar\eta^r}{\bar\kappa}
R_\theta^{t,r}H(Y^{(0),r-},w^*,\eps) \d z^r\right\|^{2+\iota}
\lesssim \E\|H(0,w^*,\eps)\|^{2+\iota} +\int_0^T \E \|G^r\|^{2+\iota} \d
r,\notag\\
\phi_{n+1}(t) &\lesssim
\E \left\|\int_0^t \frac{\bar\eta^r}{\bar\kappa}R_{\theta}^{t,r}
\Big(H(Y^{(n),r-},w^*,\eps) - H(Y^{(n-1),r-},w^*,\eps)\Big)
\d z^r\right\|^{2+\iota}\notag\\
&\lesssim\int_0^t \E \|Y^{(n),r}-Y^{(n-1),r}\|^{2+\iota}
\d r=\int_0^t \phi_n(r)\d r.\label{eq:Picardbound}
\end{align}
I.e., there exists a constant $C>0$ (depending on $M,T,\pRt(T),M_\eta,H$,
and the law of $(w^*,\eps)$) for which
$\phi_1(t) \leq C$ and $\phi_{n+1}(t)\leq C\int_0^t \phi_n(r)\d r$.
Then $\phi_{n}(t)\leq C^nt^{n-1}/(n-1)!$, implying that $\{Y^{(n)}\}_{n \geq 0}$
is Cauchy in $L^{2+\iota,\infty}(\Omega \times [0,T],\R^k)$ and thus has a
limit $Y \in L^{2+\iota,\infty}(\Omega \times [0,T],\R^k)$ such that
$\lim_{n \to \infty} \sup_{t \in [0,T]} \E\|Y^{(n),t}-Y^t\|^{2+\iota}=0$. Defining the process
\[I^t=\int_0^t \frac{\bar\eta^r}{\bar\kappa}R_\theta^{t,r}
H(Y^{r-},w^*,\eps)\d z^r\]
as above, we then have $\E\|Y^t-(F^tI^t+G^t)\|^{2+\iota}
=\lim_{n \to \infty} \E\|Y^{(n+1),t}-(F^tI^{(n),t}+G^t)\|^{2+\iota}=0$ for each $t \in [0,T]$,
so $\{Y^t\}_{t \in [0,T]}$ solves \eqref{eq:Picardconstruction}. The
same argument as \eqref{eq:Icontinuous} implies that 
$\{I^t\}_{t \in [0,T]}$ and $\{Y^t\}_{t \in [0,T]}$ have c\`adl\`ag
modifications.

For uniqueness, let $\{\widetilde Y^t\}_{t \in [0,T]}$ be any other adapted
and c\`adl\`ag solution
to \eqref{eq:Picardconstruction}. Fix any $K>0$, define
the stopping time $\tau=\inf\{t \in [0,T]:\widetilde Y^t \geq K\}$,
and consider the stopped process 
$(\widetilde Y^\tau)^t=\widetilde Y^{\tau \wedge t}$. Then
by Lemma \ref{lem:burkholder} and a similar argument as above,
\begin{align*}
\E\|(\widetilde Y^\tau)^t\|^{2+\iota}
&=\E\left\|F^{t \wedge \tau}
\int_0^t \1\{r \leq \tau\}\frac{\bar \eta^r}{\bar\kappa}
R_\theta^{t,r}H(\widetilde Y^{r-},w^*,\eps)\d z^r
+G^{t \wedge \tau}\right\|^{2+\iota}\\
&\lesssim 1+\int_0^t
\E \|\1\{r \leq \tau\}\widetilde Y^{r-}\|^{2+\iota}\d r
\lesssim 1+\int_0^t \E\|(\widetilde Y^\tau)^r\|^{2+\iota}\d r.
\end{align*}
Note that the constant underlying $\lesssim$ is independent of $K$, so Gr\"onwall's
inequality implies that $\sup_{t \in [0,T]} \E\|(\widetilde Y^\tau)^t\|^{2+\iota}$ is uniformly
bounded independently of $K$. Taking $K \to \infty$
shows that $\widetilde Y \in L^{2+\iota,\infty}(\Omega \times [0,T],\R^k)$.
Then defining $\phi(t)=\E \|Y^t-\widetilde Y^t\|^{2+\iota}$,
the same argument as \eqref{eq:Picardbound} shows
$\phi(t) \lesssim \int_0^t \phi(s)\d s$, so $\sup_{t \in [0,T]} \phi(t)=0$ by Gr\"onwall's inequality, and $\widetilde Y$ is a modification of $Y$.
\end{proof}

\begin{proof}[Proof of Theorem \ref{thm:uniqueness-existence}(a)]
Let $(C_\theta,R_\theta,C_f,R_f,R_f^*,\Gamma) \in \cS$ be given, and let
$(\theta,\theta^*)$, $\eps$, $\{u^t\}_{t \in [0,T]}$, $(\{w^t\}_{t \geq
[0,T]},w^*)$, and $\{z^t\}_{t \in [0,T]}$ in the filtered space
$(\Omega,\cF,\{\cF\}_{t \in [0,T]},\P)$ also be given. The bound for
$C_f$ in Definition \ref{def:S} implies that $\sup_{t \in [0,T]} \E\|w^t\|^{2+\iota}<\infty$ for any $\iota>0$.
Then there exists a unique $\cF_t$-adapted and c\`adl\`ag
solution $\{\xi^t\}_{t \in [0,T]}$ to \eqref{eq:def-cont-xi} by
\Cref{lem:exist-unique-poisson-volterra}. Given this solution and any
deterministic process $x \in L^4([0,T],\R^k)$, define next
\[G^t={-}\D_\xi
f(\xi^t,w^*,\eps)\int_0^t \frac{\bar\eta^r}{\bar\kappa}R_\theta^{t,r}
\D_\xi f(\xi^{r-},w^*,\eps)x^r \d z^r.\]
This process $\{G^t\}_{t \in [0,T]}$ is c\`adl\`ag and satisfies $\sup_{t \in [0,T]}\|G^t\|^{2+\iota}<\infty$ for
sufficiently small $\iota>0$, by
the same arguments as in \Cref{lem:exist-unique-poisson-volterra}.
Then there exist unique $\cF_t$-adapted and c\`adl\`ag solutions
$\{r_f^{t,*}\}_{t \in [0,T]}$ and
$\{r_f^t(x^{[t]})\}_{t \in [0,T]}$ to
(\ref{eq:def-cont-deriv-xi-star}--\ref{eq:def-cont-resp-xi})
also by \Cref{lem:exist-unique-poisson-volterra}.

For \eqref{eq:def-cont-theta}, set $\tilde
\theta^t=\theta^t-\sqrt{\gamma}\,u^t$. Note that since $C_f^{0,0}=0$, we have
$u^0=0$ and $\tilde \theta^0=\theta^0$. Then
$\{\theta^t\}_{t \in [0,T]}$ solves
\eqref{eq:def-cont-theta} if and only if 
$\{\tilde \theta^t\}_{t \in [0,T]}$ solves the differential equation
\begin{align}
\frac{\d}{\d t}\tilde\theta^t&=\underbrace{-\bar\eta^t \Big({\gamma}\Gamma^t
(\tilde\theta^t + \sqrt{\gamma}\,u^t)+g(\tilde\theta^t + \sqrt{\gamma}\,
u^t)+{\gamma}R_f^t(\tilde\theta^{[t]})+ {\gamma^{3/2}}
R_f^t(u^{[t]})+{\gamma}R_f^{t,*}\theta^*\Big)}_{:=F(t,\tilde\theta^{[t]})},
\quad \tilde\theta^0=\theta^0.\label{eq:cont-theta-rewrite}
\end{align}
The bound \eqref{eq:pRf-volterra} implies that
$R_f^t$ restricts to an operator on $C([0,t],\R^k) \subset L^4([0,t],\R^k)$ that
is Lipschitz in the sup-norm $\|x\|_\infty=\sup_{s \in [0,t]} \|x^s\|_2$,
uniformly over $t \in [0,T]$. Then the function $(t,x) \mapsto F(t,x^{[t]})$
on the right side above is measurable
in $t$ and uniformly Lipschitz in $x \in C([0,T],\R^k)$, implying
(c.f.\ \cite[Theorems 2.1, 2.3]{hale2013introduction})
the existence of a unique continuous solution to \eqref{eq:cont-theta-rewrite} 
for any realization of $(\theta^0,\theta^*)$ and
$\{u^t\}_{t \in [0,T]}$. This defines, pathwise for each $\omega \in \Omega$,
a unique $\cF_t$-adapted and c\`adl\`ag solution
$\theta^t=\tilde \theta^t+\sqrt{\gamma}\,u^t$ to \eqref{eq:def-cont-theta}, with discontinuities at the jumps of $\{u^t\}$.
Given this solution $\{\theta^t\}_{t \in [0,T]}$, the same argument shows that
for each fixed $s \in [0,T]$, there is a unique continuous solution
$\{r_\theta^{t,s}\}_{t \in [s,T]}$ to
\[\frac{\d}{\d t}r_\theta^{t,s}
={-}\bar\eta^t[(\gamma \Gamma^t+\D g(\theta^t)r_\theta^{t,s}
+\gamma R_f^t(r_\theta^{[t],s})],
\quad r_\theta^{s,s}=\Id_k,\]
which is equivalent to \eqref{eq:def-cont-deriv-t}.
\end{proof}

\subsection{Closure of the mappings}

Given $(\Ct,\Rt) \in \cS_\theta$, consider the processes \eqref{eq:def-cont-xi},
\eqref{eq:def-cont-deriv-xi-star}, and \eqref{eq:def-cont-resp-xi}
defined on $(\Omega,\cF,\{\cF_t\}_{t \geq 0},\P)$ and let
$\trsfrmB(\Ct,\Rt)=(\Cf,\Rf,\Rf^*,\Gamma)$ be the mapping defined by
(\ref{eq:def-cont-C-h}--\ref{eq:def-cont-Gamma}). Similarly, given
$(\Cf,\Rf,\Rf^*,\Gamma) \in \cS_\xi$, consider the processes
\eqref{eq:def-cont-theta} and \eqref{eq:def-cont-deriv-t} and let
$\trsfrmA(\Cf,\Rf,\Rf^*,\Gamma)=(\Ct,\Rt)$ be the mapping defined by
(\ref{eq:def-cont-C-t}--\ref{eq:def-cont-R-t}).

In this section, we check that for a sufficiently large choice of constant
$C_0>0$ defining $\cS_\theta$ and $\cS_\xi$,
$\trsfrmB$ maps ${\mathcal{S}}_\theta$ into
$\mathcal{S}^{\mathrm{cont}}_\xi \subset \cS_\xi$, and 
$\trsfrmA$ maps $\mathcal{S}_\xi$ into ${\mathcal{S}}_\theta$.

\begin{lem} \label{lem:solution-in-space}
Fix any $T>0$. Then for any constant $C_0>0$ large enough (depending on $T$,
$M_f$, $M_g$, $M_\eta$, $\E\|\theta^0\|^2$, $\E\|\theta^*\|^2$, and
$\|g(0)\|^2$),
$\trsfrmA$ maps $\mathcal{S}_\xi(T,C_0)$ into ${\mathcal{S}}_\theta(T,C_0)$, and
$\mathcal{S}^{\mathrm{cont}}_\xi(T,C_0)$ into
${\mathcal{S}}^{\mathrm{cont}}_\theta(T,C_0)$.
\end{lem}
\begin{proof}
Assume $(C_f,R_f,R_f^*,\Gamma) \in S_\xi(T,C_0)$, and denote
$(C_\theta,R_\theta)=\trsfrmA(C_f,R_f,R_f^*,\Gamma)$.

We note that by its definition in \eqref{eq:def-cont-C-t},
the kernel
$\{C_\theta^{t,s}\}_{t,s \in [0,T] \cup \{*\}}$ is positive-semidefinite 
with $C_\theta^{*,*},C_\theta^{0,*},C_\theta^{0,0}$ as specified in Definition
\ref{def:S}. Recall from \eqref{eq:def-cont-theta} that
\begin{align}\label{eq:bound-theta-second-moment-1}
    \theta^t&=\theta^0-\int_0^t \bar\eta^r \Big({\gamma}\Gamma^r
\theta^r+g(\theta^r)+{\gamma}R_f^{r}(\theta^{[r]})+
{\gamma}R_f^{r,*}\theta^*\Big)\d r+\sqrt{\gamma}\,u^t.
\end{align}
Applying Cauchy-Schwarz and the bound $(\int_0^t f(s)\d s)^2 \leq t\int_0^t
f(s)^2 \d s \leq T\int_0^t f(s)^2 \d s$, we have $\E\|\theta^t\|^2 \leq
6(\nI+\nII+\nIII+\nIV+\nV+\nVI)$ where
\begin{align*}
\nI&=\E\|\theta^0\|^2\\
\nII&=\E\left\|\int_0^t
\bar\eta^r\gamma\Gamma^r\theta^r \d r\right\|^2
\leq \gamma^2 M_f^2 M_\eta^2 T\int_0^t \E\norm{\theta^r}^2 \d r,\\
\nIII&=\E\left\|\int_0^t \bar\eta^r g(\theta^r) \d
r\right\|^2\leq M_\eta^2 T\int_0^t \E\|g(\theta^r)\|^2 \d r
    \leq 2M_g^2 M_\eta^2 T\int_0^t \E\norm{\theta^r}^2\d r
+2M_\eta^2 T^2\norm{g(0)}^2,\,\\
\nIV&=\E\left\|\int_0^t \bar\eta^r\gamma
R_f^{r}(\theta^{[r]}) \d r\right\|^2 \leq \gamma^2 M_\eta^2 T\int_0^t \E\norm{R_f^{r}(\theta^{[r]})}^2 \d r\\
    &\overset{(*)}{\leq} \gamma^2 M_\eta^2 T\int_0^t
\int_0^r \pRf(r-r')\E \|\theta^{r'}\|^2\d r'\,\d r,\\
\nV&=\E\left\|\int_0^t \bar\eta^r \gamma R_f^{r,*}\theta^*
\d r\right\|^2
    \leq \gamma^2 M_\eta^2 T\,\E\norm{\theta^*}^2
\int_0^t \Phi_{R_f^*}(r)\d r,\\
\nVI&=\gamma\,\E\norm{u^t}^2=\gamma \Tr \Cf^{t,t} \leq \gamma C_0.
\end{align*}
Here, we have used the property \eqref{eq:pRf-volterra} of $R_f^t(\cdot)$
to bound $(*)$.
Collecting these bounds shows, for any constant $C_0>0$ large enough,
\begin{align*}
\E\|\theta^t\|^2&\leq C_0\left(1+6\gamma+\int_0^t \left(\E\|\theta^r\|^2
+\int_0^r \Phi_{R_f}(r-r')\E\|\theta^{r'}\|^2\d r'+\Phi_{R_f^*}(r) \right)\d
r\right).
\end{align*}
The definition of $\pCt$ in \eqref{eq:def-pCt} implies
\begin{equation}\label{eq:PhiCthetaexplicit}
\pCt(t)=C_0\left(1+6\gamma+\int_0^t \left(\pCt(r)
+\int_0^r \Phi_{R_f}(r-r')\pCt(r')\d r'+\Phi_{R_f^*}(r) \right)\d r\right),
\end{equation}
so comparing with the above shows
\[\Tr C_\theta^{t,t}=\E\|\theta^t\|^2 \leq \Phi_{\Ct}(t).\]
Next, let $D$ be the discontinuity set of $(\Cf,\Rf,\Rf^*,\Gamma) \in
\cS_\xi(T,C_0)$.
For each maximal interval $I$ of $[0,T] \setminus D$
and any $s,t\in I$, by an analogous application of Cauchy-Schwarz
and this bound $\E\|\theta^t\|^2 \leq \Phi_{C_\theta}(t)$, for any constant
$C_0>0$ large enough,
\begin{align}
&\Tr[\Ct^{t,t} - \Ct^{t,s} - \Ct^{s,t} + \Ct^{s,s}]
=\E \norm{\theta^t-\theta^s}^2\notag\\
&=\E\left\|\int_s^t \bar\eta^r \Big({\gamma}\Gamma^r
\theta^r+g(\theta^r)+{\gamma}R_f^{r}(\theta^{[r]})+
{\gamma}R_f^{r,*}\theta^*\Big)\d r+\sqrt{\gamma}(u^t-u^s)\right\|^2\notag\\
&\leq C_0|t-s|\left(1+\int_s^t \E\|\theta^r\|^2\d r
+\int_s^t\int_0^r \Phi_{R_f}(r-r')\E\|\theta^{r'}\|^2\d r'\d r
+\int_s^t\Phi_{R_f^*}(r)\d r\right)+C_0\E\|u^t-u^s\|^2\notag\\
&\leq C_0|t-s|\left(1+\int_0^T \Phi_{C_\theta}(r)\d r
+\int_0^T \int_0^r \Phi_{R_f}(r-r')\Phi_{C_\theta}(r')\d r' \d r
+\int_0^T \Phi_{R_f^*}(r)\d r\right)
+C_0\E\|u^t-u^s\|^2\notag\\
&\leq (\Phi_{C_\theta}(T)+C_0^2)|t-s|\label{eq:bound-cont-Ct}
\end{align}
where the last inequality applies
$\E\norm{u^t-u^s}^2=\Tr[\Cf^{t,t}-\Cf^{t,s}-\Cf^{s,t}+\Cf^{s,s}]
\leq C_0|t-s|$ by \eqref{eq:Cfcontinuity}, and the form of
$\Phi_{\Ct}(T)$ in \eqref{eq:PhiCthetaexplicit}.
This checks all conditions of Definition \ref{def:S} for $C_\theta$.

For $R_\theta$, recall from \eqref{eq:def-cont-deriv-t} that
\begin{align}\label{eq:rthetarestate}
r_\theta^{t,s}=\Id-\int_s^t \bar\eta^r\bigg[\Big({\gamma}\Gamma^r+\D
g(\theta^r)\Big)r_\theta^{r,s}+{\gamma}
R_f^r\Big(r_\theta^{[r],s}\Big)\bigg]\d r \text{ for } t \geq s.
\end{align}
Then $\|r_\theta^{t,s}\|^2 \leq 4(\nI+\nII+\nIII+\nIV)$ where
\begin{align*}
\nI&=\|\Id\|^2=k\\
\nII&=\left\|\int_s^t \bar\eta^r\gamma\Gamma^r r_\theta^{r,s}\d
r\right\|^2
\leq \gamma^2M_f^2M_\eta^2T\int_s^t \|r_\theta^{r,s}\|^2 \d r,\\
\nIII&=\left\|\int_s^t \bar\eta^r \D g(\theta^r)r_\theta^{r,s}\d
r\right\|^2
\leq M_g^2 M_\eta^2 T\int_s^t \|r_\theta^{r,s}\|^2 \d r,\\
\nIV&=\left\|\int_s^t \bar\eta^r \gamma R_f^r\Big(r_\theta^{[r],s}\Big)\d
r\right\|^2
\leq \gamma^2 M_\eta^2 T\int_s^t
\Big\|R_f^r\Big(r_\theta^{[r],s}\Big)\Big\|^2 \d r\\
&\leq \gamma^2 M_\eta^2 T\int_s^t \int_s^r \Phi_{R_f}(r-r')
\|r_\theta^{r',s}\|^2 \d r'\d r,\\
\end{align*}
the last inequality again using \eqref{eq:pRf-volterra} and the fact that
$r_\theta^{r',s}=0$ for $r'<s$. Collecting these bounds shows
\[\|r_\theta^{t,s}\|^2
\leq C_0\left(1+\int_s^t \left(\norm{r_\theta^{r,s}}^2
+\int_s^r \Phi_{R_f}(r-r')\|r_\theta^{r',s}\|^2\d r'\right) \d r\right)\]
for any constant $C_0>0$ large enough.
The definition of $\Phi_{R_\theta}(t)$ in \eqref{eq:def-pRt} implies
\begin{equation}\label{eq:PhiRthetaexplicit}
\Phi_{R_\theta}(t-s)=C_0\left(1+\int_0^{t-s}\left(
\Phi_{R_\theta}(r)+\int_0^r \Phi_{R_f}(r-r')\Phi_{R_\theta}(r')\d
r'\right)\d r\right),
\end{equation}
so comparing with the above shows
\[\|R_\theta^{t,s}\|^2=\|\E r_\theta^{t,s}\|^2\leq \E\|r_\theta^{t,s}\|^2
\leq \Phi_{R_\theta}(t-s).\]
Then, applying this bound
$\E\|r_\theta^{t,s}\|^2 \leq \Phi_{R_\theta}(t-s)$,
we have similarly for any $C_0>0$ large enough and any $t \geq t' \geq s$ that
\begin{align}
\|R_\theta^{t,s}-R_\theta^{t',s}\|^2
&\leq \E\|r_\theta^{t,s}-r_\theta^{t',s}\|^2\notag\\
&\leq C_0|t-t'|\left(\int_{t'}^t \E\|r_\theta^{r,s}\|^2 \d r
+\int_{t'}^t \int_s^r \Phi_{R_f}(r-r')\E \|r_\theta^{r',s}\|^2 \d r' \d
r\right)\notag\\
&\leq C_0|t-t'|\left(\int_s^t \pCt(r-s)\d r
+\int_s^t \int_s^r \Phi_{R_f}(r-r')\pCt(r'-s)\d r' \d
r\right)\notag\\
&\leq |t-t'| \cdot \pRt(t-s) \leq |t-t'| \cdot \pRt(T),\label{eq:Rcontbound}
\end{align}
where the last two inequalities apply \eqref{eq:PhiRthetaexplicit} and
monotonicity of $\pRt$.
This checks all conditions of Definition \ref{def:S} for $R_\theta$, so
$(\Ct,\Rt) \in {\mathcal{S}}_\theta(T,C_0)$.

We note that \eqref{eq:Rcontbound} holds for all $t,t',s \in [0,T]$ with
$t \geq t' \geq s$, irrespective of the discontinuity set
$D$ for $(\Cf,\Rf,\Rf^*,\Gamma) \in \cS_\xi(T,C_0)$.
If $(\Cf,\Rf,\Rf^*,\Gamma)\in \mathcal{S}^{\mathrm{cont}}_\xi(T,C_0)$, then
\eqref{eq:bound-cont-Ct} also holds for all $t,s \in [0,T]$, implying by the
above arguments that the conditions of Definition \ref{def:S} for
$(C_\theta,R_\theta)$ hold with $D=\emptyset$.
Thus also $(\Ct,\Rt) \in {\mathcal{S}}_\theta^\mathrm{cont}(T,C_0)$.
\end{proof}

\begin{lem} \label{lem:solution-in-space-2}
Fix any $T>0$. Then for any constant $C_0>0$ large enough (depending on $T$,
$M_f$, and $M_\eta$), $\trsfrmB$ maps ${\mathcal{S}}_\theta(T,C_0)$ into
$\mathcal{S}^{\mathrm{cont}}_\xi(T,C_0)$.
\end{lem}
\begin{proof}
In both the Poisson and Gaussian settings,
$\{z^t-\bar \kappa t\}_{t \in [0,T]}$ is a
martingale whose predictable quadratic variation is given by
$\langle z^t-\bar \kappa t
\rangle=\bar \kappa t$. Then for any $\cF_t$-predictable square-integrable
process $\{X^t\}_{t \in [0,T]}$,
\begin{align}
\E\left\|\int_0^t X^r \d z^r\right\|_2^2
&\leq 2\,\E\left\|\int_0^t X^r \bar \kappa\,\d r\right\|_2^2
+2\,\E\left\|\int_0^t X^r (\d z^r-\bar \kappa\,\d r)\right\|_2^2\notag\\
&=2\,\E\left\|\int_0^t X^r \bar \kappa\,\d r\right\|_2^2
+2\int_0^t \E\|X^r\|_2^2\,\bar\kappa\,\d r
\leq 2(t+1)\int_0^t \E\|X^r\|_2^2\,\bar\kappa\,\d r.\label{eq:basicdzrbound}
\end{align}

Assume $(\Ct,\Rt) \in {\mathcal{S}}_\theta(T,C_0)$, and denote
$(C_f,R_f,R_f^*,\Gamma)=\trsfrmB(C_\theta,R_\theta)$.
By its definition in \eqref{eq:def-cont-C-h},
$C_f$ is a positive-semidefinite kernel with $C_f^{0,0}=0$.
Applying \eqref{eq:basicdzrbound}, for any $t \in [0,T]$,
\[\Tr \Cf^{t,t}=\E\left\|\int_0^t
\frac{\bar\eta^r}{\bar\kappa}f(\xi^{r-},w^*,\eps)\d z^r\right\|^2
\leq 2(t+1)\int_0^t \E\left\|\frac{\bar \eta^r}{\bar \kappa}
f(\xi^r,w^*,\eps)\right\|^2\bar \kappa\,\d r
\leq C_0\]
for a large enough constant $C_0>0$. Similarly, for all $t,s \in [0,T]$,
\begin{align*}
\Tr[\Cf^{t,t} - \Cf^{t,s} - \Cf^{s,t} + \Cf^{s,s}] &=\E\left\|\int_s^t
\frac{\bar\eta^r}{\bar\kappa}f(\xi^{r-},w^*,\eps)\d z^r\right\|^2
\leq C_0|t-s|
\end{align*}
for large enough $C_0>0$.
This checks all conditions of Definition \ref{def:S} for $C_f$.

For $R_f$, given any $x \in L^4([0,T],\R^k)$, 
it is clear from the definition of $r_f^t(x^{[t]})$
in \eqref{eq:def-cont-resp-xi} that
$t \mapsto R_f^t(x^{[t]})=\E r_f^t(x^{[t]})$ is Borel-measurable.
By \eqref{eq:def-cont-resp-xi} and \eqref{eq:basicdzrbound},
we have
\begin{align}
\E\|r_f^t(x^{[t]})\|^2 &\leq M_f^2
\E\left\|\int_0^t \frac{\bar\eta^r}{\bar\kappa}
R_\theta^{t,r}\Big(r_f^{r-}(x^{[r]})+\D_\xi f(\xi^{r-},w^*,\eps)x^{r-}\Big)\d
z^r\right\|^2\notag\\
&\leq C_0\int_0^t \Phi_{R_\theta}(t-r)
\Big(\E \|r_f^r(x^{[r]})\|^2+\|x^r\|^2\Big)\d r\label{eq:bound-Acal-moment}
\end{align}
for large enough $C_0>0$. From the definition of $\pRf$ in \eqref{eq:def-pRf}, we have
\begin{equation}\label{eq:def-pRf-lem}
\pRf(t-r)=C_0\left(\int_r^t \pRt(t-r')\pRf(r'-r)\d r'
+\Phi_{R_\theta}(t-r)\right).
\end{equation}
Integrating both sides of the equation with respect to ${\norm{x^r}^2}$, we get
\begin{align*}
\int_0^t \pRf(t-r) \norm{x^r}^2\d r
&=C_0\int_0^t \prn{\int_r^t \pRt(t-r') \pRf(r'-r) \d r' + \pRt(t-r)} \norm{x^r}^2\d r\nonumber\\
    &=C_0\int_0^t \pRt(t-r)\prn{\int_0^r \pRf(r-r')\norm{x^{r'}}^2\d {r'}+
\norm{x^r}^2}\d r.
\end{align*}
Comparing with \eqref{eq:bound-Acal-moment}, this shows that
\begin{equation}\label{eq:Erfsquaredbound}
\|R_f^t(x^{[t]})\|^2
=\|\E r_f^{t}(x^{[t]})\|^2
\leq \E\norm{r_f^{t}(x^{[t]})}^2 \leq \int_0^t \pRf(t-r) \norm{x^r}^2\d r.
\end{equation}
Further bounding
$(\int_0^t \pRf(t-r) \norm{x^r}^2\d r)^2
\leq t\pRf(t)^2\int_0^t \norm{x^r}^4\d r$, this verifies that
$R_f^t$ is a bounded linear operator from $L^4([0,t],\R^k)$ to $\R^k$. This
checks all conditions of Definition \ref{def:S} for $R_f$.

For $R_f^*$, we have similarly
\begin{align*}
    \E \norm{r_f^{t,*}}^2
&\leq 2\E\left\|\D_\xi f(\xi^t,w^*,\eps) \int_0^t
\frac{\bar\eta^r}{\bar\kappa}
\Rt^{t,r}r_f^{r-,*}\d z^r\right\|^2
+2\E\|\D_{w^*} f(\xi^t,w^*,\eps)\|^2\\
&\leq C_0\left(1+\int_0^t \pRt(t-r)\E\|r_f^{r,*}\|^2\d r\right)
\end{align*}
for large enough $C_0>0$. Comparing with the definition of $\Phi_{R_f^*}$ in
\eqref{eq:def-pRfs} shows
$\|R_f^{t,*}\|^2 \leq \E\|r_f^{t,*}\|^2 \leq \pRfs(t)$. Finally, 
the measurability of $t \mapsto \Gamma^t$ and bound
$\|\Gamma_t\|_\op \leq M_f$ follow directly from the definition of $\Gamma^t$ 
in \eqref{eq:def-cont-Gamma}, and \eqref{eq:defMbounds}.
Thus $(C_f,R_f,R_f^*,\Gamma) \in \cS_\xi^{\mathrm{cont}}(T,C_0)$.
\end{proof}

\subsection{Contractivity of the mappings}\label{sec:contractivity}

We now endow $\mathcal{S}_\xi \equiv \cS_\xi(T,C_0)$ and ${\mathcal{S}}_\theta
\equiv \cS_\theta(T,C_0)$ with metrics, under which we will show that
$\trsfrm \equiv \trsfrmA \circ \trsfrmB$ is contractive.

Fix a constant $\lambda>0$ (which we will also take large enough, depending
on $T,C_0$). For covariance kernels $C_{f,1},C_{f,2}$ on $[0,T] \otimes
\R^k$ and $C_{\theta,1},C_{\theta,2}$ on $([0,T] \otimes \R^k) \times \R^{k^*}$
satisfying the conditions of Definition \ref{def:S}, define
\begin{align*}
\lbddst{C_{f,1}}{C_{f,2}}=\inf_{(u_1,u_2) \sim \gamma \in
\Gamma(C_{f,1},C_{f,2})} \sup_{t \in [0,T]} e^{-\lambda t} \big(\E {\norm{u^t_1 -
u^t_2}^4}\big)^{1/4},\\
\lbddst{C_{\theta,1}}{C_{\theta,2}}=\inf_{(w_1,w_2) \sim \gamma \in
\Gamma(C_{\theta,1},C_{\theta,2})} \sup_{t \in [0, T]} e^{-\lambda t} \big(\E
\norm{w^t_1 - w^t_2}^4\big)^{1/4},
\end{align*}
where $\Gamma(C_{f,1},C_{f,2})$ denotes the set of all couplings between the
Gaussian process laws $\GP(0,C_{f,1})$ and $\GP(0,C_{f,2})$, and similarly for
$\Gamma(C_{\theta,1},C_{\theta,2})$. Note that under the constraints and
continuity conditions for $C_f,C_\theta$ in Definition \ref{def:S}, if 
$\lbddst{C_{\theta,1}}{C_{\theta,2}}=0$, then
$C_{\theta,1}^{t,s}=C_{\theta,2}^{t,s}$ pointwise for all $t,s \in [0,T]
\cup \{*\}$, and similarly for $C_f$.

For pairs of response kernels/operators satisfying Definition \ref{def:S}, we
define also
\begin{align*}
    \lbddst{R_{f,1}}{R_{f,2}} &= \sup_{0 \leq t \leq T} e^{-\lambda t}
\normop{R_{f,1}^t - R_{f,2}^t},\\
    \lbddst{R_{f,1}^*}{R_{f,2}^*} &= \sup_{0 \leq t \leq T} e^{-\lambda t}
\norm{R_{f,1}^{t,*} - R_{f,2}^{t,*}},\\
    \lbddst{\Gamma_1}{\Gamma_2} &= \sup_{0 \leq t \leq T} e^{-\lambda t}
\norm{\Gamma_1 - \Gamma_2},\\
    \lbddst{R_{\theta,1}} {R_{\theta,2}} &= \sup_{0 \leq s\leq t \leq T}
e^{-\lambda t} \norm{R_{\theta,1}^{t,s} - R_{\theta,2}^{t,s}},
\end{align*}
where we recall that $\normop{R_f^t}$ is the operator norm as a linear map
from $L^4([0,t],\R^k)$ to $\R^k$.
Finally, for any $X_i=(C_{f,i},R_{f,i},R_{f,i}^*,\Gamma_i) \in \mathcal{S}_\xi$
and $Y_i=(C_{\theta,i},R_{\theta,i}) \in {\mathcal{S}}_\theta$, we define 
\begin{subequations}
\begin{align*}
	\lbddst{X_1}{X_2} & = \lbddst{C_{f,1}}{C_{f,2}} +
\lbddst{R_{f,1}}{R_{f,2}} + \lbddst{R_{f,1}^*}{R_{f,2}^*}
+ \lbddst{\Gamma_1}{\Gamma_2},\\
	\lbddst{Y_1}{Y_2} & = \lbddst{C_{\theta,1}}{C_{\theta,2}}
+\lbddst{R_{\theta,1}}{R_{\theta,2}}.
\end{align*}
\end{subequations}

Throughout these arguments, we will fix $T,C_0>0$, and write
$\lesssim$ to mean inequality up to a constant depending on $T,C_0$.

We first study the contractivity properties for $\trsfrmA$.

\begin{lem}\label{lem:exist-unique-theta}
Fix any $T,C_0>0$, suppose $X=(C_f,R_f,R_f^*,\Gamma) \in
\cS_\xi(T,C_0)$, and let $\{\theta^t\}_{t \in [0,T]}$ and
$\{r_\theta^{t,s}\}_{t,s \in [0,T]}$ be defined
by \eqref{eq:def-cont-theta} and \eqref{eq:def-cont-deriv-t}.
Then there is a constant $C \equiv C(T,C_0)>0$ such that for all $t,s \in [0,T]$,
\[\E\|\theta^t\|^4 \leq C,
\qquad \E\|r_\theta^{t,s}\|^4 \leq C.\]
\end{lem}

\begin{proof}
By \eqref{eq:pRf-volterra},
\[\norm{R_f^{r}(\theta^{[r]})}^4 \leq \prn{\int_0^r \pRf(r-s)\norm{\theta^s}^2\d
s }^2\lesssim \int_0^r \norm{\theta^s}^4\d s.\]
Then from the definition \eqref{eq:def-cont-theta}, we have
\begin{align*}
    \norm{\theta^t}^4&\lesssim
\norm{\theta^0}^4+\int_0^t\Big(1+\norm{\theta^r}^4+\norm{R_f^{r}(\theta^{[r]})}^4+\norm{\theta^*}^4\Big)\d r+\norm{u^t}^4\\
    &\lesssim \norm{\theta^0}^4+\int_0^t
\Big(1+\norm{\theta^r}^4+\norm{\theta^*}^4\Big)\d r+\norm{u^t}^4.
\end{align*}
Since $\E\norm{u^t}^4 \lesssim \prn{\E\norm{u^t}^2}^2 \lesssim \pCf(T)^2 \lesssim 1$,
taking expectations and applying Gr\"onwall's inequality shows
the bound for $\E\|\theta^t\|^4$. A similar argument shows the bound
for $\E\|r_\theta^{t,s}\|^4$.
\end{proof}

\begin{lem}\label{lem:contract-Ct}
Fix any $T,C_0>0$. For $i=1,2$, 
suppose $X_i=(C_{f,i},R_{f,i},R_{f,i}^*,\Gamma_i) \in \mathcal{S}_\xi(T,C_0)$,
and let $(C_{\theta,i},R_{\theta,i})=\trsfrmA(X_i)$.
Then there exist constants $C \equiv C(T,C_0)>0$ and $\lambda_0 \equiv
\lambda_0(T,C_0)>0$ such that for any $\lambda>\lambda_0$,
\begin{align}
\lbddst{C_{\theta,1}}{C_{\theta,2}} & \leq C\,\lbddst
{C_{f,1}}{C_{f,2}}\notag\\
&\hspace{0.2in}+\frac{C}{\lambda}\prn{\lbddst{\Gamma_1} {\Gamma_2} +
\lbddst{R_{f,1}}{R_{f,2}} + \lbddst{R_{f,1}^*}{R_{f,2}^*}},
\label{eq:contractCtheta}\\
    \lbddst{R_{\theta,1}}{R_{\theta,2}} &\leq
\frac{C}{\lambda}\lbddst{X_1}{X_2}.\label{eq:contractRtheta}
\end{align}
\end{lem}
\begin{proof}
By definition of $\lbddst{\cdot}{\cdot}$,
there exists a coupling $\{u_1^t,u_2^t\}_{t \in [0,T]}$ of $\GP(0,C_{f,1})$ and
$\GP(0,C_{f,2})$ such that
\begin{align*}
\sup_{t \in [0, T]} e^{-4\lambda t} \big(\mathbb{E}
\norm{u_1^t-u_2^t}^4\big)^{1/4}
\leq 2\,\lbddst{C_{f,1}}{C_{f,2}}^4.
\end{align*}
For $i=1,2$, define $\theta_i^t$ via \eqref{eq:def-cont-theta} using
$X_i,\{u_i^t\}_{t \in [0,T]}$. Then
$\E \|\theta^t_1 - \theta^t_2\|^4 \lesssim
\int_0^t((\mathrm{I})_r+(\mathrm{II})_r)\d r +(\mathrm{III})$ where 
\begin{align*}
    (\mathrm{I})_r &=\E\left\|g(\theta_1^r)-g(\theta_2^r)+\gamma
\Gamma_1^r(\theta_1^r-\theta_2^r)+\gamma
R_{f,1}^r(\theta_1^{[r]}-\theta_2^{[r]})\right\|^4\\
    &\lesssim \E\norm{\theta_1^r-\theta_2^r}^4 + \left(\int_0^r \Phi_{\Rf}(r-r')
\E\norm{\theta^{r'}_1 - \theta^{r'}_2}^2 \d r'\right)^2\\
    &\lesssim e^{4\lambda r}e^{-4\lambda r}\E\norm{\theta_1^r-\theta_2^r}^4 +\int_0^r
e^{4\lambda r'}e^{-4\lambda r'}\E\norm{\theta^{r'}_1 - \theta^{r'}_2}^4 \d r'\\
    &\lesssim e^{4\lambda r} \sup_{t\in[0,T]}e^{-4\lambda t}
\E\norm{\theta_1^t-\theta_2^t}^4,\\
    (\mathrm{II})_r&=\E\left\|(\Gamma_1^r-\Gamma_2^r)\theta_2^r
+(R_{f,1}^r-R_{f,2}^r)\theta_2^{[r]}+
(R_{f,1}^{r,*}-R_{f,2}^{r,*})\theta^*\right\|^4\\
    &\lesssim
\norm{\Gamma_1^r-\Gamma_2^r}^4\E\|\theta_2^r\|^4+\normop{{R_{f,1}^r-R_{f,2}^r}}^4\E\|\theta_2^{[r]}\|_{L^4}^4 +
\norm{R_{f,1}^{r,*}-R_{f,2}^{r,*}}^4\E\|\theta^*\|^4\\
&\overset{(*)}\lesssim e^{4\lambda r}\Big(\lbddst{\Gamma_1}{\Gamma_2}^4
+\lbddst{R_{f,1}}{R_{f,2}}^4+\lbddst{R_{f,1}^*}{R_{f,2}^*}^4\Big),\\
(\mathrm{III})&=\E \norm{u_1^t-u_2^t}^4
\leq e^{4\lambda t} \cdot 2\lbddst{C_{f,1}}{C_{f,2}}^4.
\end{align*}
For $(*)$, we have used \Cref{lem:exist-unique-theta} to bound
$\E\|\theta^t_2\|^4,\E\|\theta_2^{[r]}\|_{L^4}^4 \lesssim 1$, and applied the
definitions of $\lbddst{\cdot}{\cdot}$.
Combining these bounds, multiplying by $e^{-4\lambda t}$, and taking the
supremum of over $t \in [0,T]$, we get
\begin{align*}
    &\sup_{t\in[0,T]}e^{-4\lambda t}\E\norm{\theta_1^t-\theta_2^t}^4\\
    &\lesssim \sup_{t\in[0,T]}e^{-4\lambda t}\E\norm{\theta_1^t-\theta_2^t}^4\cdot
\sup_{t\in[0,T]}e^{-4\lambda t} \int_0^t e^{4\lambda r}\d r\\
    &\hspace{0.5in}+
\Big(\lbddst{\Gamma_1}{\Gamma_2}^4
+\lbddst{R_{f,1}}{R_{f,2}}^4+\lbddst{R_{f,1}^*}{R_{f,2}^*}^4\Big)
\sup_{t\in[0,T]}e^{-4\lambda t}\int_0^t e^{4\lambda r} \d r\\
&\hspace{0.5in}+ \lbddst{C_{f,1}}{C_{f,2}}^4\\
&\lesssim \frac{1}{\lambda}\left(
\sup_{t\in[0,T]}e^{-4\lambda t}\E\norm{\theta_1^t-\theta_2^t}^4
+\lbddst{\Gamma_1}{\Gamma_2}^4
+\lbddst{R_{f,1}}{R_{f,2}}^4+\lbddst{R_{f,1}^*}{R_{f,2}^*}^4\right)\\
&\hspace{0.5in}+\lbddst{C_{f,1}}{C_{f,2}}^4.
\end{align*}
Thus for any $\lambda>\lambda_0(T,C_0)$ large enough, we may rearrange to get
\begin{align}
\sup_{t\in[0,T]}e^{-4\lambda t}\E\norm{\theta_1^t-\theta_2^t}^4 &\lesssim
\lbddst{C_{f,1}}{C_{f,2}}^4\label{eq:theta4diffbound}\\
&\hspace{0.2in}+\frac{1}{\lambda}\prn{\lbddst{\Gamma_1}{\Gamma_2}^4 +
\lbddst{R_{f,1}}{R_{f,2}}^4 + \lbddst{R_{f,1}^*}{R_{f,2}^*}^4}.\notag
\end{align}
Letting $\{w_1^t,w_2^t\}_{t \in [0,T] \cup \{*\}}$ be a jointly Gaussian process
with the same covariance as $\{\theta_1^t,\theta_2^t\}_{t \in [0,T] \cup
\{*\}}$, this gives a coupling of $C_{\theta,1},C_{\theta,2}$, so
\begin{align*}
\lbddst{C_{\theta,1}}{C_{\theta,2}}^4
&\leq \sup_{t \in [0, T]} e^{-4\lambda t} 
\E {\norm{w^t_1 - w^t_2}^4}
\lesssim \sup_{t \in [0, T]} e^{-4\lambda t} 
(\E {\norm{w^t_1 - w^t_2}^2})^2\\
&= \sup_{t \in [0, T]} e^{-4\lambda t} 
(\E{\norm{\theta^t_1 - \theta^t_2}^2})^2
\lesssim \sup_{t \in [0, T]} e^{-4\lambda t} 
\E{\norm{\theta^t_1 - \theta^t_2}^4}.
\end{align*}
Combining the above shows \eqref{eq:contractCtheta}.

For the analysis of $R_\theta$, we have similarly
from \eqref{eq:def-cont-deriv-t} that
$\E\|r_{\theta,1}^{t,s}-r_{\theta,2}^{t,s}\|^2
\lesssim \int_s^t ((\mathrm{I})_r+(\mathrm{II})_r)\d r$ where
\begin{align*}
(\mathrm{I})_r&=\E\Big\|(\gamma \Gamma_1^r+\D g(\theta_1^r))
(r_{\theta,1}^{r,s}-r_{\theta,2}^{r,s})
+\gamma R_f^r(r_{\theta,1}^{[r],s}-r_{\theta,2}^{[r],s})\Big\|^2\\
&\lesssim \E\|r_{\theta,1}^{r,s}-r_{\theta,2}^{r,s}\|^2
+\int_s^r \Phi_{R_f}(r-r')\E\|r_{\theta,1}^{r',s}-r_{\theta,2}^{r',s}\|^2\d
r'\\
&\lesssim e^{2\lambda r} \sup_{t \in [s,T]}
e^{-2\lambda t}\E\|r_{\theta,1}^{t,s}-r_{\theta,2}^{t,s}\|^2,\\
(\mathrm{II})_r&=\E\Big\|\gamma(\Gamma^t_1 - \Gamma^t_2)r_{\theta,2}^{r,s}
+(\D g(\theta^r_1)-\D g(\theta^r_2))r_{\theta,2}^{r,s}
+ \gamma (R_{f,1}^r-R_{f,2}^r)(r_{\theta,2}^{[r],s})\Big\|^2\\
&\lesssim \|\Gamma_1^r-\Gamma_2^r\|_\op^2 \E\|r_{\theta,2}^{r,s}\|^2
+\E[\|\D g(\theta_1)-\D g(\theta_2)\|_\op^2\|r_{\theta,2}^{r,s}\|^2]
+\normop{R_{f,1}^r-R_{f,2}^r}^2\E\|r_{\theta,2}^{[r],s}\|_{L^4}^2\\
&\lesssim \|\Gamma_1^r-\Gamma_2^r\|_\op^2 \E\|r_{\theta,2}^{r,s}\|^2
+(\E\|\theta_1^r-\theta_2^r\|^4)^{1/2}(\E\|r_{\theta,2}^{r,s}\|^4)^{1/2}
+\normop{R_{f,1}^r-R_{f,2}^r}^2\E\|r_{\theta,2}^{[r],s}\|_{L^4}^2\\
&\lesssim e^{2\lambda r}\lbddst{X_1}{X_2}^2.
\end{align*}
Here, $\|r_\theta^{[r],s}\|_{L^4}$ denotes
$(\int_0^r \|r_\theta^{r',s}\|^4 \d r')^{1/4}$,
and the last step using again \Cref{lem:exist-unique-theta} 
to bound $\E\|r_{\theta,2}^{[r],s}\|_{L^4}^4 \lesssim 1$, together with the
bound $\sup_{t \in [0,T]} e^{-2\lambda t}(\E\|\theta_1^t-\theta_2^t\|^4)^{1/2}
\lesssim \lbddst{X_1}{X_2}^2$ implied by \eqref{eq:theta4diffbound}.
Multiplying by $e^{-\lambda t}$ and taking the supremum over $0 \leq s \leq t
\leq T$,
\begin{align*}
&\sup_{0 \leq s \leq t \leq T} e^{-\lambda t}
\E\|r_{\theta,1}^{t,s}-r_{\theta,2}^{t,s}\|^2\\
&\lesssim \left(\sup_{0 \leq s \leq t \leq T} e^{-\lambda t}
\E\|r_{\theta,1}^{t,s}-r_{\theta,2}^{t,s}\|^2
+\lbddst{X_1}{X_2}^2\right)
\sup_{0 \leq s \leq t \leq T} e^{-\lambda t} \int_s^t e^{\lambda r}\d r\\
&\lesssim \frac{1}{\lambda}\left(\sup_{0 \leq s \leq t \leq T} e^{-\lambda t}
\E\|r_{\theta,1}^{t,s}-r_{\theta,2}^{t,s}\|^2
+\lbddst{X_1}{X_2}^2\right),
\end{align*}
and rearranging shows \eqref{eq:contractRtheta}.
\end{proof}

We next derive the contraction properties for the transformation $\trsfrmB$.

\begin{lem} \label{lem:contract-xi}
Fix any $T,C_0>0$.
Let $(\Omega,\cF,\{\cF\}_{t \in [0,T]},\P)$ be the filtered probability
space of Theorem \ref{thm:uniqueness-existence}, and
let $\{z^t\}_{t \in [0,T]}$ be either the Poisson process \eqref{eq:poisson-z}
or Gaussian process \eqref{eq:gaussian-z}.

For $i=1,2$, suppose $Y_i=(C_{\theta,i},R_{\theta,i}) \in
{\mathcal{S}}_\theta(T,C_0)$, and let $\{w_1^t,w_2^t\}_{t \in [0,T] \cup \{*\}}$
be any coupling of $w_1 \sim \GP(0,C_{\theta,1})$
and $w_2 \sim \GP(0,C_{\theta,2})$. Define $\xi_i^t$, 
$r_{f,i}^{t,*}$, and $r_{f,i}^t(x^{[t]})$,
via \eqref{eq:def-cont-xi}, \eqref{eq:def-cont-deriv-xi-star}, and
\eqref{eq:def-cont-resp-xi} using
$Y_i,\{w_i^t\}_{t \in [0,T]}$.
Then there exists a constant $C \equiv C(T,C_0)>0$ such that for all $t \in
[0,T]$ and any process $x \in L^4([0,T],\R^k)$,
\begin{align}
\E\left\|\xi^t_1 - \xi^t_2\right\|^4 &\leq C
\left(\sup_{r \in [0,t]} \|R_{\theta,1}^{t,r}-R_{\theta,2}^{t,r}\|^4
+\E\|w_1^t-w_2^t\|^4\right),\label{eq:contract-xi}\\
\E\|r_{f,1}^{t}(x^{[t]})-r_{f,2}^{t}(x^{[t]})\|^2 &\leq C
\left(\sup_{r \in [0,t]} \|R_{\theta,1}^{t,r}-R_{\theta,2}^{t,r}\|^2
+(\E\|w_1^t-w_2^t\|^4)^{1/2}
+\Big(\int_0^t \E\|w_1^r-w_2^r\|^4 \d r\Big)^{1/2}\right)\notag\\
&\hspace{1in}\times \left(\int_0^t \|x^r\|^4 \d r\right)^{1/2},\label{eq:contract-rfxi}\\
\E\norm{r_{f,1}^{t,*}-r_{f,2}^{t,*}}^2
&\leq C\left(\sup_{r \in [0,t]} \|R_{\theta,1}^{t,r}-R_{\theta,2}^{t,r}\|^2
+(\E\|w_1^t-w_2^t\|^4)^{1/2}\right).\label{eq:contract-rfxi-star}
\end{align}
\end{lem}
\begin{proof}
Applying the form of $\xi_i^t$, Lemma \ref{lem:burkholder}, and
the boundedness and Lipschitz properties of $f$,
\begin{align*}
\E\left\|\xi^t_1 - \xi^t_2\right\|^4 &\lesssim \E\left\|\int_0^t
\frac{\bar\eta^r}{\bar\kappa}
\prn{R_{\theta,1}^{t,r}f(\xi_1^{r-},w^*,\eps)-R_{\theta,2}^{t,r}f(\xi_2^{r-},w^*,\eps)}\d z^r\right\|^4 + \E\left\|w^t_1 - w^t_2\right\|^4\nonumber\\
  &\lesssim \int_0^t
\E\left\|R_{\theta,1}^{t,r}f(\xi_1^{r},w^*,\eps)-R_{\theta,2}^{t,r}f(\xi_2^r,w^*,\eps)\right\|^4 \d r+\E\left\|w^t_1 - w^t_2\right\|^4\\
&\lesssim \sup_{r \in [0,t]} \|R_{\theta,1}^{t,r}-R_{\theta,2}^{t,r}\|^4
+\int_0^t \E\|\xi_1^r-\xi_2^r\|^4\d r+\E\left\|w^t_1 - w^t_2\right\|^4.
\end{align*}
Then Gr\"onwall's inequality shows \eqref{eq:contract-xi}.

For \eqref{eq:contract-rfxi}, applying the form of $r_{f,i}^t(x^{[t]})$
and the Lipschitz property of $f$, we have
$\E\|r_{f,1}^{t}(x^{[t]}) - r_{f,2}^{t}(x^{[t]})\|^2 \lesssim \nI+\nII+\nIII$,
where
\begin{align*}
\nI&=\E\left\|\int_{0}^{t}
\frac{\bar\eta^r}{\bar\kappa}
  \prn{R_{\theta,1}^{t,r}(r_{f,1}^{r-}(x^{[r]})-r_{f,2}^{r-}(x^{[r]}))+\prn{R_{\theta,1}^{t,r}-R_{\theta,2}^{t,r}}r_{f,2}^{r-}(x^{[r]})}
\d z^r\right\|^2\nonumber\\
\nII&=\E\left\|\int_{0}^{t}
\frac{\bar\eta^r}{\bar\kappa} \prn{R_{\theta,1}^{t,r}(\D_\xi f(\xi^{r-}_1,w^*,\eps)-\D_\xi f(\xi^{r-}_2,w^*,\eps))
+ \prn{R_{\theta,1}^{t,r}-R_{\theta,2}^{t,r}}\D_\xi f(\xi^{r-}_2,w^*,\eps)}
x^{r-} \d z^r\right\|^2\nonumber\\
\nIII&=\E\left\|\prn{\D_\xi f(\xi^t_1,w^*,\eps) - \D_\xi f(\xi^t_2,w^*,\eps)}
\int_{0}^{t}
  \frac{\bar\eta^r}{\bar\kappa} R_{\theta,2}^{t,r}\bigl(r_{f,2}^{r-}(x^{[r]})+\D_\xi f(\xi^{r-}_2,w^*,\eps)x^{r-}\bigr) \d z^r\right\|^2.
\end{align*}
Using \eqref{eq:basicdzrbound} and \eqref{eq:Erfsquaredbound},
the first term can be bounded as
\begin{align*}
\nI &\lesssim
\int_0^t\prn{\pRt(t-r)^2\E\big\|r_{f,1}^{r}(x^{[r]})-r_{f,2}^{r}(x^{[r]})\big\|^2
+ \norm{R_{\theta,1}^{t,r}-R_{\theta,2}^{t,r}}^2\E\norm{r_{f,2}^{r}(x^{[r]})}^2 }\d r\nonumber\\
&\lesssim \int_0^t \E\big\|r_{f,1}^{r}(x^{[r]})-r_{f,2}^{r}(x^{[r]})\big\|^2\d r
+\int_0^t \norm{R_{\theta,1}^{t,r}-R_{\theta,2}^{t,r}}^2 \int_0^r
\pRf(r-r')^2\norm{x^{r'}}^2\d r'\d r\nonumber\\
&\lesssim \int_0^t \E\big\|r_{f,1}^{r}(x^{[r]})-r_{f,2}^{r}(x^{[r]})\big\|^2\d
r+\sup_{r\in[0,t]}\norm{R_{\theta,1}^{t,r}-R_{\theta,2}^{t,r}}^2\Big(\int_0^t
\norm{x^r}^4\d r\Big)^{1/2}.
\end{align*}
Similarly, the second term is bounded as
\begin{align*}
\nII &\lesssim \int_0^t\E\norm{\xi_1^{r}-\xi_2^{r}}^2 \cdot \norm{x^r}^2\d r +
\sup_{r\in[0,t]}\norm{R_{\theta,1}^{t,r}-R_{\theta,2}^{t,r}}^2\int_0^t
\norm{x^r}^2\d r\\
&\lesssim 
\Big(\int_0^t \E\norm{\xi_1^{r}-\xi_2^{r}}^4 \d r\Big)^{1/2}
\Big(\int_0^t \norm{x^r}^4\d r\Big)^{1/2} +
\sup_{r\in[0,t]}\norm{R_{\theta,1}^{t,r}-R_{\theta,2}^{t,r}}^2
\Big(\int_0^t \norm{x^r}^4\d r\Big)^{1/2}
\end{align*}
The third term is bounded as
\begin{align*}
\nIII &\lesssim \E\brk{\norm{\xi_1^{t} - \xi_2^{t}}^2 \cdot \left\|
\int_{0}^{t} \frac{\bar\eta^r}{\bar\kappa} R_{\theta,2}^{t,r}\bigl(r_{f,2}^{r-}(x^{[r]})+\D_\xi
f(\xi^{r-}_2,w^*,\eps)x^{r-}\bigr)
\d z^r\right\|^2}\nonumber\\
&\leq (\E \norm{\xi_1^{t} - \xi_2^{t}}^4)^{1/2}\left(\E\left\|
\int_{0}^{t} \frac{\bar\eta^r}{\bar\kappa} R_{\theta,2}^{t,r}\bigl(r_{f,2}^{r-}(x^{[r]})+\D_\xi
f(\xi^{r-}_2,w^*,\eps)x^{r-}\bigr)
\d z^r\right\|^4\right)^{1/2}.
\end{align*}
Applying Lemma \ref{lem:burkholder} to bound the second factor, we have
\[\nIII \lesssim
(\E \norm{\xi_1^{t} - \xi_2^{t}}^4)^{1/2}\cdot \left(\int_0^t
\prn{\E\norm{r_{f,2}^{r}(x^{[r]})}^4 + \norm{x^{r}}^4} \d
r\right)^{1/2}.\]
Applying Lemma \ref{lem:burkholder} again to the definition of
$r_f^t(\cdot)$ shows
\[\E\norm{r_f^{t}(x^{[t]})}^4 \lesssim \int_0^t
\left(\E\norm{r_f^{r}(x^{[r]})}^4+\norm{x^r}^4\right)\d r,\]
so Gr\"onwall's inequality implies
$\E\norm{r_f^{t}(x^{[t]})}^4 \lesssim \int_{0}^{t}\norm{x^r}^4\d r$.
Therefore
\begin{align*}
\nIII &\leq  (\E\norm{\xi_1^{t} - \xi_2^{t}}^4)^{1/2}\left(\int_0^t
\norm{x^{r}}^4 \d r\right)^{1/2}.
\end{align*}
Combining the above bounds, we have
\begin{align*}
&\E\big\|r_{f,1}^{t}(x^{[t]}) - r_{f,2}^{t}(x^{[t]})\big\|^2
\lesssim
\int_0^t \E\big\|r_{f,1}^{r}(x^{[r]}) - r_{f,2}^{r}(x^{[r]})\|^2\d r \\
&\hspace{0.2in}+\left(\sup_{r\in[0,t]}\norm{R_{\theta,1}^{t,r}-R_{\theta,2}^{t,r}}^2+(\E\norm{\xi_1^t-\xi_2^t}^4)^{1/2}
+\Big(\int_0^t \E\norm{\xi_1^r-\xi_2^r}^4\d r\Big)^{1/2}\right)
\left(\int_{0}^{t}\norm{x^{r}}^{4}\d r \right)^{1/2}.
\end{align*}
Applying again Gr\"onwall's inequality and \eqref{eq:contract-xi} to bound
$\E\norm{\xi_1^r-\xi_2^r}^4$, we obtain \eqref{eq:contract-rfxi}.

Finally, from the form of $r_{f,i}^{t,*}$, we have
$\E\norm{r_{f,1}^{t,*}-r_{f,2}^{t,*}}^2 \lesssim \nI+\nII+\nIII$,
where
\begin{align*}
\nI&=\E\left\|\int_0^t \frac{\bar\eta^r}{\bar\kappa}
\prn{R_{\theta,1}^{t,r}r_{f,1}^{r-,*}-R_{\theta,2}^{t,r}r_{f,2}^{r-,*}}\d z^r\right\|^2, \\
\nII&= \E\left\|\Big(\D_\xi f(\xi_1^t,w^*,\eps)-\D_\xi f(\xi_2^t,w^*,\eps)\Big)
\int_0^t \frac{\bar\eta^r}{\bar\kappa}
R_{\theta,1}^{t,r}r_{f,1}^{r-,*}\d z^r\right\|^2,\\
\nIII&=\E\norm{\D_{w^*} f(\xi_1^t,w^*,\eps)-\D_{w^*} f(\xi_2^t,w^*,\eps)}^2.
\end{align*}
By similar arguments as above, we may bound each term to obtain
\[\E\norm{r_{f,1}^{t,*}-r_{f,2}^{t,*}}^2
\lesssim \int_0^t \E\norm{r_{f,1}^{r,*}-r_{f,2}^{r,*}}^2 \d r
+\sup_{r \in [0,t]}\|R_{\theta,1}^{t,r}-R_{\theta,2}^{t,r}\|^2
+(\E\|\xi_1^t-\xi_2^t\|^4)^{1/2}.\]
Applying Gr\"onwall's inequality and \eqref{eq:contract-xi} to bound
$\E\norm{\xi_1^t-\xi_2^t}^4$ shows \eqref{eq:contract-rfxi-star}.
\end{proof}

\begin{lem}\label{lem:contract-fsys}
Fix any $T,C_0>0$. For $i=1,2$, suppose
$Y_i=(C_{\theta,i},R_{\theta,i}) \in {\mathcal{S}}_\theta(T,C_0)$, and
let $(C_{f,i},R_{f,i},R_{f,i}^*,\Gamma_i)=\trsfrmB(Y_i)$. Then there exist
constants $C \equiv C(T,C_0)>0$ and $\lambda_0 \equiv \lambda_0(T,C_0)>0$
such that for any $\lambda>\lambda_0$,
\[\begin{gathered}
\lbddst{C_{f,1}}{C_{f,2}} \leq \frac{C}{\sqrt{\lambda}}\,\lbddst{Y_1}{Y_2}, \\
\lbddst{R_{f,1}}{R_{f,2}},\lbddst{R_{f,1}^*}{R_{f,2}^*},
\lbddst{\Gamma_1}{\Gamma_2} \leq C\,\lbddst{Y_1}{Y_2}.
\end{gathered}\]
\end{lem}
\begin{proof}
Let $\{w^t_1,w^t_2\}_{t \in [0,T] \cup \{*\}}$ be a coupling
of $\GP(0,C_{\theta,1})$ and $\GP(0,C_{\theta,2})$ that satisfies
\begin{equation}\label{eq:wcoupling}
\sup_{t \in [0, T]} e^{-\lambda t} (\mathbb{E}\norm{w^t_1 -
w^t_2}^4)^{1/4}  \leq 2\,\lbddst{C_{\theta,1}}{C_{\theta,2}}.
\end{equation}
Then by \eqref{eq:contract-xi} and the definition of $\lbddst{\cdot}{\cdot}$,
\[e^{-4\lambda t}
\E\|\xi_1^t-\xi_2^t\|^4 \lesssim 
\lbddst{R_{\theta,1}}{R_{\theta,2}}^4+\lbddst{C_{\theta,1}}{C_{\theta,2}}^4
\lesssim \lbddst{Y_1}{Y_2}^4.\]
For $i=1,2$, define
\[h_i^t=\int_0^t \frac{\bar\eta^r}{\bar\kappa}f(\xi_i^{r-},w^*,\eps) \d z^r.\]
Then by the Lipschitz property of $f$ and \eqref{eq:basicdzrbound},
\[\E\|h_1^t-h_2^t\|^2
\lesssim \int_0^t \E\|\xi_1^r-\xi_2^r\|^2 \d r
\leq \sup_{r \in [0,t]}
e^{-2\lambda r}\,\E\|\xi_1^r-\xi_2^r\|^2
\cdot \int_0^t e^{2\lambda r} \d r
\lesssim \lbddst{Y_1}{Y_2}^2 \cdot \int_0^t e^{2\lambda r} \d r,\]
so
\[\sup_{t \in [0,T]} e^{-2\lambda t}
\E\|h_1^t-h_2^t\|^2 \lesssim \frac{1}{\lambda}\,\lbddst{Y_1}{Y_2}^2.\]
Letting $\{u_1^t,u_2^t\}_{t \in [0,T]}$ be a jointly Gaussian process with the
same covariance as $\{h_1^t,h_2^t\}_{t \in [0,T]}$, this gives a coupling of
$\GP(0,C_{f,1})$ and $\GP(0,C_{f,2})$, so
\begin{align*}
\lbddst{C_{f,1}}{C_{f,2}}^4
&\leq \sup_{t \in [0, T]} e^{-4\lambda t} 
\E {\norm{u^t_1 - u^t_2}^4}
\lesssim \sup_{t \in [0, T]} e^{-4\lambda t} 
(\E {\norm{u^t_1 - u^t_2}^2})^2\\
&=\sup_{t \in [0, T]} e^{-4\lambda t} 
(\E{\norm{h^t_1 - h^t_2}^2})^2
\lesssim \frac{1}{\lambda^2}\,\lbddst{Y_1}{Y_2}^4.
\end{align*}

For $R_f$, note that by \eqref{eq:contract-rfxi},
\begin{align*}
&\normop{R_{f,1}^t-R_{f,2}^t}^2
=\sup_{x:\|x\|_{L^4} \leq 1}
\|\E r_{f,1}^t(x^{[t]})-\E r_{f,2}^t(x^{[t]})\|^2\\
&\lesssim \sup_{r \in [0,t]}
\|R_{\theta,1}^{t,r}-R_{\theta,2}^{t,r}\|^2
+(\E\|w_1^t-w_2^t\|^4)^{1/2}
+\Big(\int_0^t \E\|w_1^r-w_2^r\|^4 \d r\Big)^{1/2}\\
&\lesssim \sup_{r \in [0,t]}
\|R_{\theta,1}^{t,r}-R_{\theta,2}^{t,r}\|^2
+\sup_{r \in [0,t]} 
(e^{-4\lambda r}\E \|w_1^r-w_2^r\|^4)^{1/2}
\left(e^{2\lambda t}+\left(\int_0^t e^{4\lambda r}\d r\right)^{1/2}\right).
\end{align*}
Multiplying by $e^{-2\lambda t}$ and applying \eqref{eq:wcoupling} shows
\[\lbddst{R_{f,1}}{R_{f,2}}^2
\lesssim
\lbddst{R_{\theta,1}}{R_{\theta,2}}^2+\lbddst{C_{\theta,1}}{C_{\theta,2}}^2
\lesssim \lbddst{Y_1}{Y_2}^2.\]
A similar argument using \eqref{eq:contract-rfxi-star} 
and $\|\Gamma_1^t-\Gamma_2^t\| \lesssim \E\|\xi_1^t-\xi_2^t\|$ shows
$\lbddst{R_{f,1}^*}{R_{f,2}^*} \lesssim \lbddst{Y_1}{Y_2}$
and $\lbddst{\Gamma_1}{\Gamma_2} \lesssim \lbddst{Y_1}{Y_2}$.
\end{proof}

We now conclude the proof of \Cref{thm:uniqueness-existence}.

\begin{proof}[Proof of \Cref{thm:uniqueness-existence}(b)]
Fix $T>0$, and any constant $C_0>0$ large enough.
By Lemmas \ref{lem:solution-in-space} and \ref{lem:solution-in-space-2},
\[\trsfrm:=\trsfrmA \circ \trsfrmB\]
defines a mapping from $\cS_\theta^\text{cont} \equiv
\cS_\theta^\text{cont}(T,C_0)$ to itself. For $i=1,2$, consider any
$Y_i=(C_{\theta,i},R_{\theta,i}) \in {\mathcal{S}}_\theta^\text{cont}$,
let $X_i=(C_{f,i},R_{f,i},R_{f,i}^*,\Gamma_i)=\trsfrmB(Y_i)$.
Then by Lemmas \ref{lem:contract-Ct} and \ref{lem:contract-fsys},
for any large enough $\lambda>0$,
\begin{align}\label{eq:T-contraction}
    &\lbddst{\trsfrm(Y_1)}{\trsfrm(Y_2)}
\lesssim 
\lbddst{C_{f,1}}{C_{f,2}}+\frac{1}{\lambda}\lbddst{X_1}{X_2}
\lesssim \frac{1}{\sqrt{\lambda}}\lbddst{Y_1}{Y_2}.
\end{align}
Thus for large enough $\lambda>0$, we have
$\lbddst{\trsfrm(Y_1)}{\trsfrm(Y_2)} \leq \frac{1}{2}\,\lbddst{Y_1}{Y_2}$,
implying that $\trsfrm:\cS_\theta^\text{cont} \to \cS_\theta^\text{cont}$
is contractive in the metric $\lbddst{\cdot}{\cdot}$.

It is clear that $\lbddst{R_{\theta,1}}{R_{\theta,2}}$ defines a complete metric
on the space of Borel-measurable $\R^{k \times k}$-valued
processes $\{R_\theta^{t,s}\}_{t,s \in [0,T]}$, and the conditions in
Definition \ref{def:S} for
$\{R_\theta^{t,s}\}_{t,s \in [0,T]}$ to belong to $\cS_\theta^{\text{cont}}$ are
closed under $\lbddst{R_{\theta,1}}{R_{\theta,2}}$. For the metric
$\lbddst{C_{\theta,1}}{C_{\theta,2}}$, define also
\[\|C_{\theta,1}-C_{\theta,2}\|_\infty
=\sup_{t,s \in [0,T] \cup \{*\}} |C_{\theta,1}^{t,s}-C_{\theta,2}^{t,s}|.\]
We have by Cauchy-Schwarz that
for any $C_{\theta,1},C_{\theta,2}$ satisfying the conditions of
$\cS_\theta^{\text{cont}}$,
\begin{align*}
\|C_{\theta,1}-C_{\theta,2}\|_\infty &\leq
\inf_{w_1,w_2 \sim \gamma \in \Gamma(C_{\theta,1},C_{\theta,2})}
\sqrt{\E \brk{\norm{w_1^t}^2} \cdot \E \brk{\norm{w_1^s - w_2^s}^2}} + \sqrt{\E
\brk{\norm{w_2^s}^2} \cdot \E \brk{\norm{w_1^t - w_2^t}^2}} \nonumber \\
	& \leq 2\sqrt{\pCt(T)} e^{\lambda T} \cdot
\lbddst{R_{\theta,1}}{R_{\theta,2}}.
\end{align*}
Conversely, by the condition \eqref{eq:Cthetacontinuity} and
\cite[Lemma D.1]{fan2025dynamical2}, there exists a jointly
Gaussian coupling $(w_1,w_2) \sim \gamma \in \Gamma(C_{\theta,1},C_{\theta,2})$
for which
\begin{align*}
\lbddst{C_{\theta,1}}{C_{\theta,2}}
&\leq \sup_{t \in [0,T]} e^{-\lambda t}(\E\|w_1^t-w_2^t\|^4)^{1/4}
\leq C\sup_{t \in [0,T]} e^{-\lambda t}(\E\|w_1^t-w_2^t\|^2)^{1/2}\\
&\leq C'\Big(\sqrt{\|C_{\theta,1}-C_{\theta,2}\|_\infty}
+\|C_{\theta,1}-C_{\theta,2}\|_\infty\Big)^{1/2},
\end{align*}
where $C,C'>0$ are some constants depending on $T,C_0,k$. Thus, the metrics
$\lbddst{C_{\theta,1}}{C_{\theta,2}}$ and
$\|C_{\theta,1}-C_{\theta,2}\|_\infty$ are uniformly equivalent over the space
of covariance kernels $\{C_\theta^{t,s}\}_{t,s \in [0,T] \cup \{*\}}$
satisfying the conditions of $\cS_\theta^{\text{cont}}$. Since this space is
complete under $\|C_{\theta,1}-C_{\theta,2}\|_\infty$, it is also complete under
$\lbddst{C_{\theta,1}}{C_{\theta,2}}$.

Then $\lbddst{\cdot}{\cdot}$ defines a complete metric on
$\cS_\theta^{\text{cont}}$, so by the Banach fixed-point theorem, there exists a
unique fixed point $Y \in \mathcal{S}^{\mathrm{cont}}_\theta$ such that
$\mathcal{T}(Y)=Y$. Since the image of $\mathcal{T}:\cS_\theta \to \cS_\theta$
is contained in $\cS_\theta^\text{cont}$, this fixed point is also unique in
$\cS_\theta$. Finally, this implies that for any
$C_0>0$ large enough, there is a unique fixed point
$(C_\theta,R_\theta,C_f,R_f,R_f^*,\Gamma)
\in \cS \equiv \cS(T,C_0)$ satisfying
(\ref{eq:def-cont-C-t}--\ref{eq:def-cont-Gamma}).
\end{proof}


\section{DMFT approximation in discrete time}

In this section, we introduce and analyze
a time-discretized version of the SGD and SME
processes: Fixing a discretization size $\delta>0$, let $\{\z_\delta^t\}_{t
\geq 0}$ be i.i.d.\ random vectors in $\R^n$ encapsulating the data sampling
noise, with law
\[\begin{cases}
\z_\delta^t=\bPi \bP_\delta^t & \text{for SGD},\\
\z_\delta^t \sim
\cN(\delta \bar\kappa1_n,\delta\bar\kappa\,\Id_{n\times n}) & \text{for SME}.
\end{cases}\]
In the SGD setting, $\bP_\delta^t \in \R^{{n\choose \kappa}}$ is a random vector
with i.i.d.\ $\Pois(\delta n^{1-\alpha} / {n\choose\kappa})$ entries,
and $\bPi\in\R^{n\times {n\choose \kappa}}$ is an incidence matrix having
entries
\[\Pi_{i,S}=\1_{i \in S} \text{ for all $i \in [n]$ and $S \subset \cS$}.\]
Note that marginally, each entry of $\bPi\bP_\delta^t$ has law
$\Pois(\binom{n-1}{\kappa-1} \cdot \delta n^{1-\alpha} /
{n\choose\kappa})=\Pois(\delta \kappa/n^\alpha)$,
where $\delta \kappa/n^\alpha \approx \delta\bar\kappa$
for large $n$ under Assumption \ref{asp:SGD}.

For both SGD and SME, we consider the discrete-time dynamics
\begin{align}
{\btheta}_\delta^{t+1}&={\btheta}_\delta^{t}
-\frac{\bar\eta^{t\delta}}{\bar\kappa}
\sum_{i=1}^n \x_i \otimes \Big(f(\x_i^\top \btheta_\delta^t,\x_i^\top
\btheta^*,\eps_i) \cdot z_{\delta,i}^t\Big)
-\delta\bar\eta^{t\delta} g(\btheta_\delta^t)\notag\\
&={\btheta}_\delta^{t}
-\frac{\bar\eta^{t\delta}}{\bar\kappa} \X^\top \big(f(\X \btheta_\delta^t,
\X\btheta^*,\beps) \odot \z_\delta^t\big)
- \delta\bar\eta^{t\delta} g({\btheta}_\delta^{t}).\label{eq:discretedynamics}
\end{align}
Here and throughout, $\X=[\x_1,\ldots,\x_n]^\top \in \R^{n \times d}$,
$f:\R^{n \times (k+k^*+1)} \to \R^{n \times k}$
and $g:\R^{d \times k} \to \R^{d \times k}$ denote the applications of
$f(\cdot)$ and $g(\cdot)$ row-wise, and $\odot$ denotes the row-wise
scalar product. We will write also
\[{\bxi}_\delta^{t}=\X {\btheta}_\delta^{t},
\qquad \bxi^*=\X \btheta^*.\]
We remark that these discrete-time dynamics \eqref{eq:discretedynamics}
may be understood as a variant of SGD/SME where the stochastic
gradient or drift/diffusion coefficients are only updated in (rescaled) time
increments of $\delta$, i.e.\ they are computed in \eqref{eq:SGDrescaled}
and \eqref{eq:SME} using $\btheta^{k\delta}$ instead of $\btheta^t$ if $t
\in [k\delta,(k+1)\delta)$.

We show in Section \ref{sec:discreteDMFT} that
these discrete-time dynamics are characterized by a system of discrete-time
DMFT equations as $n,d \to \infty$, and in Section
\ref{sec:deltaDMFTconvergence} that this discrete-time DMFT system converges
to the continuous-time system defined in Section \ref{sec:DMFTlimits} as
$\delta \to 0$. This proof strategy was introduced in \cite{celentano2021high}, and our
arguments will follow closely its implementation in \cite{fan2025dynamical1}.

\subsection{Discrete DMFT approximation via AMP}\label{sec:discreteDMFT}

Given the same variables $\theta^0,\theta^*,\eps$ as in
\eqref{eq:thetaeps}, the discrete DMFT system is defined by the primary
processes
\begin{align}
\theta^t_\delta&=
\theta^0-\sum_{r=0}^{t-1} \delta\bar\eta^{r\delta}\prn{{\gamma}\Gamma_\delta^r \theta^r_\delta+g(\theta^r_\delta)+\gamma\sum_{q=0}^{r-1}R_{f,\delta}^{r,q}
\theta^q_\delta
+ {\gamma} R_{f,\delta}^{r,*}\theta^*}+\sqrt{\gamma}u_\delta^t,
\label{eq:def-dis-theta}\\
\xi^t_\delta&={-}\sum_{r=0}^{t-1}\frac{\bar\eta^{r\delta}}{\bar\kappa}
R_{\theta,\delta}^{t,r}f(\xi_\delta^r,w^*,\eps)z^r_\delta+w^t_\delta,\label{eq:def-dis-xi}
\end{align}
and the response processes
\begin{align}
r_{\theta,\delta}^{t,s}&=\Id_k
-\sum_{r=s}^{t-1} \delta\bar\eta^{r\delta}\bigg[\left({\gamma}\Gamma_\delta^r+
\D g(\theta_\delta^r)\right)\cdot r_{\theta,\delta}^{r,s}
+{\gamma} \sum_{q=s+1}^{r-1}  R_{f,\delta}^{r,q}
r_{\theta, \delta}^{q,s}\bigg] \text{ for } s<t,\notag\\
&\hspace{1in} r_{\theta,\delta}^{t,s}=0 \text{ for } s \geq t,\label{eq:def-dis-deriv-t}\\
r_{f,\delta}^{t,*}&={-}\D_\xi f(\xi^t_\delta,w^*,\eps) \sum_{r=0}^{t-1}
\frac{\bar\eta^{r\delta}}{\bar\kappa}
R_{\theta,\delta}^{t,r}r_{f,\delta}^{r,*}z^r_\delta
+\D_{w^*} f(\xi^t_\delta,w^*,\eps),\label{eq:def-dis-deriv-xi-star}\\
    r_{f,\delta}^{t,s}&=-\D_\xi
f(\xi_\delta^{t},w^*,\eps)\prn{\sum_{r=s+1}^{t-1}\frac{\bar\eta^{r\delta}}{\bar\kappa}
R_{\theta,\delta}^{t,r} r_{f,\delta}^{r,s}z^r_\delta +
\frac{\bar\eta^{s\delta}}{\bar\kappa} R_{\theta,\delta}^{t,s}\cdot\D_\xi
f(\xi_\delta^{s},w^*,\eps)z^s_\delta} \text{ for } s<t,\notag\\
&\hspace{1in}r_{f,\delta}^{t,s}=0 \text{ for } s \geq t,\label{eq:def-dis-deriv-xi}
\end{align}
where
\[\{u_\delta^t\}_{t \geq 0} \sim \GP(0,C_{\theta,\delta}),
\qquad (\{w_\delta^t\}_{t \geq 0},w^*) \sim \GP(0,C_{f,\delta})\]
are discrete-time Gaussian processes,
\[\{z_\delta^t\}_{t \geq 0} \overset{iid}{\sim} \begin{cases}
\Pois(\delta\bar\kappa) & \text{ for SGD}, \\
\cN(\delta\bar\kappa,\delta\bar\kappa) & \text{ for SME},
\end{cases}\]
and these are independent of each other and of $\theta^0,\theta^*,\eps$.
The deterministic correlation and response kernels are defined as
expectations with respect to the above laws,
\begin{align}
C_{\theta,\delta}^{t,s}&=\E[\theta^t_\delta \otimes \theta^s_\delta] \text{ for } t,s \in [0,\infty)
\cup \{*\}\label{eq:def-dis-C-t}\\
C_{f,\delta}^{t,s}&=\E\brk{\sum_{r=0}^t \frac{\bar\eta^{r\delta}}{\bar\kappa}
f(\xi_\delta^r,w^*,\eps)z^r_\delta
\otimes \sum_{r=0}^s \frac{\bar\eta^{s\delta}}{\bar\kappa}
f(\xi_\delta^s,w^*,\eps)z^s_\delta}\label{eq:def-dis-C-h}\\
 R_{\theta,\delta}^{t,s}&=\E[ r_{\theta,\delta}^{t,s}]
\label{eq:def-dis-R-t}\\
R_{f,\delta}^{t,s}&=\E[r_{f,\delta}^{t,s}]\label{eq:def-dis-R-h}\\
R_{f,\delta}^{t,*}&=\E[ r_{f,\delta}^{t,*}]
\label{eq:def-dis-R-h-star}\\
\Gamma^t_\delta&=\E[\D_\xi f(\xi^t_\delta,w^*,\eps)]\label{eq:def-dis-Gamma}
\end{align}
We clarify that in contrast to the continuous-time setting, it is immediate to
see that these processes/kernels are uniquely defined recursively in time via
\[\begin{gathered}
\{\theta_\delta^t,r_{\theta,\delta}^{t,s}\}_{s,t=0}
\mapsto \{C_{\theta,\delta}^{t,s},R_{\theta,\delta}^{t,s}\}_{s,t=0}
\mapsto \{\xi_\delta^t,r_{f,\delta}^{t,s},r_{f,\delta}^{t,*}\}_{s,t=0}
\mapsto
\{C_{f,\delta}^{t,s},R_{f,\delta}^{t,s},R_{f,\delta}^{t,*},\Gamma_\delta^t\}_{s,t=0}
\mapsto\\
\{\theta_\delta^t,r_{\theta,\delta}^{t,s}\}_{s,t \leq 1}
\mapsto \{C_{\theta,\delta}^{t,s},R_{\theta,\delta}^{t,s}\}_{s,t \leq 1}
\mapsto \{\xi_\delta^t,r_{f,\delta}^{t,s},r_{f,\delta}^{t,*}\}_{s,t \leq 1}
\mapsto
\{C_{f,\delta}^{t,s},R_{f,\delta}^{t,s},R_{f,\delta}^{t,*},\Gamma_\delta^t\}_{s,t \leq 1} \mapsto \ldots
\end{gathered}\]

The main result of this section is the following lemma, which shows that
this discrete-time DMFT system approximates the dynamics
\eqref{eq:discretedynamics} in the large $n,d$ limit.

\begin{lem}\label{lem:finite_dim_converge}
Let $\{\theta_{\delta,j}^t,\theta_j^*\}$ be the rows of
$\btheta_\delta^t,\btheta^*$, and let $\{\xi_{\delta,i}^t,\xi_i^*,\eps_i\}$ be 
those of $\bxi_\delta^t,\bxi^*,\beps$.
Then for any fixed integer $T \geq 0$, almost surely as $n,d\rightarrow\infty$,
\begin{align*}
\frac{1}{d}\sum_{j=1}^d \delta_{(\theta^0_{\delta,j},\ldots,
\theta^T_{\delta,j},\theta^\ast_j)}
\Rightarrow \Law\big(\theta^0_\delta,\ldots, \theta^T_\delta,\theta^\ast\big),
\quad \frac{1}{n}\sum_{i=1}^n \delta_{(\xi^0_{\delta,i}, \ldots,
\xi^T_{\delta,i},\xi_i^*,\eps_i)}
\Rightarrow \Law\big(\xi^0_\delta,\ldots,\xi^T_\delta,w^*,\eps\big)
\end{align*}
weakly and in Wasserstein-2, where the right side denotes the joint law of these
variables defined via the discrete-time DMFT recursions
(\ref{eq:def-dis-theta}--\ref{eq:def-dis-Gamma}).
\end{lem}

We proceed to prove Lemma \ref{lem:finite_dim_converge},
by re-expressing the dynamics \eqref{eq:discretedynamics} as an Approximate
Message Passing (AMP) procedure
for which rigorous state evolution results are available. 

\begin{lem}\label{lem:initialconditions}
Let $\{\theta_j^0,\theta_j^*\}$ and
$\{\eps_i,z_{\delta,i}^t\}$ denote the rows of
$\btheta^0,\btheta^*$ and $\beps,\z_\delta^t$.
Then for any fixed integer $T \geq 0$, almost surely as $n,d \to \infty$,
\[\frac{1}{d}\sum_{j=1}^d
\delta_{(\theta_j^0,\theta_j^*)} \Rightarrow \Law(\theta^0,\theta^*)
\qquad \frac{1}{n}\sum_{i=1}^n
\delta_{(\eps_i,z_{\delta,i}^0,\ldots,z_{\delta,i}^T)}
\Rightarrow \Law(\eps,z_\delta^0,\ldots,z_\delta^T)\]
weakly and in Wasserstein-$p$ for each fixed order $p \geq 1$.
\end{lem}
\begin{proof}
The convergence for $(\btheta^0,\btheta^*)$ is assumed in
Assumption \ref{asp:model}(c).
In the SME setting, the convergence for
$(\beps,\z_\delta^0,\ldots,\z_\delta^T)$ is also immediate from Assumption
\ref{asp:model}(c) and the independence of entries of $\z_\delta^t$,
see e.g.\ \cite[Proposition E.1]{fan2022approximate}.

In the setting of SGD, fix any $k \geq 0$ and
monic monomial $q:\R^{T+1} \to \R$, and consider 
\[Q=\frac{1}{n}\sum_{i=1}^n \eps_i^k q(z_{\delta,i}^0,\ldots,z_{\delta,i}^T)\]
as a polynomial function of the i.i.d.\ Poisson variables
$\bP_\delta^0,\ldots,\bP_0^T$ defining
$\z_\delta^t=\bPi\bP_\delta^t$. Note that since $\bPi$ has nonnegative entries,
each term $q(z_{\delta,i}^0,\ldots,z_{\delta,i}^T)$ is then some polynomial of
the coordinates of $\bP_\delta^0,\ldots,\bP_\delta^T$ with
nonnegative coefficients. By the concentration inequality of
\cite[Theorem 1.4, Lemma 7.6]{schudy2012concentration}, for any $\lambda \geq 0$, we have
\begin{equation}\label{eq:poissonconcentration}
\P[|Q-\E Q| \geq \lambda]
\leq C\max\left\{\max_{r=1,\ldots,\deg(q)}
e^{-\frac{c\lambda^2}{\mu_0\mu_r}},
\max_{r=1,\ldots,\deg(q)}
e^{-\left(\frac{c\lambda}{\mu_r}\right)^{1/r}}\right\}
\end{equation}
for some $(k,q)$-dependent constants $C,c>0$. Here $\deg(q)$ is the degree of
$q$, and (c.f.\ \cite[Section 1.4]{schudy2012concentration})
\[\mu_0 \leq \frac{1}{n}\sum_{i=1}^n
|\eps_i|^k \E q(z_{\delta,i}^0,\ldots,z_{\delta,i}^T),
\qquad \mu_r \leq \max_{|S|=r} \E\left[\partial_S 
\frac{1}{n}\sum_{i=1}^n 
|\eps_i|^k q(z_{\delta,i}^0,\ldots,z_{\delta,i}^T)\right]\]
where this bound for $\mu_r$ denotes the maximum expected mixed partial
derivative of order $r$ in any combination of the coordinates of
$\bP_\delta^0,\ldots,\bP_\delta^T$. Since
each coordinate of $\bP_\delta^0,\ldots,\bP_\delta^T$ can influence at most
$\kappa$ samples $i=1,\ldots,n$, it is clear by Assumption \ref{asp:model}(c) that 
for all large $n,d$, each $r \geq 1$, and some $(k,q,r)$-dependent constants
$C,C'>0$, we have
\[\mu_0 \leq C,
\quad \mu_r \leq \max_{S \subset [n]:|S|=\kappa} \frac{C}{n}\sum_{i \in S}
|\eps_i|^k
\leq \max_{S \subset [n]:|S|=\kappa}
C\left(\frac{1}{n}\sum_{i=1}^n |\eps_i|^{2k}\right)^{1/2}
\left(\frac{1}{n}\sum_{i=1}^n \1_{i \in S}\right)^{1/2}
\leq C'\sqrt{\frac{\kappa}{n}}.\]
Since $\kappa \asymp n^\alpha$
and $\alpha<1$, applying these bounds to \eqref{eq:poissonconcentration} shows,
by the Borel-Cantelli lemma,
that there exists a constant $\iota>0$ for which $|Q-\E Q|<n^{-\iota}$
a.s.\ for all large $n,d$. On the other hand, letting
$\{\tilde z_\delta^t\}_{t \geq 0} \overset{iid}{\sim}
\Pois(\delta\kappa/n^\alpha)$ and 
$\{z_\delta^t\}_{t \geq 0} \overset{iid}{\sim}
\Pois(\delta\bar \kappa)$, by Assumptions \ref{asp:SGD} and \ref{asp:model}
we have
\[\E Q=\E q(\tilde z_\delta^0,\ldots,\tilde z_\delta^T)
\cdot \frac{1}{n}\sum_{i=1}^n \eps_i^k
=\E q(z_\delta^0,\ldots,z_\delta^T) \cdot \E \eps^k+o(1).\]
This shows that $Q$ converges a.s.\ to $\E q(z_\delta^0,\ldots,z_\delta^T) \cdot
\E \eps^k$, i.e.\ each mixed moment of the empirical measure
$\frac{1}{n}\sum_{i=1}^n
\delta_{(\eps_i,z_{\delta,i}^0,\ldots,z_{\delta,i}^T)}$
converges a.s.\ to that of $\Law(\eps,z_\delta^0,\ldots,z_\delta^T)$.
Since $\eps$ has finite moment generating function around 0, so
does $(\eps,z_\delta^0,\ldots,z_\delta^T)$. Then this convergence of mixed
moments implies convergence weakly and in Wasserstein-$p$
\cite[Thm 6.9]{villani2008optimal}, establishing the lemma.
\end{proof}

\begin{proof}[Proof of Lemma \ref{lem:finite_dim_converge}]
The proof relies on mapping the dynamics \eqref{eq:discretedynamics} to an
AMP algorithm: Let us set
\[\V^0=[\V_1^0,\V_2^0]:=[\btheta^0,\btheta^*] \in \R^{d \times (k+k^*)},
\quad \V^t=[\V_1^t,\V_2^t]:=[\btheta_\delta^t-\btheta_\delta^{t-1},0]
\in \R^{d \times (k+k^*)} \text{ for } t \geq 1,\]
\[\Y^t=[\Y_1^t,\Y_2^t]:=\left[{-}\frac{\bar\eta^{t\delta}}{\bar\kappa}f(\X\btheta_\delta^t,\X\btheta^*,\beps)
\odot \z_\delta^t,\,0\right] \in \R^{n
\times (k+k^*)} \text{ for } t \geq 0.\]
Then
\begin{equation}\label{eq:thetaV}
[\btheta_\delta^t,\btheta^*]=\sum_{r=0}^t \V^r.
\end{equation}
For any matrices $A^{t,q},B^{t,q} \in \R^{(k+k^*) \times
(k+k^*)}$, the dynamics \eqref{eq:discretedynamics} are then equivalent to the
iterations, for $t=0,1,2,\ldots$
\begin{align}
\W^t&=\X\V^t-\sum_{r=0}^{t-1} \Y^r (A^{t,r})^\top,\label{eq:AMPW}\\
\Y^t&=\left[{-}\frac{\bar\eta^{t\delta}}{\bar\kappa}
f\left(\sum_{q=0}^t \left(\W^q+\sum_{r=0}^{q-1} \Y^r
(A^{q,r})^\top\right),\beps\right) \odot \z_\delta^t,\;0\right],\label{eq:AMPY}\\
\U^t&=\X^\top \Y^t-\sum_{q=0}^t \V^q (B^{t,q})^\top,\label{eq:AMPU}\\
\V^{t+1}&=\U^t+\sum_{q=0}^t \V^q (B^{t,q})^\top
-\left[\delta\bar\eta^{\delta t}g\left(\sum_{q=0}^t
\V_1^q\right),\,0\right].\label{eq:AMPV}
\end{align}

We may identify \eqref{eq:AMPY} and \eqref{eq:AMPV} as two mappings
\begin{equation}\label{eq:ytvt}
\Y^t=y_t(\W^0,\ldots,\W^t,\beps,\z_\delta^0,\ldots,\z_\delta^t),
\qquad \V^{t+1}=v_{t+1}(\U^0,\ldots,\U^t,\V^0)
\end{equation}
defined recursively for $t=0,1,2,\ldots$ via
\begin{align}
y_t(W^0,\ldots,W^t,\eps,z_\delta^0,\ldots,z_\delta^t)
&=\begin{pmatrix} {-}\frac{\bar\eta^{t\delta}}{\bar\kappa}z_\delta^t 
f\left(\sum_{q=0}^t \left(W^q+\sum_{r=0}^{q-1}
A^{q,r}y_r(W^0,\ldots,W^r,\eps,z_\delta^0,\ldots,z_\delta^r)\right),\eps\right)
\\ 0 \end{pmatrix},\label{eq:yt}\\
v_{t+1}(U^0,\ldots,U^t,V^0)
&=U^t+B^{t,0}V^0+\sum_{q=1}^t B^{t,q}v_q(U^0,\ldots,U^{q-1},V^0)\notag\\
&\hspace{1in}-\begin{pmatrix} \delta\bar\eta^{\delta t}g\left(V_1^0+\sum_{q=1}^t
[v_q(U^0,\ldots,U^{q-1},V^0)]_1\right) \\ 0 \end{pmatrix},\label{eq:vt}
\end{align}
and applied row-wise to their inputs in \eqref{eq:ytvt}.
Let $\theta^0,\theta^*,\eps,\{z_\delta^t\}_{t \geq 0}$ be as defined in the
discrete-time DMFT system. We set
\[V^0=(\theta^0,\theta^*) \in \R^{k+k^*}\]
and define recursively for $t=0,1,2,\ldots$
\begin{equation}\label{eq:AMPstateevolution}
\begin{gathered}
A^{t,q}=\E\D_q V^t \text{ for all } q \in \{0,\ldots,t-1\},
\quad \Sigma^t=\E[(V^0,\ldots,V^t) \otimes (V^0,\ldots,V^t)],\\
(W^0,\ldots,W^t) \sim \cN(0,\Sigma^t) \text{ (independent of
$\eps,\{z_\delta^t\}_{t \geq 0}$)}, \quad
Y^t=y_t(W^0,\ldots,W^t,\eps,z_\delta^0,\ldots,z_\delta^t),\\
B^{t,q}=\gamma\,\E\D_q Y^t \text{ for all } s \in \{0,\ldots,t\}, \quad
\Omega^t=\gamma\,\E[(Y^0,\ldots,Y^t) \otimes (Y^0 \ldots,Y^t)],\\
(U^0,\ldots,U^t) \sim \cN(0,\Omega^t) \text{ (independent of $\theta^0,\theta^*$)},
\quad V^{t+1}=v_{t+1}(U^0,\ldots,U^t,V^0),
\end{gathered}
\end{equation}
where $\D_q V^t \equiv \D_{U^q} v_t(U^0,\ldots,U^{t-1},V^0) \in \R^{(k+k^*)
\times (k+k^*)}$ denotes the partial derivative (Jacobian) in $U^q$, and
similarly $\D_q Y^t \equiv \D_{W^q}
y_t(W^0,\ldots,W^t,\eps,z_\delta^0,\ldots,z_\delta^t) \in
\R^{(k+k^*) \times (k+k^*)}$ denotes that in $W^q$. Then, using these choices of
$\{A^{t,q},B^{t,s}\}$, (\ref{eq:AMPW}--\ref{eq:AMPV}) is a standard form of an
AMP algorithm with state evolution covariances $\{\Sigma^t,\Omega^t\}$, see
e.g.\ \cite{javanmard2013state}.

We apply the result of \cite[Theorem 2.21]{wang2024universality} for this AMP
algorithm: It is clear that
$T$ iterations of (\ref{eq:AMPW}--\ref{eq:AMPV}) may be directly mapped to
$(k+k^*)T$ vector iterations of
\cite[Eqs.\ (2.14) and (D.1--D.2)]{wang2024universality} with
side information vectors $(\beps,\z_\delta^0,\ldots,\z_\delta^T)$
and the notational identification $(1/\sqrt{\gamma})\X \leftrightarrow \W^\top$.
Each function
$y_t(\cdot)$ and $v_{t+1}(\cdot)$ recursively defined via \eqref{eq:AMPY} and
\eqref{eq:AMPV} has polynomial growth and is Lipschitz in its first $t+1$
arguments $W^0,\ldots,W^t$ and $U^0,\ldots,U^t$, by the conditions for
$f(\cdot)$ and $g(\cdot)$ in Assumption \ref{asp:lipschitz}.\footnote{Here,
$(W^0,\ldots,W^t) \mapsto y(W^0,\ldots,W^t,\eps,z_\delta^0,\ldots,z_\delta^t)$
is $(\prod_{s=0}^t C\max(1,|z_\delta^s|))$-Lipschitz where $C>0$ depends on the
Lipschitz constant of $f$. The argument of \cite[Theorem
2.21]{wang2024universality} implicitly assumes that
$(W^0,\ldots,W^t) \mapsto y(W^0,\ldots,W^t,\eps,z_\delta^0,\ldots,z_\delta^t)$
is $L$-Lipschitz for a deterministic constant $L>0$, but it is clear from the
proof that the argument holds as long as this is
$L(\eps,z_\delta^0,\ldots,z_\delta^t)$-Lipschitz where
$\E|L(\eps,z_\delta^0,\ldots,z_\delta^t)|^{2+\iota}<\infty$ for some $\iota>0$.}
The conditions of \cite[Definition 2.18]{wang2024universality} hold for
the data matrix
$(1/\sqrt{\gamma})\X$ by Assumption \ref{asp:model}, with constant 
variance profile across all entries, and \cite[Assumption
2.17]{wang2024universality} is verified by Lemma \ref{lem:initialconditions}.
Then by \cite[Theorem 2.21]{wang2024universality}, for any fixed $T \geq 1$,
almost surely as $n,d\to\infty$,
\begin{equation}\label{eq:state_evolution}
\begin{gathered}
\frac{1}{d}\sum_{j=1}^d \delta_{(U_j^0,\ldots,U_j^T,V_j^0,\ldots,V_j^T)} \Rightarrow
\Law(U^0,\ldots,U^T,V^0,\ldots,V^T),\\
\frac{1}{n}\sum_{i=1}^n
\delta_{(W_i^0,\ldots,W_i^T,Y_i^0,\ldots,Y_i^T,\eps_i,z_{\delta,i}^0,\ldots,z_{\delta,i}^T)} \Rightarrow
\Law(W^0,\ldots,W^T,Y^0,\ldots,Y^T,\eps,z_\delta^0,\ldots,z_\delta^T)
\end{gathered}
\end{equation}
weakly and in Wasserstein-2.

Finally, we show that this implies the statements of the lemma.
Writing the two components of $V^t,Y^t \in \R^{k \times k^*}$ as
$V^t=(V_1^t,V_2^t)$ and $Y^t=(Y_1^t,Y_2^t)$, note that by the definitions
\eqref{eq:yt} and \eqref{eq:vt} we have
\[V^t=(V_1^t,0) \text{ for all } t \geq 1,
\quad Y^t=(Y_1^t,0) \text{ for all } t \geq 0,\]
and hence also
\[A^{t,q}=\begin{pmatrix} (A^{t,q})_{11} & (A^{t,q})_{12} \\
0 & 0 \end{pmatrix}, \; W_2^t=0 \text{ for all } t \geq 1,\]
\[B^{t,q}=\begin{pmatrix} (B^{t,q})_{11} & (B^{t,q})_{12} \\
0 & 0 \end{pmatrix}, \; U_2^t=0 \text{ for all } t \geq 0.\]
Let us define
\begin{align}
(\theta_\delta^t,\theta^*)&:=\sum_{r=0}^t V^r
=\left(\sum_{r=0}^t V_1^r,\,V_2^0\right),\label{eq:AMPthetadef}\\
(\xi_\delta^t,w^*)&:=\sum_{q=0}^t \left(W^q+\sum_{r=0}^{q-1}
A^{q,r}Y^r\right)
=\left(\sum_{q=0}^t W_1^q+\sum_{q=1}^t \sum_{r=0}^{q-1} (A^{q,r})_{11}Y_1^r,\,
W_2^0\right).\label{eq:AMPxidef}
\end{align}
Comparing these definitions with
\eqref{eq:thetaV} and \eqref{eq:AMPW}, the convergence
\eqref{eq:state_evolution} implies that also
\begin{equation}\label{eq:AMPDMFTlimits}
\frac{1}{d}\sum_{j=1}^d \delta_{(\theta^0_{\delta,j}, \ldots,
\theta^T_{\delta,j},\theta^\ast_j)} \Rightarrow
\Law\big(\theta^0_\delta,\ldots, \theta^T_\delta,\theta^\ast\big),
\quad \frac{1}{n}\sum_{i=1}^n \delta_{(\xi^0_{\delta,i}, \ldots,
\xi^T_{\delta,i},\xi_i^*,\eps_i)}
\Rightarrow \Law\big(\xi^0_\delta,\ldots,\xi^T_\delta,w^*,\eps\big).
\end{equation}
So it remains to check that these definitions of
$\theta_\delta^t,\xi_\delta^t,w^*$ in (\ref{eq:AMPthetadef}--\ref{eq:AMPxidef})
coincide with those defined by the discrete DMFT recursions
(\ref{eq:def-dis-theta}--\ref{eq:def-dis-Gamma}).

Differentiating \eqref{eq:yt}
and \eqref{eq:vt} by the chain rule and applying $A^{q,r}=\E \D_rV^q$
and $B^{t,s}=\gamma\,\E \D_s Y^t$, observe that
\begin{align}
\D_s Y^t&:=\begin{pmatrix} (\D_s Y^t)_{11} & (\D_s Y^t)_{12} \\ 0 & 0
\end{pmatrix}\notag\\
&={-}\frac{\bar\eta^{t\delta}}{\bar\kappa}z_\delta^t \begin{pmatrix} 
\D_\xi f(\xi_\delta^t,w^*,\eps)
& \D_{w^*} f(\xi_\delta^t,w^*,\eps)\\ 0 & 0 \end{pmatrix}
\left(\Id_{k+k^*}+\sum_{q=s+1}^t \sum_{r=s}^{q-1}(\E \D_r V^q)(\D_s
Y^r)\right),\label{eq:AMPchainrule1}\\
\D_s V^{t+1}&:=\begin{pmatrix} (\D_s V^{t+1})_{11} & (\D_s V^{t+1})_{12}
\\ 0 & 0 \end{pmatrix}\notag\\
&=\Id_{k+k^*}\1_{s=t}
+\sum_{q=s+1}^t (\gamma \E \D_q Y^t)(\D_s V^q)\notag\\
&\hspace{1in}- \delta \bar\eta^{\delta t}\begin{pmatrix}\D g(\theta_\delta^t)
\sum_{q=s+1}^t (\D_s V^q)_{11} & \D g(\theta_\delta^t)
\sum_{q=s+1}^t (\D_s V^q)_{12} \\ 0 & 0 \end{pmatrix}.
\label{eq:AMPchainrule2}
\end{align}
In particular, for $s=t$ we have $(\D_t
Y^t)_{11}={-}\frac{\bar\eta^{t\delta}}{\bar\kappa}z_\delta^t \cdot \D_\xi
f(\xi_\delta^t,w^*,\eps)$. Let us define
\begin{equation}\label{eq:AMPresponsedefs}
\begin{gathered}
(\D_sY^t-\D_{s+1}Y^t)_{11}
:={-}\frac{\bar\eta^{t\delta}}{\bar\kappa}z_\delta^t r_{f,\delta}^{t,s},
\quad (\D_0Y^t)_{12}:=
{-}\frac{\bar\eta^{t\delta}}{\bar\kappa}z_\delta^t r_{f,\delta}^{t,*},
\quad R_{f,\delta}^{t,s}=\E r_{f,\delta}^{t,s},
\quad R_{f,\delta}^{t,*}=\E r_{f,\delta}^{t,*},\\
\sum_{r=s+1}^t (\D_s V^r)_{11}
=\sum_{r=s}^{t-1} (\D_s V^{r+1})_{11}:=r_{\theta,\delta}^{t,s},
\quad R_{\theta,\delta}^{t,s}=\E r_{\theta,\delta}^{t,s},
\quad \Gamma_\delta^t=\E \D_\xi f(\xi_\delta^t,w^*,\eps),
\end{gathered}
\end{equation}
and apply the identities
\begin{align}
\sum_{q=s+1}^t \sum_{r=s}^{q-1} [(\E \D_r V^q)(\D_s Y^r)]_{1i}
&=\sum_{r=s}^{t-1}\left(\sum_{q=r+1}^t \E(\D_r V^q)_{11} \right)(\D_s
Y^r)_{1i}\notag\\
&=\sum_{r=s}^{t-1} R_{\theta,\delta}^{t,r}(\D_s Y^r)_{1i}
\text{ for } i=1,2,\label{eq:sumidentity1}\\
\sum_{q=s+1}^t [(\E \D_q Y^t)(\D_s V^q)]_{11}
&=\sum_{q=s+1}^t \left(\E(\D_t Y^t)_{11}+\sum_{r=q}^{t-1}
\E(\D_r Y^t-\D_{r+1} Y^t)_{11}\right)(\D_s V^q)_{11}\notag\\
&=\E(\D_t Y^t)_{11} \sum_{q=s+1}^t (\D_s V^q)_{11}
+\sum_{r=s+1}^{t-1}
\E(\D_r Y^t-\D_{r+1} Y^t)_{11}\left(\sum_{q=s+1}^r (\D_s V^q)_{11}\right)\notag\\
&={-}\delta\bar\eta^{t\delta} \Gamma_\delta^t r_{\theta,\delta}^{t,s}
-\delta\bar \eta^{t\delta}\sum_{r=s+1}^{t-1} R_{f,\delta}^{t,r}
r_{\theta,\delta}^{r,s}.
\label{eq:sumidentity2}
\end{align}
In the last step, we have used $\E z_\delta^t=\delta\bar\kappa$ and the fact
that $z_\delta^t$ is independent of
$\D_\xi f(\xi_\delta^t,w^*,\eps)$ and $r_{f,\delta}^{t,s}$.
Then, applying \eqref{eq:AMPresponsedefs},
\eqref{eq:sumidentity1}, \eqref{eq:sumidentity2} to
(\ref{eq:AMPchainrule1}--\ref{eq:AMPchainrule2}), we obtain the recursions
\begin{align*}
r_{f,\delta}^{t,s}&=\D_\xi f(\xi_\delta^t,w^*,\eps)
\left(\sum_{r=s+1}^{t-1}
R_{\theta,\delta}^{t,r}(\D_s Y^r-\D_{s+1}Y^r)_{11}
+R_{\theta,\delta}^{t,s}(\D_s Y^s)_{11}\right)\\
&={-}\D_\xi f(\xi_\delta^t,w^*,\eps)
\left(\sum_{r=s+1}^{t-1}
\frac{\bar\eta^{r\delta}}{\bar\kappa}
R_{\theta,\delta}^{t,r}r_{f,\delta}^{r,s}z_\delta^r
+\frac{\bar\eta^{s\delta}}{\bar\kappa} R_{\theta,\delta}^{t,s} \cdot
\D_\xi f(\xi_\delta^s,w^*,\eps)z_\delta^s\right),\\
r_{f,\delta}^{t,*}&=
\D_{w^*} f(\xi_\delta^t,w^*,\eps)
+\D_\xi f(\xi_\delta^t,w^*,\eps)\sum_{r=0}^{t-1}R_{\theta,\delta}^{t,r}
(\D_0 Y^r)_{12}\\
&={-}\D_\xi f(\xi_\delta^t,w^*,\eps)
\sum_{r=0}^{t-1} \frac{\bar\eta^{r\delta}}{\bar\kappa}
R_{\theta,\delta}^{t,r}r_{f,\delta}^{r,*}z_\delta^r
+\D_{w^*} f(\xi_\delta^t,w^*,\eps),\\
r_{\theta,\delta}^{t,s}&=\sum_{r=s}^{t-1} \left(
\Id_k\1_{s=r}-\gamma \delta\bar\eta^{r\delta} \Gamma_\delta^r
r_{\theta,\delta}^{r,s}
-\gamma\delta\bar\eta^{r\delta}\sum_{q=s+1}^{r-1}
R_{f,\delta}^{r,q}r_{\theta,\delta}^{q,s}\right)
-\delta \bar\eta^{r\delta}
\D g(\theta_\delta^r)r_{\theta,\delta}^{r,s}\\
&=\Id_k-\sum_{r=s}^{t-1} \delta\bar\eta^{r\delta}
\bigg[(\gamma \Gamma_\delta^r+\D g(\theta_\delta^r)) \cdot
r_{\theta,\delta}^{r,s}
+\gamma \sum_{q=s+1}^{r-1} R_{f,\delta}^{r,q}r_{\theta,\delta}^{q,s}\bigg]
\end{align*}
which are precisely (\ref{eq:def-dis-deriv-t}--\ref{eq:def-dis-deriv-xi}).
Furthermore, let us define
\begin{equation}\label{eq:AMPuwdef}
\sum_{r=0}^t U_1^r:=\sqrt{\gamma}u_\delta^t,
\qquad \sum_{q=0}^t W_1^q:=w_\delta^t.
\end{equation}
Then by \eqref{eq:AMPthetadef} and \eqref{eq:vt},
\begin{align*}
\theta_\delta^t=V_1^0+\sum_{r=1}^t V_1^r
&=\theta^0+\sum_{r=0}^{t-1}\left(U_1^r+\sum_{q=0}^r (B^{r,q})_{11}V_1^r
+(B^{r,0})_{12}V_2^0-\delta\bar\eta^{r\delta} g(\theta_\delta^r)\right)\\
&=\theta^0+\sum_{r=0}^{t-1} \left(U_1^r+\gamma\sum_{q=0}^r
(\E\D_q Y^r)_{11}V_1^r+\gamma (\E \D_0 Y^r)_{12}V_2^0
-\delta\bar\eta^{r\delta} g(\theta_\delta^r)\right)\\
&\overset{(*)}{=}\theta^0+\sum_{r=0}^{t-1}\left(U_1^r+\gamma
\left({-}\delta\bar\eta^{r\delta}\Gamma_\delta^r \theta_\delta^r-
\delta\bar\eta^{r\delta}
\sum_{q=0}^{r-1} R_{f,\delta}^{r,q}
\theta_\delta^q -\delta\bar\eta^{r\delta} R_{f,\delta}^{t,*}\theta^*\right)
-\delta\bar\eta^{r\delta} g(\theta_\delta^r)\right)\\
&=\theta^0-\sum_{r=0}^{t-1}\delta\bar\eta^{r\delta}
\left(\gamma \Gamma_\delta^r \theta_\delta^r
+g(\theta_\delta^r)+\gamma \sum_{q=0}^{r-1}
R_{f,\delta}^{r,q}\theta_\delta^q\right)+\sqrt{\gamma} u_\delta^t,
\end{align*}
where $(*)$ applies an identity analogous to \eqref{eq:sumidentity2}. Similarly,
applying \eqref{eq:AMPxidef} and \eqref{eq:yt},
\begin{align*}
\xi_\delta^t=\sum_{q=0}^t W_1^q+\sum_{q=1}^t \sum_{r=0}^{q-1}
(A^{q,r})_{11}Y_1^r
&=\sum_{q=0}^t W_1^q+\sum_{q=1}^t \sum_{r=0}^{q-1}
(\E \D_r V^q)_{11}Y_1^r\\
&\overset{(*)}{=}
\sum_{q=0}^t W_1^q+\sum_{r=0}^{t-1} R_{\theta,\delta}^{t,r} Y_1^r
={-}\sum_{r=0}^{t-1} \frac{\bar\eta^{r\delta}}{\bar\kappa}
R_{\theta,\delta}^{t,r}f(\xi_\delta^r,w^*,\eps)z_\delta^r+w_\delta^t,
\end{align*}
where $(*)$ applies an identity analogous to \eqref{eq:sumidentity1}. These
recursions are precisely (\ref{eq:def-dis-theta}--\ref{eq:def-dis-xi}).
Finally, by \eqref{eq:AMPstateevolution},
we note that $\{u_\delta^t\}_{t \geq 0}$ and $(\{w_\delta^t\}_{t \geq 0},
w_\delta^*)$ defined via \eqref{eq:AMPxidef} and \eqref{eq:AMPuwdef} are
Gaussian with covariance
\begin{align*}
\E[u_\delta^t \otimes u_\delta^s]
&=\gamma^{-1}\E\left[\sum_{r=0}^t U_1^r \otimes \sum_{r=0}^s U_1^r\right]
=\E\left[\sum_{r=0}^t Y_1^t \otimes \sum_{r=0}^s Y_1^r\right]\\
&=\E\left[\sum_{r=0}^t \frac{\bar\eta^{r\delta}}{\bar\kappa}
f(\xi_\delta^r,w^*,\eps)z_\delta^r
\otimes \sum_{r=0}^s \frac{\bar\eta^{r\delta}}{\bar\kappa}
f(\xi_\delta^r,w^*,\eps)z_\delta^r\right],\\
\E[w_\delta^t \otimes w_\delta^s]
&=\E\left[\sum_{q=0}^t W_1^q \otimes \sum_{q=0}^s W_1^q\right]
=\E\left[\sum_{q=0}^t V_1^q \otimes \sum_{q=0}^s V_1^q\right]
=\E[\theta_\delta^t \otimes \theta_\delta^s],\\
\E[w_\delta^t \otimes w_\delta^*]
&=\E\left[\sum_{q=0}^t W_1^q \otimes W_2^0\right]
=\E\left[\sum_{q=0}^t V_1^q \otimes V_2^0\right]
=\E[\theta^t \otimes \theta^*],\\
\E[w_\delta^* \otimes w_\delta^*]
&=\E\left[W_2^0 \otimes W_2^0\right]
=\E\left[V_2^0 \otimes V_2^0\right]
=\E[\theta^* \otimes \theta^*],
\end{align*}
which is precisely (\ref{eq:def-dis-C-t}--\ref{eq:def-dis-C-h}). This verifies
that the limits in \eqref{eq:AMPDMFTlimits} coincide with the definitions from
(\ref{eq:def-dis-theta}--\ref{eq:def-dis-Gamma}), showing the lemma.
\end{proof}

\subsection{Convergence to the continuous DMFT system}\label{sec:deltaDMFTconvergence}

We now show that a continuous-time embedding of the preceding
discrete-time DMFT system converges to (the unique fixed
point of) the continuous-time DMFT system defined in Section
\ref{sec:DMFTlimits}, as the discretization step size $\delta \to 0$.

To ease notation in this section, for all times $t \in [0,T]$,
we will use the time-discretization conventions
\begin{equation}\label{eq:flrcil}
\flr{t} = \max\{i\delta: i\delta\leq t, i\in \Z_+\} \in \delta\Z_+,
\quad \cil{t} = \flr{t} + \delta \in \delta\Z_+
\end{equation}
Let $\{\theta_\delta^t\}_{t \in \Z_+}$, $\{\xi_\delta^t\}_{t \in \Z_+}$,
and the correlation and response kernels
\[C_{\theta,\delta},R_{\theta,\delta},C_{f,\delta},R_{f,\delta},R_{f,\delta}^*,\Gamma_\delta\]
be the components of the discrete-time DMFT system defined via the recurions
(\ref{eq:def-dis-theta}--\ref{eq:def-dis-Gamma}).
We define their continuous-time embeddings
$\{\bar\theta_\delta^t\}_{t \in [0,T]}$, 
$\{\bar \xi_\delta^t\}_{t \in [0,T]}$,
$\{\bar C_{\theta,\delta}^{t,s}\}_{t,s \in [0,T]}$,
$\{\bar C_{f,\delta}^{t,s}\}_{t,s \in [0,T]}$,
$\{\bar R_{f,\delta}^{t,*}\}_{t \in [0,T]}$, and
$\{\bar \Gamma_\delta^t\}_{t \in [0,T]}$ by
\begin{equation}\label{eq:embeddings}
\bar\theta_\delta^t=\theta_\delta^k,
\; \bar\xi_\delta^t=\xi_\delta^k,
\; \bar C_{\theta,\delta}^{t,s}=C_{\theta,\delta}^{k,j},
\; \bar C_{f,\delta}^{t,s}=C_{f,\delta}^{k,j},
\; \bar R_{f,\delta}^{t,*}=R_{f,\delta}^{k,*},
\; \bar \Gamma_{f,\delta}^t=\Gamma_{f,\delta}^k
\text{ if } \flr{t}=k \text{ and } \flr{s}=j.
\end{equation}
These embeddings are piecewise-constant and right-continuous with jumps at
$\delta\Z_+$ (henceforth abbreviated as $\delta$-p.c.r.c.)
We define also the $\delta$-p.c.r.c.\ embedding
$\{\bar R_{\theta,\delta}^{t,s}\}_{t,s \in [0,T]}$ by
\[\bar R_{\theta,\delta}^{t,s}=\begin{cases} 0 & \text{ if } s>t \\
\Id & \text{ if } s \leq t \text{ and } \flr{s}=\flr{t} \\
R_{\theta,\delta}^{k,j} & \text{ if } \flr{s}=j<\flr{t}=k,\end{cases}\]
and a family of linear operators $\{\bar R_{f,\delta}^t\}_{t \in [0,T]}$
with $\bar R_{f,\delta}^t:L^4([0,t],\R^k) \to \R^k$ by
\begin{equation}\label{eq:Rfembed}
\bar R_{f,\delta}^t(x)
=\int_0^{\flr{t}} \bar R_{f,\delta}^{t,s} x^s\d s,
\quad \bar R_{f,\delta}^{t,s}=\delta^{-1}R_{f,\delta}^{k,j} \text{ if }
\flr{s}=j<\flr{t}=k.
\end{equation}
We denote the tuples of these embedded correlation and response functions by
\[X^\delta=(\bar C_{f,\delta},\bar R_{f,\delta},
\bar R_{f,\delta}^*,\bar \Gamma_\delta),
\quad Y^\delta=(\bar C_{\theta,\delta},\bar R_{\theta,\delta}).\]

The main result of this section is the following lemma, where we recall
the domain $\cS(T,C_0)$ and metric $\lbddst{\cdot}{\cdot}$ defined in Section
\ref{sec:existenceuniqueness}.

\begin{lem}\label{lem:dmft_discretization_error}
\begin{enumerate}[(a)]
\item For any sufficiently large constant $C_0 \equiv C_0(T)>0$ and any
$\delta>0$,
\[(Y^\delta,X^\delta) \in \cS(T,C_0).\]
\item For any sufficiently large constant $\lambda \equiv \lambda(T,C_0)>0$,
\[\lim_{\delta \to 0} \lbddst{X^\delta}{X}=0,
\quad \lim_{\delta \to 0} \lbddst{Y^\delta}{Y}=0\]
where $(Y,X) \in \cS^\text{cont}(T,C_0)$ is the fixed point given by Theorem
\ref{thm:uniqueness-existence}.
\item For any fixed $m \geq 0$ and $0 \leq t_1 \leq \ldots \leq t_m \leq T$,
as $\delta \to 0$,
\begin{align}
\Law(\bar \theta^{t_1}_\delta,\ldots,\bar \theta^{t_m}_\delta,\theta^*)
&\Rightarrow
\Law(\theta^{t_1},\ldots,\theta^{t_m},\theta^*),\,\label{eq:discretedmftthetaconv}\\
\Law(\bar \xi_\delta^{t_1},\ldots,\bar \xi_\delta^{t_m},w^*,\eps) &\Rightarrow
\Law(\xi^{t_1},\ldots,\xi^{t_m},w^*,\eps)\label{eq:discretedmftetaconv}
\end{align}
weakly and in Wasserstein-2.
\end{enumerate}
\end{lem}

\begin{proof}[Proof of Lemma \ref{lem:dmft_discretization_error}(a)]
We define also the embeddings
\begin{equation}\label{eq:embeddings2}
\bar u_\delta^t=u_\delta^k,
\;\bar w_\delta^t=w_\delta^k,
\;\bar r_{f,\delta}^{t,*}=r_{f,\delta}^{k,*} \text{ if } \flr{t}=k,
\quad
\bar r_{\theta,\delta}^{t,s}=\begin{cases} 0 & \text{ if } s>t \\
\Id & \text{ if } s \leq t \text{ and } \flr{s}=\flr{t} \\
r_{\theta,\delta}^{k,j} & \text{ if } \flr{s}=k<\flr{t}=j.\end{cases}
\end{equation}
Let $D([0,T],\R^k)$ denote the space of $\delta$-p.c.r.c.\ $\R^k$-valued
processes $\{x^t\}_{t \in [0,T]}$, i.e.\ satisfying $x^t=x^{\flr{t}}$
for all $t \in [0,T]$. We embed
$r_{f,\delta}^{k,j}$ as a family of linear operators
$\bar r_{f,\delta}^t:D([0,T],\R^k) \to \R^k$ given by
\begin{equation}\label{eq:rfembed}
\bar r_{f,\delta}^t(x^{[t]})
=\sum_{j=0}^{k-1} r_{f,\delta}^{k,j}x^{j\delta} \text{ if } \flr{t}=k.
\end{equation}
Finally, we identify the variables $\{z_\delta^t\}_{t \geq 0}$ of the
discrete-time DMFT system with the increments of the continuous-time Gaussian
or Poisson process $\{z^t\}_{t \geq 0}$ defined by
\eqref{eq:poisson-z}/\eqref{eq:gaussian-z},
\[z_\delta^t=z^{(t+1)\delta}-z^{t\delta}.\]
Then the discrete-time DMFT recursions
(\ref{eq:def-dis-theta}--\ref{eq:def-dis-Gamma})
imply the following equations for these embeddings:
For any $x \in D([0,T],\R^k)$,
\begin{align}
\bar\theta^t_\delta&=\bar\theta^0_\delta-\int_0^{\flr{t}} \bar\eta^{\flr{r}} \Big({\gamma}\bar\Gamma^r_\delta
\bar\theta^r_\delta+g(\bar\theta^r_\delta)+{\gamma}\bar R_{f,\delta}^r(\bar\theta^{[r]}_\delta)+
{\gamma}\bar R_{f,\delta}^{r,*}\theta^*\Big)\d r+\sqrt{\gamma}\,\bar u^t_\delta,
\label{eq:def-dist-theta}\\
\bar\xi^t_\delta&={-}\int_0^{\flr{t}} \frac{\bar\eta^{\flr{r}}}{\bar\kappa} \bar R_{\theta,\delta}^{t,r}
f(\bar\xi^{r-}_\delta,w_\delta^*,\eps)\d z^r+\bar
w^t_\delta,\label{eq:def-dist-xi}\\
\bar r_{\theta,\delta}^{t,s}&=\Id- 
\mathbf{1}_{\cil{s} \leq \flr{t}} \int_{\cil{s}}^{\flr{t}} \bar\eta^{\flr{r}}\bigg[\Big({\gamma}\bar\Gamma^r_\delta+\D
g(\bar\theta^r_\delta)\Big)\bar r_{\theta,\delta}^{r,s}+{\gamma}
\bar R_{f,\delta}^r\Big(\bar r_{\theta,\delta}^{[r],s}\Big)\bigg]\d r
\text{ for } s \leq t,\notag\\
&\hspace{1in}r_{\theta,\delta}^{t,s}=0 \text{ for } s>t,
\label{eq:def-dist-deriv-t}\\
\bar r_{f,\delta}^{t,*}&={-}\D_\xi f(\bar\xi^t_\delta,w_\delta^*,\eps) \int_0^{\flr{t}}
\frac{\bar\eta^{\flr{r}}}{\bar\kappa}
\bar R_{\theta,\delta}^{t,r}\bar r_{f,\delta}^{r-,*}\d z^r
+\D_{w^*} f(\bar\xi^t_\delta,w_\delta^*,\eps),\label{eq:def-dist-deriv-xi-star}\\
\bar r_{f,\delta}^t(x^{[t]})&={-}\D_\xi
f(\bar\xi^t_\delta,w_\delta^*,\eps) \int_0^{\flr{t}} \frac{\bar\eta^{\flr{r}}}{\bar\kappa}
\bar R_{\theta,\delta}^{t,r}\bigg(\bar r_{f,\delta}^{r-}(x^{[r]})+\D_\xi
f(\bar\xi^{r-}_\delta,w_\delta^*,\eps)x^{r-}\bigg)\d z^r,\label{eq:def-dist-resp-xi}
\end{align}
and
\begin{align}
\bar C_{\theta,\delta}^{t,s}&=\E[\bar\theta^t_\delta \otimes \bar\theta^s_\delta] \text{ for } t,s \in [0,\infty)
\cup \{*\}\label{eq:def-dist-C-t}\\
\bar R_{\theta,\delta}^{t,s}&=\E[\bar r_{\theta,\delta}^{t,s}]
\label{eq:def-dist-R-t}\\
\bar C_{f,\delta}^{t,s}&=\E\brk{\int_0^{\flr{t}}
\frac{\bar\eta^{\flr{r}}}{\bar\kappa}f(\bar\xi^{r-}_\delta,w_\delta^*,\eps)\d z^r
\otimes \int_0^{\flr{s}}
\frac{\bar\eta^{\flr{r}}}{\bar\kappa}f(\bar\xi^{r-}_\delta,w_\delta^*,\eps)\d z^r}\label{eq:def-dist-C-h}\\
\bar R_{f,\delta}^t(x^{[t]})&=\E\big[\bar r_{f,\delta}^t(x^{[t]})\big]\label{eq:def-dist-R-h}\\
\bar R_{f,\delta}^{t,*}&=\E[\bar r_{f,\delta}^{t,*}]
\label{eq:def-dist-R-h-star}\\
\bar \Gamma^t_\delta&=\E[\D_\xi f(\bar\xi^t_\delta,w_\delta^*,\eps)].\label{eq:def-dist-Gamma}
\end{align}
For example, since $\{\bar\theta_\delta^t\}$ is $\delta$-p.c.r.c., we have
by \eqref{eq:Rfembed} for $\flr{t}=k$ that
\[\bar R_{f,\delta}^t(\bar\theta_\delta^{[t]})
=\int_0^{\flr{t}} \bar R_{f,\delta}^{t,s}\bar\theta_\delta^s\d s
=\delta\sum_{q=0}^{k-1} (\delta^{-1}R_{f,\delta}^{k,q})\theta_\delta^q
=\sum_{q=0}^{k-1} R_{f,\delta}^{k,q} \theta_\delta^q.\]
Then, since $\{\bar\theta_\delta^t,\bar u_\delta^t\}$ are both
$\delta$-p.c.r.c.\ and defined by \eqref{eq:embeddings} and
\eqref{eq:embeddings2}, we have that \eqref{eq:def-dist-theta} is equivalent to
the discrete-time equation \eqref{eq:def-dis-theta},
\[\theta_\delta^k=\theta_\delta^0
-\delta \sum_{r=0}^k
\bar\eta^{r\delta}\Big(\gamma \Gamma_\delta^r\theta_\delta^r
+g(\theta_\delta^r)+\gamma \sum_{q=0}^{r-1} R_{f,\delta}^{r,q}
\theta_\delta^q+\gamma R_{f,\delta}^{r,*}\theta^*\Big)
+\sqrt{\gamma}u_\delta^k.\]
Similar arguments show that
(\ref{eq:def-dis-xi}--\ref{eq:def-dis-deriv-xi-star}) are equivalent to
(\ref{eq:def-dist-xi}--\ref{eq:def-dist-deriv-xi-star}), and multiplying
\eqref{eq:def-dis-deriv-xi} on both sides by $x^s$ and summing over
$s=0,\ldots,t-1$ shows \eqref{eq:def-dist-resp-xi}. It is clear from the
definitions that
(\ref{eq:def-dis-C-t}--\ref{eq:def-dis-Gamma}) are equivalent to
(\ref{eq:def-dist-C-t}--\ref{eq:def-dist-Gamma}), where \eqref{eq:def-dist-R-h}
follows from \eqref{eq:Rfembed} and \eqref{eq:rfembed}.

Now let $\cD_\theta^\delta$ denote the space of all tuples
$Y^\delta=(\bar C_{\theta,\delta},\bar R_{\theta,\delta})$ where
\begin{itemize}
\item $\bar C_{\theta,\delta} \equiv \{\bar C_{\theta,\delta}^{t,s}\}_{t,s \in
[0,T]}$ is a covariance kernel on $(\R^k \otimes [0,T]) \times \R^{k^*}$
that is $\delta$-p.c.r.c.\ in $(t,s)$.
\item $\bar R_{\theta,\delta} \equiv \{\bar R_{\theta,\delta}^{t,s}\}_{t,s \in
[0,T]}$ is a $\R^{k \times k}$-valued process satisfying
$\bar R_{\theta,\delta}^{t,s}=0$ if $s>t$,
$\bar R_{\theta,\delta}^{t,s}=\Id$ if $s \leq t$ with $\flr{s}=\flr{t}$,
and $\bar R_{\theta,\delta}^{t,s}$ is $\delta$-p.c.r.c.\ in $(t,s)$.
\end{itemize}
Likewise, let $\cD_\xi^\delta$ denote the space of all tuples
$X^\delta=(\bar C_{f,\delta},\bar R_{f,\delta}, \bar R_{f,\delta}^*,
\bar \Gamma_\delta)$ where
\begin{itemize}
\item $\bar C_{f,\delta} \equiv \{\bar C_{f,\delta}^{t,s}\}_{s,t \in [0,T]}$
is a covariance kernel on $\R^k \otimes [0,T]$ that is $\delta$-p.c.r.c.\ in
$(t,s)$,
\item $\bar R_{f,\delta}^* \equiv \{\bar R_{f,\delta}^{t,*}\}_{t \in [0,T]}$
and $\bar \Gamma_\delta \equiv \{\bar \Gamma_\delta^t\}_{t \in [0,T]}$ are
$\R^{k \times k^*}$-valued and $\R^{k \times k}$-valued processes,
respectively, that are $\delta$-p.c.r.c.\ in $t$,
\item $\bar R_{f,\delta} \equiv \{\bar R_{f,\delta}^t\}_{t \in [0,T]}$ is a
family of linear operators $\bar R_{f,\delta}^t:D([0,t],\R^k) \to \R^k$ having
the forms, for some matrix-valued coefficients $R_{f,\delta}^{k,j}$,
\begin{equation}\label{eq:barRform}
\bar R_{f,\delta}^t(x^{[t]})=
\sum_{j=0}^{k-1} R_{f,\delta}^{k,j}x^{j\delta} \text{ if } \flr{t}=k.
\end{equation}
\end{itemize}
Given any $X^\delta=(\bar C_{f,\delta},\bar R_{f,\delta},
\bar R_{f,\delta}^*,\bar \Gamma_\delta) \in \cD_\xi^\delta$, note that
$\{\bar u_\delta^t\} \sim \GP(0,\bar C_{f,\delta})$ is $\delta$-p.c.r.c.\ with
probability 1. Then it may be checked inductively in $k=0,1,2,\ldots$ that 
\eqref{eq:def-dist-theta} and \eqref{eq:def-dist-deriv-t} have
unique solutions over $t,s \in [0,k\delta)$, which are $\delta$-p.c.r.c. Then
$Y^\delta=(\bar C_{\theta,\delta},\bar R_{\theta,\delta})$ defined by
(\ref{eq:def-dist-C-t}--\ref{eq:def-dist-R-t}) belongs to $\cD_\theta^\delta$,
i.e.\ \eqref{eq:def-dist-theta}, \eqref{eq:def-dist-deriv-t}, and 
(\ref{eq:def-dist-C-t}--\ref{eq:def-dist-R-t}) define a mapping
\[\trsfrmDA:\cD_\xi^\delta \to \cD_\theta^\delta.\]
Similarly, given any
$Y^\delta=(\bar C_{\theta,\delta},\bar R_{\theta,\delta})
\in \cD_\theta^\delta$, it may be checked inductively that 
\eqref{eq:def-dist-xi} and \eqref{eq:def-dist-deriv-xi-star} have unique
solutions which are $\delta$-p.c.r.c.\ in $t$, and that there exist (random)
matrix coefficients $r_{f,\delta}^{k,j}$ such that for any
$x \in D([0,T],\R^k)$, the solution to \eqref{eq:def-dist-resp-xi} takes a form
\[\bar r_{f,\delta}^t(x^{[t]})
=\sum_{j=0}^{k-1} r_{f,\delta}^{k,j}x^{j\delta}
\text{ if } \flr{t}=k.\]
Then $X^\delta=(\bar C_{f,\delta},\bar R_{f,\delta},
\bar R_{f,\delta}^*,\bar \Gamma_\delta)$ defined by
(\ref{eq:def-dist-C-h}--\ref{eq:def-dist-Gamma}) belongs to $\cD_\xi^\delta$,
i.e.\ \eqref{eq:def-dist-xi},
(\ref{eq:def-dist-deriv-xi-star}--\ref{eq:def-dist-resp-xi}), and
(\ref{eq:def-dist-C-h}--\ref{eq:def-dist-Gamma}) define a mapping
\[\trsfrmDB:\cD_\theta^\delta \to \cD_\xi^\delta.\]
These mappings may be understood as
discretized versions of $\trsfrmA$ and $\trsfrmB$ from Section
\ref{sec:existenceuniqueness}. Now denoting by
\[X^\delta=(\bar C_{f,\delta},\bar R_{f,\delta},
\bar R_{f,\delta}^*,\bar \Gamma_\delta),
\quad Y^\delta=(\bar C_{\theta,\delta},\bar R_{\theta,\delta})\]
the specific elements of $\cD_\xi^\delta$ and $\cD_\theta^\delta$ that
correspond to the previously defined embeddings of the discrete-time DMFT
recursions,
since these satisfy (\ref{eq:def-dist-theta}--\ref{eq:def-dist-Gamma}), they
are a fixed point of these mappings, i.e.
\[\trsfrmDA(X^\delta)=Y^\delta,
\quad \trsfrmDB(Y^\delta)=X^\delta.\]

Lemma \ref{lem:dmft_discretization_error}(a) then holds by the following
reasoning: On one hand, $(X^\delta,Y^\delta) \in \cD_\xi^\delta \times
\cD_\theta^\delta$ is the unique such
fixed point of these mappings, because the equations
(\ref{eq:def-dist-theta}--\ref{eq:def-dist-Gamma}) uniquely determine,
recursively in time,
\[\begin{gathered}
\{\bar C_{\theta,\delta}^{t,s},\bar R_{\theta,\delta}^{t,s}\}_{t,s:\flr{t},\flr{s}=0}
\Rightarrow
\{\bar C_{f,\delta}^{t,s},\bar
R_{f,\delta}^t,\bar R_{f,\delta}^{t,*},\bar \Gamma_\delta^t\}_{t,s:\flr{t},\flr{s}=0}
\Rightarrow\\
\{\bar C_{\theta,\delta}^{t,s},\bar
R_{\theta,\delta}^{t,s}\}_{t,s:\flr{t},\flr{s} \in \{0,\delta\}}
\Rightarrow
\{\bar C_{f,\delta}^{t,s},\bar
R_{f,\delta}^t,\bar R_{f,\delta}^{t,*},\bar \Gamma_\delta^t\}_{t,s:\flr{t},\flr{s} \in \{0,\delta\}} \Rightarrow\\
\{\bar C_{\theta,\delta}^{t,s},\bar
R_{\theta,\delta}^{t,s}\}_{t,s:\flr{t},\flr{s} \in \{0,\delta,2\delta\}}
\Rightarrow \{\bar C_{f,\delta}^{t,s},\bar
R_{f,\delta}^t,\bar R_{f,\delta}^{t,*},\bar
\Gamma_\delta^t\}_{t,s:\flr{t},\flr{s} \in \{0,\delta,2\delta\}} \Rightarrow
\ldots
\end{gathered}\]
On the other hand, let us extend each operator
$\bar R_{f,\delta}^t:D([0,t],\R^k) \to \R^k$ having the form
\eqref{eq:barRform} to an operator
$\bar R_{f,\delta}^t:L^4([0,t],\R^k) \to \R^k$ given by
\[\bar R_{f,\delta}^t(x^{[t]})=\int_0^{\flr{t}} \bar R_{f,\delta}^{t,s}x^s\d s,
\quad \bar R_{f,\delta}^{t,s}=\delta^{-1}R_{f,\delta}^{k,j} \text{ if }
\flr{s}=j<\flr{t}=k.\]
In particular, this extension of $\bar R_{f,\delta}^t:D([0,t],\R^k) \to \R^k$
defined via \eqref{eq:def-dist-R-h} and
\eqref{eq:rfembed} is precisely the embedded operator we have defined in
\eqref{eq:Rfembed}. In this way,
we may identify $\cD_\xi^\delta \cap \cS_\xi(T,C_0)$ as a subset of the
space $\cS_\xi(T,C_0)$ of Section \ref{sec:existenceuniqueness}, where the
discontinuity set $D \subset [0,T]$ of Definition \ref{def:S}
is $\delta\Z_+ \cap [0,T]$.

Then for any sufficiently large constants $C_0 \equiv C_0(T)>0$
and $\lambda \equiv \lambda(T,C_0)>0$,
the same argument as in \Cref{lem:solution-in-space} shows that 
if $X^\delta \in \cD_\xi^\delta \cap \cS_\xi(T,C_0)$, then 
$Y^\delta=\trsfrmDA(X^\delta)$ defined via
\eqref{eq:def-dist-theta}, \eqref{eq:def-dist-deriv-t}, and 
(\ref{eq:def-dist-C-t}--\ref{eq:def-dist-R-t}) satisfies
$Y^\delta \in \cD_\theta^\delta \cap \cS_\theta(T,C_0)$;
the same argument as in \Cref{lem:solution-in-space-2} shows that if
$Y^\delta \in \cD_\theta^\delta \cap \cS_\theta(T,C_0)$, then 
$X^\delta=\trsfrmDB(Y^\delta)$ defined via
\eqref{eq:def-dist-xi},
(\ref{eq:def-dist-deriv-xi-star}--\ref{eq:def-dist-resp-xi}), and
(\ref{eq:def-dist-C-h}--\ref{eq:def-dist-Gamma}) satisfy
$X^\delta \in \cD_\xi^\delta \cap \cS_\xi(T,C_0)$;
and the same arguments as in Lemmas \ref{lem:contract-Ct} and
\ref{lem:contract-fsys} show that the combined map
\[\trsfrmD \equiv \trsfrmDA \circ \trsfrmDB\]
is contractive in the metric $\lbddst{\cdot}{\cdot}$ on
$\cD_\theta^\delta \cap \cS_\theta(T,C_0)$. The conditions that
$(C_\theta,R_\theta) \in \cS_\theta(T,C_0)$ are $\delta$-p.c.r.c.\ and
hence belong to $\cD_\theta^\delta$ are closed under this metric, so
$\cD_\theta^\delta \cap \cS_\theta(T,C_0)$ is also complete 
in $\lbddst{\cdot}{\cdot}$. Then by the Banach fixed-point theorem,
there exists a fixed point of $\trsfrmD:\cD_\theta^\delta \cap \cS_\theta(T,C_0)
\to \cD_\theta^\delta \cap \cS_\theta(T,C_0)$,
which must coincide with $Y^\delta$
by the above claim of uniqueness of this fixed point in $\cD_\theta^\delta$.
This shows that $Y^\delta \in 
\cD_\theta^\delta \cap \cS_\theta(T,C_0)$, and hence also
$X^\delta \in \cD_\xi^\delta \cap \cS_\xi(T,C_0)$,
i.e.\ $(Y^\delta,X^\delta) \in \cS(T,C_0)$ as claimed.
\end{proof}

\begin{proof}[Proof of Lemma \ref{lem:dmft_discretization_error}(b)]
We compare the outputs of the continuous map $\trsfrmA$
and the discretized map $\trsfrmDA$ when both are
given the same discrete input $X^\delta \in \mathcal D_\xi^\delta \cap
\cS_\xi(T,C_0)$. Given any such input
$X^\delta=(\bar C_{f,\delta},\bar R_{f,\delta},\bar
R_{f,\delta}^*,\bar\Gamma_\delta)$, consider the continuous
process $\{\theta^t\}_{t \in [0,T]}$ and the embedded discrete process
$\{\bar\theta^t_\delta\}_{t \in [0,T]}$ coupled using the same underlying
realization of $\{\bar u_\delta^t\}_{t \in [0,T]} \sim \GP(0,\bar C_{f,\delta})$:
\begin{align*}
    \theta^t&=\theta^0-\int_0^t \bar\eta^r \Big({\gamma}\bar\Gamma^r_\delta
\theta^r+g(\theta^r)+{\gamma}\bar R_{f,\delta}^r(\theta^{[r]})+
{\gamma}\bar R_{f,\delta}^{r,*}\theta^*\Big)\d r+\sqrt{\gamma}\,\bar u^{t}_\delta\, , \\
    \bar\theta^t_\delta&=\theta^0-\int_0^{\flr{t}} \bar\eta^{\flr{r}} \Big({\gamma}\bar\Gamma^r_\delta
\bar\theta^r_\delta+g(\bar\theta^r_\delta)+{\gamma}\bar R_{f,\delta}^r(\bar\theta^{[r]}_\delta)+
{\gamma}\bar R_{f,\delta}^{r,*}\theta^*\Big)\d r+\sqrt{\gamma}\,\bar u^t_\delta \, .
\end{align*}
Then
\begin{align*}
    \bar\theta^t_\delta - \theta^t &= - \int_0^{\flr{t}} \bar\eta^r \prn{{\gamma}\bar\Gamma_\delta^{r}(\bar \theta^{r}_\delta - \theta^r)+ (g(\bar \theta^{r}_\delta)-g(\theta^r)) + {\gamma} \bar R_{f,\delta}^{r}\prn{\prn{\bar\theta_\delta - \theta}^{[r]}} } \d r \nonumber\\
    &\qquad+ \int_0^{\flr{t}} \prn{\bar\eta^r - \bar\eta^{\flr{r}}}\prn{g(\bar \theta^{r}_\delta) +  {\gamma}\prn{\bar\Gamma_\delta^{r}\bar \theta^{r}_\delta+ \bar R_{f,\delta}^{r}\prn{{\bar\theta_\delta}^{[r]}} + \bar R_{f,\delta}^{r,*}\theta^* }} \d r \nonumber\\
    &\qquad+ \int_{\flr{t}}^t \bar\eta^r \Big({\gamma}\bar\Gamma^r_\delta
\theta^r+g(\theta^r)+{\gamma}\bar R_{f,\delta}^r(\theta^{[r]})+
{\gamma}\bar R_{f,\delta}^{r,*}\theta^*\Big)\d r.
\end{align*}
As a consequence, $\E\brk{\norm{\bar\theta^t_\delta - \theta^t}^2} \leq
4((\mathrm{I})+(\mathrm{II})+(\mathrm{III})+(\mathrm{IV}))$ where
\begin{align*}
(\mathrm{I})&=\E\left\|\int_0^{\flr{t}}
\bar\eta^r \prn{{\gamma}\bar\Gamma_\delta^{r}(\bar \theta^{r}_\delta -
\theta^r)+ (g(\bar \theta^{r}_\delta)-g(\theta^r))}\d r\right\|^2,\\
(\mathrm{II})&=\E\left\|\int_0^{\flr{t}}
\bar\eta^r {{\gamma} \bar R_{f,\delta}^{r}\prn{\prn{\bar\theta_\delta -
\theta}^{[r]}} } \d r\right\|^2,\\
(\mathrm{III})&=\E\left\|\int_0^{\flr{t}}
\prn{\bar\eta^r - \bar\eta^{\flr{r}}}\prn{g(\bar \theta^{r}_\delta) +
{\gamma}\prn{\bar\Gamma_\delta^{r}\bar \theta^{r}_\delta+ \bar
R_{f,\delta}^{r}\prn{{\bar\theta_\delta}^{[r]}} + \bar
R_{f,\delta}^{r,*}\theta^* }} \d r \right\|^2,\\
(\mathrm{IV})&= \E\left\|\int_{\flr{t}}^t \bar\eta^r \Big({\gamma}\bar\Gamma^r_\delta
\theta^r+g(\theta^r)+{\gamma}\bar R_{f,\delta}^r(\theta^{[r]})+
{\gamma}\bar R_{f,\delta}^{r,*}\theta^*\Big)\d r\right\|^2.
\end{align*}
By the Lipschitz continuity of $g(\cdot)$ and bound \eqref{eq:pRf-volterra}
for $\bar R_{f,\delta}$, the first two terms can be bounded by
\[\nI+\nII \lesssim
\frac{1}{\lambda}e^{2\lambda t}\cdot\sup_{s\in[0,T]} e^{-2\lambda s}
\E\brk{\norm{\theta^s-\bar \theta^s_\delta}^2}\]
where $\lesssim$ denotes inequality up to a constant not depending on $\delta$.
For the third term, by the Lipschitz continuity of
$\{\bar \eta^t\}_{t \in [0,T]}$
in Assumption \ref{asp:SGD}, we have $\int_0^T (\bar \eta^r-\bar
\eta^{\flr{r}})^2 \d r \lesssim \delta^2$ for a constant $C>0$, and hence also
$\nIII \lesssim \delta^2$ by Cauchy-Schwarz. Finally $\nIV \lesssim \delta^2$
by the conditions for $X^\delta \in \cS_\xi(T,C_0)$. Then
choosing large enough $\lambda$ yields $\sup_{s\in[0,T]} e^{-2\lambda s}
\E\norm{\theta^s-\bar \theta^s_\delta}^2 \lesssim \delta^2$, which implies
$\lbddst{C_\theta}{\bar C_{\theta,\delta}} \lesssim \delta$.
A similar argument shows $\lbddst{R_\theta}{\bar R_{\theta,\delta}} \lesssim
\delta$, establishing
\begin{equation}\label{eq:trsfrmAclose}
\lbddst{\trsfrmA(X^\delta)}{\trsfrmDA(X^\delta)} \lesssim \delta
\text{ for all } X^\delta \in \cD_\xi^\delta \cap \cS_\xi(T,C_0).
\end{equation}

Analogously, we may compare $\trsfrmB$
and $\trsfrmDB$ given the same discrete input $Y^\delta=(\bar C_{\theta,\delta},\bar R_{\theta,\delta}) \in \cD_\theta^\delta \cap \cS_\theta(T,C_0)$. Coupling by the same realizations of $(\{\bar w_\delta^t\}_{t \in [0,T]},w^*_\delta) \sim \GP(0,\bar C_{\theta,\delta})$ an $\{z^t\}_{t \in [0,T]}$, we have
\begin{align*}
    \xi^t&={-}\int_0^{{t}} \frac{\bar\eta^{{r}}}{\bar\kappa} \bar R_{\theta,\delta}^{t,r}
f(\xi^{r-},w_\delta^*,\eps)\d z^r+\bar
w^t_\delta\, , \\
    \bar\xi^t_\delta&={-}\int_0^{\flr{t}} \frac{\bar\eta^{\flr{r}}}{\bar\kappa} \bar R_{\theta,\delta}^{t,r}
f(\bar\xi^{r-}_\delta,w_\delta^*,\eps)\d z^r+\bar
w^t_\delta \, .
\end{align*}
Then 
\begin{align*}
    \xi^t - \bar\xi^t_\delta &= - \underbrace{\int_0^{\flr{t}} \frac{\bar\eta^{{r}}}{\bar\kappa} \bar R_{\theta,\delta}^{t,r} \prn{f(\xi^{r-},w_\delta^*,\eps) - f(\bar\xi^{r-}_\delta,w_\delta^*,\eps)} \d z^r}_{(\mathrm{I})} \nonumber\\
    &\qquad- \underbrace{\int_0^{\flr{t}} \frac{\bar\eta^r - \bar\eta^{\flr{r}}}{\bar\kappa}R_{\theta,\delta}^{t,r} {f(\bar\xi^{r-}_\delta,w_\delta^*,\eps)} \d z^r}_{(\mathrm{II})} -  \underbrace{\int_{\flr{t}}^t \frac{\bar\eta^{{r}}}{\bar\kappa} \bar R_{\theta,\delta}^{t,r}
f(\xi^{r-},w_\delta^*,\eps)\d z^r}_{(\mathrm{III})}. 
\end{align*}
We can bound $\E[\norm{\xi^t - \bar\xi^t_\delta}^2] \leq 3(\E[(\mathrm{I})^2]+ \E[(\mathrm{II})^2] + \E[(\mathrm{III})^2])$.
For $\E[(\mathrm{I})^2]$, by using Lemma \ref{lem:burkholder}, boundedness of $\bar R_{\theta,\delta}$, and Lipschitz continuity of $f$, as in the above argument, we have
\[\E[(\mathrm{I})^2] \lesssim
\frac{1}{\lambda}e^{2\lambda t}\cdot\sup_{s\in[0,T]} e^{-2\lambda s}
\E\brk{\norm{\xi^s-\bar \xi^s_\delta}^2}.\]
Using Lemma \ref{lem:burkholder}, we have similarly to the proof of \eqref{eq:trsfrmAclose} that $\E[(\mathrm{II})^2] \lesssim \delta^2 $ and $\E[(\mathrm{III})^2] \lesssim \delta $ up to some constants not depending on $\delta$. Then choosing $\lambda$ large enough and rearranging shows $\sup_{s\in[0,T]} e^{-2\lambda s}
\E\norm{\xi^s-\bar \xi^s_\delta}^2 \lesssim \delta$. Then by an analogous argument,
\[
\lbddst{C_f}{\bar C_{f,\delta}}^2 \lesssim \sup_{t\in [0,T]} e^{-2\lambda t}\E \left\|\int_0^t \frac{\bar\eta^r}{\bar \kappa}f(\xi^{r-},w^*_\delta,\eps) \d z^r- \int_0^{\flr{t}} \frac{\bar\eta^{\flr{r}}}{\bar \kappa}f(\bar\xi^{r-}_\delta,w^*_\delta,\eps)\d z^r\right\|^2 \lesssim \delta\,,
\]
that is,
$\lbddst{C_f}{\bar C_{f,\delta}} \lesssim \delta^{1/2}$.
Similarly, one may analyze $R_f$ and $R_f^*$ to show that
\begin{equation}\label{eq:trsfrmBclose}
\lbddst{\trsfrmB(Y^\delta)}{\trsfrmDB(Y^\delta)} \lesssim \delta^{1/2}
\text{ for all } Y^\delta \in \cD_\theta^\delta \cap \cS_\theta(T,C_0),
\end{equation}
and we omit this verification for brevity.

Now let $Y^\delta \in \cD_\theta^\delta \cap \cS_\theta(T,C_0)$ be the fixed
point of $\trsfrmD$, and let $Y \in \cS_\theta^\text{cont}(T,C_0)$ be the fixed
point of $\trsfrm$. Then
\begin{align*}
	\lbddst{Y}{Y^\delta} & =
\lbddst{\mathcal{T}(Y)}{\mathcal{T}^\delta(Y^\delta)}  \leq
\underbrace{\lbddst{\mathcal{T}(Y)}{\mathcal{T}(Y^\delta)}}_{\mathrm{(I)}} +
\underbrace{\lbddst{\mathcal{T}(Y^\delta)}{
\mathcal{T}^\delta(Y^\delta)}}_{\mathrm{(II)}}.
\end{align*}
By the contractivity \eqref{eq:T-contraction} of $\trsfrm$, we can choose
$\lambda$ large enough such that
\begin{align}
	\mathrm{(I)} \leq \frac 1 2 \lbddst{Y}{Y^\delta} \, . \label{eq:mid-integral-differential-approximation-2}
\end{align}
To control $\mathrm{(II)}$, we have
\begin{align*}
\nII&=\lbddst{\trsfrmA \circ \trsfrmB(Y^\delta)}{\trsfrmDA \circ
\trsfrmDB(Y^\delta)}\\
&\leq \lbddst{\trsfrmA \circ \trsfrmB(Y^\delta)}{\trsfrmA \circ
\trsfrmDB(Y^\delta)}\\
&\hspace{1in}+\lbddst{\trsfrmA \circ \trsfrmDB(Y^\delta)}{\trsfrmDA \circ
\trsfrmDB(Y^\delta)} \lesssim \delta^{1/2}
\end{align*}
where the last inequality holds
by \eqref{eq:trsfrmAclose}, \eqref{eq:trsfrmBclose}, and Lemma
\ref{lem:contract-Ct}. Thus $\lbddst{Y}{Y^\delta} \lesssim \delta^{1/2}$.
Setting $X=\trsfrmB(Y)$ and $X^\delta=\trsfrmDB(Y^\delta)$, this
implies by Lemma \ref{lem:contract-fsys} that also
$\lbddst{X}{X^\delta} \lesssim \delta^{1/2}$, so
\[\lim_{\delta \to 0} \lbddst{X}{X^\delta}=0,
\qquad \lim_{\delta \to 0} \lbddst{Y}{Y^\delta}=0.\]
\end{proof}

\begin{proof}[Proof of Lemma \ref{lem:dmft_discretization_error}(c)]
The convergence $\lim_{\delta\to 0}\lbddst{C_f}{\bar C_{f,\delta}}=0$
implies that there exists a coupling of $\{u^t\}_{t \in [0,T]} \sim
\GP(0,C_f)$ and $\{\bar u^t_\delta\}_{t \in [0,T]} \sim \GP(0,\bar C_{f,\delta})$ for which
\[\lim_{\delta \to 0} \sup_{t \in [0,T]} \E\norm{u^t-\bar u^t_\delta}^4=0.\]
Then defining $\{\theta^t\}_{t \in [0,T]}$ via $(C_f,R_f,R_f^*,\Gamma)$
and $\{\bar \theta_\delta^t\}_{t \in [0,T]}$ via
$(\bar C_{f,\delta},\bar R_{f,\delta},\bar R_{f,\delta}^*,\bar
\Gamma_\delta)$ using this coupling of $\{u^t\}$ and $\{\bar u_\delta^t\}$, the same arguments as leading to (\ref{eq:theta4diffbound}) show
\[\lim_{\delta \to 0} \sup_{t \in [0,T]} e^{-4\lambda t}
\E\norm{\theta^t-\bar \theta^t_\delta}^4=0.\]
In particular,
$\lim_{\delta \to 0} \E[\norm{\theta^{t_1}-\bar
\theta^{t_1}_\delta}^2+\ldots+\norm{\theta^{t_m}-\bar \theta^{t_m}_\delta}^2]
=0$, which implies the Wasserstein-2 convergence
(\ref{eq:discretedmftthetaconv}). A similar argument shows
(\ref{eq:discretedmftetaconv}).
\end{proof}

\section{Discretization Error of SGD and SME}\label{sec:disc-sgd-sme}

In this section, we conclude the proof of Theorem \ref{thm:DMFT-limit} by
analyzing the original SGD/SME dynamics \eqref{eq:SGDrescaled} and
\eqref{eq:SME}, and
showing that they are also well-approximated by the discrete-time dynamics
\eqref{eq:discretedynamics} as $\delta \to 0$.

\subsection{Discretization of SGD}

Recall the SGD update
\[\bar\btheta^{k+1}-\bar\btheta^k=-\eta^k {\left(\frac{1}{\kappa}
\sum_{i \in S^k} \x_i \otimes
f(\x_i^\top \bar\btheta^k,\x_i^\top \btheta^*,\eps_i)
+\frac{1}{n}\,g(\bar\btheta^k)\right)}, \qquad |S^k|=\kappa,\]
where $g(\cdot)$ is applied row-wise. We will analyze a Poissonized
version of this process: Recall $\set=\{S \subset [n]:|S|=\kappa\}$, and let
\[\bN^t=\{N_S^t\}_{S \in \set}\]
be a $\R^{\binom{n}{\kappa}}$-valued process indexed by $\set$, whose
coordinates are independent homogeneous Poisson jump processes with rate
$n^{1-\alpha}/{n \choose \kappa}$. Set
\[N^t=\1^\top \bN^t=\sum_{S \in \set} N_S^t.\]
Marginally, $N^t$ is a Poisson jump process with rate $n^{1-\alpha}$. We define
a Poissonized version of the SGD dynamics $\{\bar\btheta^k\}_{k \geq 0}$ of
\eqref{eq:SGD} (without time rescaling) by
\begin{equation}\label{eq:SGDPoissonized}
\bnu^t=\bar\btheta^{N^t}, \qquad \bnu^0=\btheta^0.
\end{equation}
Then $\bnu^t$ is a continuous-time jump process given by
\begin{align*}
\bnu^t&=\bnu^0-\int_0^t \eta^{N^{s-}} \sum_{S \in \set}\prn{\frac{1}{\kappa}
\sum_{i \in S} \x_i \otimes f(\x_i^\top \bnu^{s-},\x_i^\top \btheta^*,\eps_i)
+ \frac{1}{n} g(\bnu^{s-})} \d N^s_S\\
&=\bnu^0-\int_0^t \eta^{N^{s-}} \prn{\frac{1}{\kappa}\sum_{i=1}^n
\x_i \otimes f(\x_i^\top \bnu^{s-},\x_i^\top \btheta^*,\eps_i)\sum_{S \in
\set}
\1_{i\in S} \d N^s_S+\frac{1}{n} g(\bnu^{s-}) \sum_{S \in \set}
\d N^s_S}\\
&=\bnu^0-\int_0^t \eta^{N^{s-}} \prn{\frac{1}{\kappa}\,\X^\top
\diag(\bPi\,\d\bN^s)f(\X\bnu^{s-},\X\btheta^*,\beps))
+\frac{1}{n} g(\bnu^{s-}) \d N^s}.
\end{align*} 
Here $f(\cdot)$ is also applied row-wise, and $\bPi\in\R^{n\times{n\choose
\kappa}}$ is the incidence matrix with entries $\Pi_{i,S}=\1_{i \in S}$.

Fixing a discretization step size $\delta>0$, in this section, we continue
to write $\flr{t}$ and $\cil{t}$ for the time discretization notation
previously defined in \eqref{eq:flrcil}. We then consider an approximation
$\{\bnu_\delta^t\}_{t \geq 0}$ to the above Poissonized process, defined by
\begin{equation}\label{eq:deltaSGD}
\bnu_\delta^t = \bnu^0 - \int_0^{\flr{t}}
\frac{\bar\eta^{\flr{s}}}{\bar\kappa} \X^\top \diag(\bPi\,\d\bN^s)
f(\X\bnu_\delta^{s-},\X\btheta^*,\beps)
-\int_0^{\flr{t}} \bar \eta^{\flr{s}} g(\bnu_\delta^{s}) \d s.
\end{equation}
It may be checked that $\{\bnu_\delta^t\}_{t \geq 0}$ is a
piecewise-constant embedding of the discrete-time dynamics
\eqref{eq:discretedynamics} in the SGD setting of
$\z_\delta^k=\bPi\bP_\delta^k$,
\[\bnu_\delta^t=\btheta_\delta^k \text{ if } \flr{t}=k,\]
upon identifying $\bP_\delta^k=\int_{k\delta}^{(k+1)\delta} \d\bN^s
=\bN^{(k+1)\delta}-\bN^{k\delta}$.

The main result of this section is the following comparison lemma.

\begin{lem}\label{lem:disc-sgd}
For any fixed $T>0$, $\delta \in (0,1)$,
and some constant $C \equiv C(T,\alpha)>0$ not depending on $\delta$,
almost surely as $n,d \to \infty$,
\[\limsup_{n,d\to
\infty}\sup_{t\in[0,T]}\frac{\norm{\bnu^t-\bnu_\delta^t}^2}{d} \leq
C\delta.\]
\end{lem}

We proceed to prove Lemma \ref{lem:disc-sgd}.
The following first establishes a high-probability bound on
the norm of $\{\bnu^t\}_{t \in [0,T]}$, which
will be important for controlling error terms in the subsequent analysis.

\begin{lem}\label{lem:l2-bound-sgd}
For any fixed $T>0$ and some constant $C \equiv C(T,\alpha)>0$,
almost surely for all large $n,d$,
\[\sup_{t\in [0,T]}\norm{\bnu^t} \leq C\sqrt{n}.\]
\end{lem}
\begin{proof}
We continue to write $\|\cdot\|$ for the vector $\ell_2$-norm and matrix
Frobenius norm, denote $\|\M\|_\infty=\max_{i,j} |M_{ij}|$, and write $\lesssim$
for inequality up to a constant possibly depending on $T,\alpha$.

Using the boundedness of $f(\cdot)$ and Lipschitz continuity of $g(\cdot)$ in
Assumption \ref{asp:lipschitz} and the scalings $\eta^t \lesssim n^\alpha$
and $\kappa \asymp n^\alpha$ in Assumption \ref{asp:SGD}, we have
\begin{align*}
\|\bnu^t\| &\lesssim \|\bnu^0\|+\|\X\|_\op\left\|\int_0^t \diag(\bPi\,\d\bN^s)
f(\X\bnu_\delta^{s-},\X\btheta^*,\beps)\right\|
+n^{\alpha-1}\left\|\int_0^t g(\bnu^{s-})\1^\top \d\bN^s\right\|\\
&\lesssim \|\bnu^0\|+\|\X\|_\op\left\|\int_0^t \bPi\,\d\bN^s\right\|
\|f(\X\bnu_\delta^{s-},\X\btheta^*,\beps)\|_\infty
+n^{\alpha-1}\int_0^t \|g(\bnu^{s-})\|\1^\top \d\bN^s\\
&\lesssim \|\bnu^0\|+\|\X\|_\op\|\bPi\,\bN^T\|
+n^{\alpha-1}\int_0^t (\|\bnu^{s-}\|+\sqrt{n}) \d N^s.
\end{align*}
Here, $\bN^T$ has i.i.d.\ coordinates with law
$\Pois(Tn^{1-\alpha}/\binom{n}{\kappa})$. Then
applying the concentration inequality \eqref{eq:poissonconcentration} to
$Q=\|\bPi\,\bN^T\|^2$, where
\[\E Q=\left(\frac{Tn^{1-\alpha}}{\binom{n}{\kappa}}\right)^2
\|\bPi \1\|^2
+\frac{Tn^{1-\alpha}}{\binom{n}{\kappa}}
\Tr \bPi^\top \bPi
=\left(\frac{Tn^{1-\alpha}}{\binom{n}{\kappa}}\right)^2
\cdot n\binom{n-1}{\kappa-1}^2
+\frac{Tn^{1-\alpha}}{\binom{n}{\kappa}} \binom{n}{\kappa}\kappa
\lesssim n,\]
$\mu_0 \lesssim \E Q \lesssim n$, and $\mu_2 \lesssim \|\bPi^\top \bPi\|_\infty
\leq \kappa \lesssim n^\alpha$, we have $Q \lesssim n$ a.s.\ for all large
$n,d$. Since $N^T=\1^\top \bN^T \sim \Pois(n^{1-\alpha})$ where
$\alpha<1$, by a standard Poisson tail bound, also
$N^T \lesssim n^{1-\alpha}$ a.s.\ for all large $n,d$.
Under Assumption \ref{asp:model}, 
$\|\X\|_\op \lesssim 1$ a.s.\ for all large $n,d$, and
$\|\bnu^0\|=\|\btheta^0\| \lesssim \sqrt{n}$. Then on the
intersection $\mathcal{E}$ of these events, which holds a.s.\ for all large $n,d$,
\[\|\bnu^t\| \lesssim \sqrt{n}+n^{\alpha-1}\int_0^t \|\bnu^{s-}\| \d N^s.\]
Letting $t_1,t_2,\ldots$ be the jumps of $\{N^t\}_{t \in [0,T]}$, this means
that for a constant $C \equiv C(T,\alpha)>0$,
\[\|\bnu^{t_{k+1}}\| \leq
C\sqrt{n}+Cn^{\alpha-1}(\|\bnu^{t_1}\|+\ldots+\|\bnu^{t_k}\|).\]
A discrete Gr\"onwall inequality then shows that
$\|\bnu^{t_k}\| \leq C\sqrt{n} e^{Ckn^{\alpha-1}}$ for all $k=1,\ldots,N^T$.
Since the number of jumps $N^T$
satisfies $N^T \lesssim n^{1-\alpha}$ on $\mathcal{E}$, this shows the lemma.
\end{proof}

\begin{proof}[Proof of Lemma \ref{lem:disc-sgd}]
Let
\[\bZ^t=\bN^t-t\frac{n^{1-\alpha}}{{n\choose \kappa}}\mathbf{1}_{n\choose
\kappa}, \qquad Z^t=\1^\top\bZ^t\]
be the centered versions of $\bN^t$ and $N^t=\1^\top \bN^t$, where $\1_{n \choose \kappa}$ denotes the
all-1's vector in $\R^{n \choose \kappa}$. Then
\begin{equation}\label{eq:PiN}
\bPi\,\d\bN^s
=\bPi\,\d\bZ^s+n^{1-\alpha}\frac{\binom{n-1}{\kappa-1}}{\binom{n}{\kappa}}
\1_{n \choose \kappa} \d s
=\bPi\,\d\bZ^s+\kappa n^{-\alpha}\1_{n \choose \kappa} \d s
\end{equation}
where we have used
$\bPi\cdot \mathbf{1}_{n\choose \kappa} = {n-1 \choose \kappa-1}
\mathbf{1}_{n\choose \kappa}$. We introduce an intermediary process
\begin{equation}\label{eq:tildenu}
\tilde\bnu^t = \bnu^0 - \int_0^t \frac{\bar\eta^{s}}{\bar\kappa} \X^\top 
\diag(\bPi\,\d\bN^s)f(\X\tilde\bnu^{s-},\X\btheta^*,\beps) - \int_0^t
{\bar\eta^{s}} g(\tilde\bnu^{s}) \d s
\end{equation}
and proceed to first compare $\{\bnu^t\}_{t \in [0,T]}$ to $\{\tilde\bnu^t\}_{t
\in [0,T]}$, and then compare $\{\tilde\bnu^t\}_{t \in [0,T]}$ to the final
discretized process $\{\bnu_\delta^t\}_{t \in [0,T]}$.\\

\paragraph{Step 1: Bounding $\norm{\bnu^t - \tilde \bnu^t}$}

We decompose the difference into four terms
\begin{align*}
\bnu^t - \tilde \bnu^t&=\underbrace{\int_0^t
\frac{\bar\eta^{s}}{\bar\kappa} \X^\top
\diag(\bPi\,\d\bN^s)f(\X\tilde\bnu^{s-},\X\btheta^*,\beps) - \int_0^t
\frac{\eta^{N^{s-}}}{\kappa} \X^\top
\diag(\bPi\,\d\bN^s)f(\X\bnu^{s-},\X\btheta^*,\beps)}_{\mathrm{(I)} + \mathrm{(II)}}\\
    &\qquad+\underbrace{\int_0^t \bar\eta^{s} g(\tilde\bnu^{s}) \d s - \int_0^t
\eta^{N^{s-}} \frac{1}{n} g(\bnu^{s-}) \d N^s}_{\mathrm{(III)} + \mathrm{(IV)}}\,,
\end{align*}
applying \eqref{eq:PiN} and $\d N^s=\d Z^s+n^{1-\alpha}\,\d s$
to separate $\nI+\nII$ and $\nIII+\nIV$ into their drift and martingale
components:
\begin{align*}
    \mathrm{(I)} &= \int_0^t \frac{\bar\eta^{s}}{\bar\kappa} \X^\top
(f(\X\tilde\bnu^s,\X\btheta^*,\beps)-f(\X\bnu^s,\X\btheta^*,\beps)) \kappa
n^{-\alpha} \d s\\
&\hspace{1in} - \int_0^t \prn{\frac{\eta^{N^s}}{\kappa} -
\frac{\bar\eta^{s}}{\bar\kappa}} \X^\top f(\X\bnu^s,\X\btheta^*,\beps) \kappa
n^{-\alpha}\d s\, ,\\
    \mathrm{(II)} &= \underbrace{\int_0^t \frac{\bar\eta^{s}}{\bar\kappa} \X^\top
\diag(\bPi\,\d\bZ^s)(f(\X\tilde\bnu^{s-},\X\btheta^*,\beps)-f(\X\bnu^{s-},\X\btheta^*,\beps))}_{\mathrm{(II.a)}}\\
&\hspace{1in} - \underbrace{\int_0^t \prn{\frac{\eta^{N^{s-}}}{\kappa} -
\frac{\bar\eta^{s}}{\bar\kappa}} \X^\top
\diag(\bPi\,\d\bZ^s)f(\X\bnu^{s-},\X\btheta^*,\beps)}_{\mathrm{(II.b)}},\\
    \mathrm{(III)} & = \int_0^t \bar\eta^{s} g(\tilde\bnu^{s}) \d s - \int_0^t
{\frac{\eta^{N^{s-}}}{n^\alpha} g(\bnu^{s-})  } \d s\, ,\\
\mathrm{(IV)} &= - \int_0^t \eta^{N^{s-}} \frac{1}{n} g(\bnu^{s-})\d Z^s.
\end{align*}

By Doob's $L^p$-maximal inequality and a standard bound for the central moments
of a Poisson random variable, for any fixed $p>1$,
$\E \sup_{t \in [0,T]} |\frac{N^t}{n^{1-\alpha}}-t|^p
\lesssim \E |\frac{N^T}{n^{1-\alpha}}-T|^p \lesssim n^{-(1-\alpha)p/2}$. Thus,
fixing any sufficiently small constant $\iota>0$,
Markov's inequality and the Borel-Cantelli lemma imply
that a.s.\ for all large $n,d$,
\begin{equation}\label{eq:Poissontimeapprox}
\sup_{t \in [0,T]} \left|\frac{Z^t}{n^{1-\alpha}}\right|
=\sup_{t \in [0,T]} \left|\frac{N^t}{n^{1-\alpha}}-t\right| \leq
n^{-\frac{1-\alpha}{2}+\iota}<n^{-\iota}.
\end{equation}
By the scaling and continuity assumptions for
$\kappa,\eta$ in Assumption \ref{asp:SGD}, we then
have for any fixed $p \geq 1$ that
\begin{equation}\label{eq:Poissonprocapprox}
\int_0^T \abs{\frac{\eta^{N^s} }{\kappa} -  \frac{\bar\eta^{s}}{\bar\kappa}}^p
\d s=\int_0^T \abs{\frac{\bar\eta^{N^s/n^{1-\alpha}}}{\bar \kappa} -
\frac{\bar\eta^{s}}{\bar\kappa}}^p \d s+o(1)=o(1).
\end{equation}
Thus a.s.\ for all large $n,d$, by the boundedness and Lipschitz continuity of $f$,
\begin{align}
    \norm{\mathrm{(I)}} &\lesssim  \int_0^t \norm{\X}^2 \norm{\tilde\bnu^{s}
- \bnu^{s}} \d s + \norm{\X} \sqrt{n} \int_0^t \abs{\frac{\eta^{N^s} }{\kappa}
-  \frac{\bar\eta^{s}}{\bar\kappa}} \d s\notag\\
    &\lesssim \int_0^t \norm{\tilde\bnu^{s} - \bnu^{s}}  \d s
+o(\sqrt{n}).\label{eq:SGDIbound}
\end{align}
Since $g$ is Lipschitz and $\sup_{t \in [0,T]} \norm{\bnu^t} \lesssim \sqrt{n}$
a.s.\ for all large $n,d$ from \Cref{lem:l2-bound-sgd},
we can similarly bound
\begin{equation}\label{eq:SGDIIIbound}
\norm{\nIII} \lesssim \int_0^t \norm{\tilde\bnu^{s} - \bnu^{s}}  \d s +
o(\sqrt{n}).
\end{equation}
Also $\sup_{t \in [0,T]} \eta^t \lesssim n^\alpha$ by Assumption \ref{asp:SGD},
and $\sup_{t \in [0,T]} \norm{g(\bnu^t)} \lesssim \sqrt{n}$ a.s.\ for all
large $n,d$ by \Cref{lem:l2-bound-sgd} and the Lipschitz continuity of $g$, so
\begin{equation}\label{eq:SGDIVbound}
\norm{\nIV} \lesssim n^{\alpha-\frac{1}{2}} \int_0^t \d Z^s
\lesssim n^{\alpha-\frac{1}{2}} \sup_{t \in [0,T]} |Z^t| \lesssim 
n^{\alpha-\frac{1}{2}} \cdot n^{\iota+\frac{1-\alpha}{2}}
\lesssim n^{\frac{\alpha}{2}+\iota}=o(\sqrt{n}),
\end{equation}
the last bound holding for $\iota>0$ small enough.

For the term $\nII$, let $\{\bE^t\}_{t \geq 0}$ be any $\R^{m \times
\binom{n}{\kappa}}$-valued predictable process, and define
\[\M^t=\int_0^t \bE^s \d \bZ^s.\]
By It\^o's formula, for any continuously-differentiable function $F:\R^m \to
\R$,
\begin{equation}\label{eq:Ito-jump}
F(\M^t)-F(\mathbf{0})=\int_0^t \nabla F(\M^{s-})^\top \d \M^s + \sum_{0<s\leq t}
(F(\M^s)-F(\M^{s-})-\nabla F(\M^{s-})^\top \Delta \M^s)
\end{equation}
where $\Delta\M^s=\M^s-\M^{s-}$. In particular,
\[
\norm{\M^t}^2 = 2\int_0^t {\M^{s-}}^{\top}\bE^s\d\bZ^s + \sum_{0<s\leq
t}\norm{\Delta \M^s}^2\, .
\]
Here, each jump of $\M^s$ occurs at a jump of $N_S$ for some $S \in \set$, in
which case $\Delta \M^s = \bE^s \e_S$ where $\e_S$ is the standard basis
vector corresponding to the coordinate $S \in \set$. Hence
\begin{align*}
\sum_{0<s\leq t}\norm{\Delta \M^s}^2 &= \sum_{S \in \set} \int_0^t
\norm{\bE^s \e_S}^2 \d {N}^s_S = \sum_{S \in \set} \int_0^t \norm{\bE^s
\e_S}^2 \d Z^s_S + \frac{n^{1-\alpha}}{{n\choose \kappa}}\sum_{S \in \set}
\int_0^t \norm{\bE^s \e_S}^2 \d s\nonumber\\
&=\int_0^t \Tr \bE^s \diag(\d \bZ^s){\bE^s}^\top + \frac{n^{1-\alpha}}{{n\choose
\kappa}}\int_0^t \Tr \bE^s{\bE^s}^\top \d s\,.
\end{align*}
Thus
\[\|\M^t\|^2=\underbrace{\frac{n^{1-\alpha}}{{n\choose
\kappa}}\int_0^t \Tr \bE^s{\bE^s}^\top \d s}_{:=P^t}
+\underbrace{2\int_0^t {\M^{s-}}^{\top}\bE^s\d\bZ^s 
+\int_0^t \Tr \bE^s \diag(\d \bZ^s){\bE^s}^\top}_{:=Q^t}.\]

We apply this to the first component $\mathrm{(II.a)}$ of $\mathrm{(II)}$,
identifying $\R^{d \times k} \equiv \R^m$ with $m=dk$ and noting that
\[\mathrm{(II.a)}=\int_0^t \frac{\bar\eta^{s}}{\bar\kappa} \X^\top
\diag(\bPi\,\d\bZ^s)(f(\X\tilde\bnu^{s-},\X\btheta^*,\beps)-f(\X\bnu^{s-},\X\btheta^*,\beps))=\int_0^t
\bE^s\d\bZ^s\]
where
\[\bE^s=\frac{\bar\eta^{s}}{\bar\kappa} \A^s\bPi,
\qquad \A^s=[\a_1^s,\ldots,\a_n^s] \in \R^{dk \times n},\]
\[\a_i^s=\vec(\x_i \otimes [f(\x_i^\top \tilde
\bnu^{s-},\X\btheta^*,\beps) - f(\x_i^\top \bnu^{s-},\X\btheta^*,\beps)]).\]
This gives the decomposition
\begin{equation}\label{eq:IIabound}
\|\mathrm{(II.a)}\|^2=P^t+Q^t.
\end{equation}
For this choice of $\bE^s$, we may bound
$P^t$ using $\bPi \bPi^\top = ({n-1 \choose \kappa-1} - {n-2 \choose
\kappa-2}) \Id_n + {n-2 \choose \kappa-2}
\mathbf{1}_n\mathbf{1}_n^\top$,
where ${n-1 \choose \kappa-1} - {n-2 \choose
\kappa-2}=(1-\frac{\kappa-1}{n-1})\binom{n-1}{\kappa-1} \lesssim
\frac{\kappa}{n}\binom{n}{\kappa} \lesssim n^{\alpha-1}\binom{n}{\kappa}$ and
$\binom{n-2}{\kappa-2} \lesssim \frac{\kappa^2}{n^2}\binom{n}{\kappa}
\lesssim n^{2\alpha-2}\binom{n}{\kappa}$. Then
\begin{align*}
\Tr \bE^s{\bE^s}^\top & \leq \prn{\frac{\bar\eta^{t}}{\bar\kappa}}^2
\left(\binom{n-1}{\kappa-1}-\binom{n-2}{\kappa-2}\right)
\Tr \A^s{\A^s}^\top
+\prn{\frac{\bar\eta^{t}}{\bar\kappa}}^2
\binom{n-2}{\kappa-2}\|\A^s\1_n\|^2\\
&\lesssim n^{\alpha-1}\binom{n}{\kappa}
\sum_{i=1}^n \|\x_i\|^2\|f(\x_i^\top \tilde
\bnu^{s-},\X\btheta^*,\beps) - f(\x_i^\top \bnu^{s-},\X\btheta^*,\beps)\|^2\\
&\hspace{1in}+n^{2\alpha-2}\binom{n}{\kappa}\|\X^\top(f(\X\tilde
\bnu^{s-},\X\btheta^*,\beps) - f(\X\bnu^{s-},\X\btheta^*,\beps))\|^2\\
&\lesssim n^{\alpha-1}\binom{n}{\kappa}\|f(\X\tilde
\bnu^{s-},\X\btheta^*,\beps) - f(\X\bnu^{s-},\X\btheta^*,\beps)\|^2
\lesssim n^{\alpha-1}\binom{n}{\kappa}\|\tilde \bnu^{s-}-\bnu^{s-}\|^2
\end{align*}
where the last two inequalities hold on a $\X$-dependent event a.s.\ for all
large $n$. So
\begin{equation}\label{eq:Ptbound}
P_t \lesssim \int_0^t \norm{\tilde\bnu^s - \bnu^s}^2\d s.
\end{equation}

For the martingale part $Q^t$, we may write
\[Q^t=\int_0^t {\H^s}^\top \d\bZ^s,
\quad \H^s=(H_S^s)_{S \in \set},
\quad H_S^s=2{\M^{s-}}^\top \bE^s\e_S+\|\bE^s\e_S\|^2.\]
Then by the Burkholder-Davis-Gundy type inequality of
\cite[Theorem 2.11]{kunita2004stochastic}, 
identifying $\bZ^s$ as a compensated Poisson random measure on $(0,T] \times
\set$, for any $p \geq 1$ and some constant $C_p>0$,
\begin{equation}\label{eq:Qtmomentbound}
\E \sup_{t \in [0,T]} (Q^t)^{2p}
\leq C_p\left(\E\left(\int_0^T \|\H^s\|^2\,
\frac{n^{1-\alpha}}{\binom{n}{\kappa}}\d s\right)^p
+\E\int_0^T \sum_{S \in \set} |H_S^s|^{2p}\,
\frac{n^{1-\alpha}}{\binom{n}{\kappa}}\d s\right).
\end{equation}
To bound $\H^s$, note that
\begin{align*}
\norm{\bE^s\e_S} &\lesssim \|\A^s\|_\op \cdot \|\bPi\e_S\|
\lesssim \normop{\X}\norm{f}_\infty \cdot \sqrt{\kappa}  \lesssim n^{\alpha/2},
\end{align*}
and the operator norm of this matrix can be bounded similarly by 
\begin{align*}
    \norm{\bE^s}_\op &\lesssim \normop{\X}\norm{f}_\infty \cdot \norm{\bPi}_\op \lesssim
\sqrt{\|\bPi\|_{\ell_1 \to \ell_1} \|\bPi\|_{\ell_\infty \to \ell_\infty}}
\lesssim \sqrt{\kappa \binom{n-1}{\kappa-1}}
\lesssim
{\kappa}\sqrt{\frac{{n\choose \kappa}}{n}} \lesssim
n^{\alpha}\sqrt{\frac{{n\choose \kappa}}{n}}.
\end{align*}
Thus $|H_S^s| \lesssim \norm{\M^{s-}}\,n^{\alpha/2} + n^{\alpha}$, and
\begin{align*}
\|\H^s\|^2 &\lesssim \sum_{S \in \set} \brk{\prn{{\M^{s-}}^\top \bE^s \e_S}^2
+ n^{2\alpha}} \lesssim \prn{{\M^{s-}}^\top \bE^s{\bE^s}^\top \M^{s-}
+ n^{2\alpha}{n\choose \kappa}}\\
&\lesssim \prn{\norm{\M^{s-}}^2 n^{2\alpha-1}{n\choose \kappa} + n^{2\alpha}{n\choose \kappa}}\,.
\end{align*}
As a consequence, applying these bounds and H\"older's inequality to
\eqref{eq:Qtmomentbound},
\begin{align}
    \E\sup_{t\in[0,T]} (Q^t)^{2p}
    &\lesssim \E\prn{n^{\alpha}\int_0^T \norm{\M^{s-}}^2 \d s +
n^{\alpha+1}}^{p} + \frac{n^{1-\alpha}}{{n\choose\kappa}}\E\int_0^T \sum_{S \in
\set} \prn{\norm{\M^{s-}}n^{\alpha /2} + n^{\alpha }}^{2p} \d s\nonumber\\
    &\lesssim n^{\alpha p}\,\E\int_0^T \norm{\M^s}^{2p} \d s +
n^{p(\alpha+1)} + n^{1-\alpha}\left(n^{\alpha p}\,\E \int_0^T
\norm{\M^s}^{2p} \d s+n^{2\alpha p}\right)\nonumber\\
    &\lesssim n^{\alpha p + 1 - \alpha}\int_0^T \E \norm{\M^s}^{2p} \d s +
n^{\alpha p + p}\label{eq:kunita-BDG}
\end{align}
for $p \geq 1$. To bound $\E \|\M^s\|^{2p}$, observe that taking expectations
in It\^o's formula \eqref{eq:Ito-jump} and recalling that each jump of
$\M^s$ is given by $\nabla \M^s=\bE^s\e_S$, corresponding to the jump of
some coordinate $N_S$ which has rate $\frac{n^{1-\alpha}}{\binom{n}{\kappa}}$,
\begin{align}
    \frac{\d}{\d t} \E F(\M^{t}) &=
\frac{n^{1-\alpha}}{{n\choose\kappa}}\E\sum_{S \in \set}
\prn{F(\M^{t-}+\bE^t\e_S) - F(\M^{t-}) - \nabla F(\M^{t-})^\top \bE^t
\e_S} \nonumber\\
    &= \frac{n^{1-\alpha}}{{n\choose\kappa}}\E\sum_{S \in \set} \int_0^1
(1-h)\,\e_S^\top {\bE^t}^\top \cdot \nabla^2 F(\M^{t-}+h\bE^t \e_S) \cdot \bE^t \e_S \d h\nonumber\\
    &\lesssim \frac{n}{{n\choose\kappa}}\E\sum_{S \in \set} \int_0^1
(1-h)\norm{\nabla^2 F(\M^{t-}+h\bE^t \e_S)}\d h\, .
\end{align}
Applying this with $F(\x)=\norm{\x}^{2p}$, we know $\norm{\nabla^2 F(\x)}
\lesssim \norm{\x}^{2p-2}$. Thus
\begin{align*}
    \frac{\d}{\d t} \E \norm{\M^{t}}^{2p} &\lesssim
\frac{n}{{n\choose\kappa}}\E\sum_{S \in \set} \int_0^1
(1-h)\norm{\M^{t-}+h\bE^t \e_S}^{2p-2}\d h
    \lesssim \frac{n}{{n\choose\kappa}}\E\sum_{S \in \set}
\prn{\norm{\M^{t-}}^{2p-2}+\norm{\bE^t \e_S}^{2p-2}} \nonumber\\
    &\lesssim n  \prn{\E \norm{\M^{t-}}^{2p-2}+n^{\alpha (p-1)}}
\lesssim n (\E \norm{\M^{t-}}^{2p})^{\frac{p-1}{p}}+n^{\alpha(p-1)+1}
   \lesssim \E\norm{\M^{t-}}^{2p}+n^p.
\end{align*}
Equivalently, in integral form,
$\E \norm{\M^{t}}^{2p} \lesssim \int_0^t \prn{\E\norm{\M^s}^{2p}+n^p}\d s$.
Then by Gr\"onwall's inequality, $\E\norm{\M^{t}}^{2p} \lesssim n^p$.
Applying this to \eqref{eq:kunita-BDG} shows $\E \sup_{t\in[0,T]} (Q^t)^{2p} \lesssim
n^{(\alpha + 1)p + (1 - \alpha)}$. For any fixed and sufficiently small $\iota>0$, choosing $p \geq 1$
large enough and applying Markov's inequality and the Borel-Cantelli lemma, this
implies that a.s.\ for all large $n,d$,
\begin{equation}\label{eq:Qtbound}
    \sup_{t\in[0,T]} |Q^t|<n^{\frac{\alpha+1}{2}+\iota}=o(n).
\end{equation}
Applying \eqref{eq:Ptbound} and \eqref{eq:Qtbound} to \eqref{eq:IIabound},
this shows
\begin{align*}
\|\mathrm{(II.a)}\|^2
\lesssim \int_0^t \norm{\bnu^s-\tilde\bnu^s}^2 \d s + o(n).
\end{align*}

For the second term $\mathrm{(II.b)}$, we have analogously
\[\mathrm{(II.b)}
=\int_0^t \bE^s\,\d\bZ^s\]
where now
\[\bE^s=\left(\frac{\eta^{N^{s-}}}{\kappa}-\frac{\bar \eta^s}{\bar\kappa}
\right)\A^s \bPi, \qquad \A^s=[\a_1^s,\ldots,\a_n^s] \in \R^{dk \times n}, \qquad
\a_i^s=\vec(\x_i \otimes f(\x_i^\top \bnu^{s-},\x_i^\top \btheta^*,\beps)).\]
Then we have an analogous decomposition $\|\mathrm{(II.b)}\|^2=P^t+Q^t$. For
$P^t$, similar arguments as above show
\[P^t \lesssim \frac{n^{1-\alpha}}{\binom{n}{\kappa}}
\cdot n^{\alpha-1}\binom{n}{\kappa}
\int_0^t \left(\frac{\eta^{N^{s-}}}{\kappa}-\frac{\bar \eta^s}{\bar\kappa}
\right)^2 \|f(\X\bnu^{s-},\X\btheta^*,\beps)\|^2\,\d s
=o(n),\]
the last statement applying boundedness of $f$ and
\eqref{eq:Poissonprocapprox}. For $Q^t$, applying
$|\frac{\eta^{N^{s-}}}{\kappa}-\frac{\bar \eta^s}{\bar\kappa}| \lesssim 1$
in the bounds for $\bE^s$, arguments identical to the above show
\[\sup_{t \in [0,T]} |Q^t|<n^{\frac{\alpha+1}{2}+\iota}=o(n).\]
Combining these bounds for $\mathrm{(II.a)}$ and $\mathrm{(II.b)}$ shows, on an
event holding a.s.\ for all large $n,d$,
\begin{equation}\label{eq:SGDIIbound}
    \norm{\mathrm{(II)}}^2 \lesssim \int_0^t \norm{\bnu^s-\tilde\bnu^s}^2 \d s
+o(n).
\end{equation}
Then combining \eqref{eq:SGDIbound}, \eqref{eq:SGDIIbound}, \eqref{eq:SGDIIIbound},
and \eqref{eq:SGDIVbound}, we have $\norm{\bnu^t - \tilde\bnu^t}^2
\lesssim \int_0^t \norm{\bnu^s-\tilde\bnu^s}^2 \d s + o(n)$, and hence
by Gr\"onwall's inequality,
\begin{equation}\label{eq:disc-sgd1}
\sup_{t \in [0,T]} \|\bnu^t-\tilde\bnu^t\|^2 \lesssim o(n).
\end{equation}

\paragraph{Step 2: Bounding $\norm{\tilde\bnu^t - \bnu_\delta^t}$}
Next, we compare $\tilde\bnu^t$ to the discretized process $\bnu_\delta^t$: 
\begin{align*}
    \tilde\bnu^t &= \bnu^0 - \int_0^t \frac{\bar\eta^{s}}{\bar\kappa} \X^\top
\diag(\bPi\,\d\bN^s) f(\X\tilde\bnu^{s-},\X\btheta^*,\beps) - \int_0^t
{\bar\eta^{s}} g(\tilde\bnu^{s}) \d s,\\
    \bnu_\delta^t &= \bnu^0 - \int_0^{\flr{t}}
\frac{\bar\eta^{\flr{s}}}{\bar\kappa} \X^\top
\diag(\bPi\,\d\bN^s)f(\X\bnu_\delta^{s-},\X\btheta^*,\beps) - \int_0^{\flr{t}}
\bar\eta^{\flr{s}} g(\bnu_\delta^{s}) \d s.
\end{align*}
We may decompose 
    $\tilde \bnu^t -  \bnu_\delta^t=\mathrm{(I)} + \mathrm{(II)} +
\mathrm{(III)} + \mathrm{(IV)} + \mathrm{(V)}  + \mathrm{(VI)}$
where
\begin{align*}
    \mathrm{(I)} &= - \int_0^{\lfloor t\rfloor} \frac{\bar\eta^{s}}{\bar\kappa}
\X^\top
\diag(\bPi\,\d\bN^s)\big(f(\X\tilde\bnu^{s-},\X\btheta^*,\beps)-f(\X\bnu_\delta^{s-},\X\btheta^*,\beps)\big), \\
    \mathrm{(II)} &= -\int_0^{\lfloor t\rfloor} \frac{\bar\eta^{s} -
\bar\eta^{\lfloor s\rfloor}}{\bar\kappa} \X^\top
\diag(\bPi\,\d\bN^s)f(\X\bnu_\delta^{s-},\X\btheta^*,\beps), \\ \mathrm{(III)}
&= - \int_0^{\lfloor t\rfloor} \bar\eta^{s}\big(g(\tilde\bnu^{s}) -
g(\bnu_\delta^{s})\big)\d s,\\
\mathrm{(IV)}&= - \int_0^{\lfloor t\rfloor} \big(\bar\eta^{s} -
\bar\eta^{\lfloor s\rfloor}\big) g(\bnu_\delta^{s})\d s,\\
\mathrm{(V)} &= - \int_{\lfloor t\rfloor}^{t} \frac{\bar\eta^{s}}{\bar\kappa}
\X^\top \diag(\bPi\,\d\bN^s)f(\X\tilde\bnu^{s-},\X\btheta^*,\beps),\\
\mathrm{(VI)}&=- \int_{\lfloor t\rfloor}^{t} \bar\eta^{s} g(\tilde\bnu^{s})\d s.
\end{align*}
Note that $t-\flr{t} \lesssim \delta$, and
by the Lipschitz continuity of $\bar\eta$ in Assumption \ref{asp:SGD},
for any $p \geq 1$ and $\delta \in (0,1)$,
\begin{equation}\label{eq:etadiscrapprox}
\int_0^t |\bar \eta^s-\bar\eta^{\flr{s}}|^p \d s \lesssim \delta.
\end{equation}
Then on the event $\sup_{t \in [0,T]} \|g(\tilde \bnu^t)\|
\leq C\sqrt{n}$ which holds a.s.\ for all large $n,d$, we have
\begin{align*}
    \norm{\mathrm{(III)}} \lesssim \int_0^t
\norm{\tilde\bnu^{s} - \bnu_\delta^{s}}\d s,
\qquad
    \norm{\mathrm{(IV)}}^2\lesssim \delta\left(\int_0^t
\norm{\bnu_\delta^s - \tilde\bnu^s }^2\d s + n\right),
\qquad
\norm{\mathrm{(VI)}}\lesssim \delta\sqrt{n}.
\end{align*}
For the remaining stochastic integration terms, applying
\eqref{eq:etadiscrapprox} and similar arguments as in Step 1 (omitted here for
brevity), on an event holding a.s.\ for all large $n,d$,
\begin{align*}
\norm{\mathrm{(I)}}^2\lesssim \int_0^t \norm{\tilde\bnu^s-\bnu_\delta^s}^2
\d s + o(n),
\qquad \norm{\mathrm{(II)}}^2 \lesssim  \delta n + o(n),
\qquad \norm{\mathrm{(V)}}^2\lesssim \delta n + o(n).
\end{align*}
Combining this gives
$\norm{\bnu_\delta^t - \tilde\bnu^t}^2\lesssim \int_0^t
\norm{\bnu_\delta^s-\tilde\bnu^s}^2 \d s + \delta n+o(n)$,
so Gr\"onwall's inequality implies
\begin{equation}\label{eq:disc-sgd2}
\sup_{t \in [0,T]}
\norm{\bnu_\delta^t - \tilde\bnu^t}^2 \lesssim \delta n+o(n).
\end{equation}
The lemma follows from combining these bounds \eqref{eq:disc-sgd1} and
\eqref{eq:disc-sgd2}.
\end{proof}

\subsection{Discretization of SME}
We now apply a similar analysis to the Stochastic Modified Equation. The
SME process \eqref{eq:SME} may be written concisely as
\begin{align*}
\btheta^t &= \btheta^0-\int_0^t \bar\eta^s \prn{\X^\top
f(\X\btheta^s,\X\btheta^*,\beps) + g(\btheta^s)} \d s + \int_0^t
\frac{\bar\eta^s}{\sqrt{\bar \kappa}} \X^\top
\diag(\d\B^s)f(\X\btheta^{s},\X\btheta^*,\beps)
\end{align*}
where $\B^s$ is a standard $n$-dimensional Brownian motion.
We compare this to the discretized process $\{\bnu_\delta^t\}_{t \geq 0}$
defined by
\begin{equation}\label{eq:deltaSME}
\bnu_\delta^t = \btheta^0-\int_0^{\flr{t}} \bar\eta^{\flr{s}}
\prn{\X^\top f(\X\bnu_\delta^s,\X\btheta^*,\beps) +
g(\bnu_\delta^s)} \d s + \int_0^{\flr{t}}
\frac{\bar\eta^{\flr{s}}}{\sqrt{\bar \kappa}} \X^\top \diag(\d\B^s)
f(\X\bnu_\delta^s,\X\btheta^*,\beps).
\end{equation}
One may again check that $\{\bnu_\delta^t\}_{t \geq 0}$ is precisely a
piecewise-constant embedding of the discrete-time dynamics
\eqref{eq:discretedynamics} in the SME setting of $\z_\delta^k \sim
\cN(\delta\bar\kappa\,\1_n,\delta\bar\kappa\,\Id_n)$,
\[\bnu_\delta^t=\btheta_\delta^k \text{ if } \flr{t}=k,\]
upon identifying $\z_\delta^k=\int_{k\delta}^{(k+1)\delta}
(\bar\kappa \1_n \d t+\sqrt{\bar\kappa}\,\d\B^t)$.

The following lemma is analogous to Lemma \ref{lem:disc-sgd}.

\begin{lem}\label{lem:disc-sme}
For any fixed $T>0$, $\delta \in (0,1)$, and some constant $C \equiv
C(T)>0$ not depending on $\delta$, almost surely as $n,d \to \infty$,
\begin{equation}
\limsup_{n,d\to \infty} \sup_{t\in[0,T]}
\frac{\norm{\btheta^t-\bnu_\delta^t}^2}{d} \leq C\delta.
\end{equation}
\end{lem}

\begin{proof}
We decompose the difference $\btheta^t - \bnu_\delta^t$ into six terms:
\begin{align*}
\btheta^t - \bnu_\delta^t
&=-\underbrace{\int_0^{\flr{t}}\bar\eta^{\flr{s}} \prn{\X^\top
(f(\X\btheta^s,\X\btheta^*,\beps)-f(\X\bnu_\delta^s,\X\btheta^*,\beps)) +
g(\btheta^s)-g(\bnu_\delta^s)} \d s}_{(\mathrm{I})}\\
&\qquad- \underbrace{\int_{\flr{t}}^t \bar\eta^{\flr{s}} \prn{\X^\top
f(\X\btheta^s,\X\btheta^*,\beps) + g(\btheta^s)} \d s}_{(\mathrm{II})}  \nonumber\\
&\qquad+ \underbrace{\int_0^{\flr{t}}\frac{\bar\eta^{\flr{s}}}{\sqrt{\bar \kappa}} \X^\top
\diag(\d\B^s)\big(f(\X\btheta^s,\X\btheta^*,\beps)-f(\X\bnu_\delta^s,\X\btheta^*,\beps)\big)}_{(\mathrm{III})}\\
&\qquad+ \underbrace{\int_{\flr{t}}^t\frac{\bar\eta^{\flr{s}}}{\sqrt{\bar \kappa}} \X^\top
\diag(\d\B^s)f(\X\btheta^s,\X\btheta^*,\beps)}_{(\mathrm{IV})}\\
&\qquad- \underbrace{\int_0^t \prn{\bar\eta^{{s}} - \bar\eta^{\flr{s}}} \prn{\X^\top
f(\X\btheta^s,\X\btheta^*,\beps) + g(\btheta^s)} \d s}_{\mathrm{(V)}}\\
&\qquad+ \underbrace{\int_0^t
\frac{\bar\eta^{{s}} - \bar\eta^{\flr{s}}}{\sqrt{\bar \kappa}} \X^\top
\diag(\d\B^s) f(\X\btheta^s,\X\btheta^*,\beps)}_{\mathrm{(VI)}}\, ,
\end{align*}
% This gives the bound  $\norm{\btheta^t - \btheta_\eps^t}^2\leq 4 (\norm{(\mathrm{I})}^2+\norm{(\mathrm{II})}^2+\norm{(\mathrm{III})}^2+\norm{(\mathrm{IV})}^2)$.

\paragraph{Bounding the martingale terms (III) and (IV)}

Identify $\R^{d \times k} \equiv \R^m$ for $m=dk$, and denote
\[\M^t=\int_0^t \X^\top \diag(\d\B^s)\bD^s \in \R^{dk},
\qquad \bD^s=\frac{\bar\eta^{\flr{s}}}{\sqrt{\bar\kappa}}
\big(f(\X\btheta^s,\X\btheta^*,\beps)
-f(\X\bnu_\delta^s,\X\btheta^*,\beps)\big).\]
Note that this may be written as
\[\d\M^t=\A^t \d\B^t,\qquad
\A^t=[\a_1^t,\ldots,\a_n^t] \in \R^{dk \times n},\]
\[\a_i^t=\frac{\bar\eta^{\flr{t}}}{\sqrt{\bar\kappa}}
\vec(\x_i \otimes [f(\x_i^\top \btheta^s,\x_i^\top \btheta^*,\beps)
-f(\x_i^\top \bnu_\delta^s,\x_i^\top \btheta^*,\beps)]).\]
By It\^o's formula,
\begin{align*}
&\norm{\M^t}^2=\int_0^t \|\A^s\|^2 \d s
+\int_0^t 2{\M^s}^\top \A^s \d\B^s\\
&=\underbrace{\int_0^t \frac{(\bar\eta^{\flr{s}})^2}{\bar\kappa}
\left(\sum_{i=1}^n  \|\x_i\|^2 \cdot \|f(\x_i^\top \btheta^s,\x_i^\top
\btheta^*,\beps)-f(\x_i^\top \bnu_\delta^s,\x_i^\top
\btheta^*,\beps)\|^2\right)\d s}_{:=P^t}
+\underbrace{\int_0^s 2{\M^s}^\top \X^\top \diag(\d\B^s)\bD_s}_{:=Q^t}.
\end{align*}
By the Lipschitz continuity of $f$, on an event holding a.s.\ for all large
$n,d$, we have
\[P^t \lesssim \int_0^t \|f(\X \btheta^s,\X \btheta^*,\beps)-f(\X
\bnu_\delta^s,\X \btheta^*,\beps)\|^2 \d s
\lesssim \int_0^t \|\btheta^s-\bnu_\delta^s\|^2 \d s.\]
By the Burkholder-Davis-Gundy inequality
\cite[Theorem 2.11]{kunita2004stochastic}, for any $p \geq 1$,
\begin{align}
    \E \sup_{t\in[0,T]} (Q^t)^{2p} \leq C_p\,
\E\left(\int_0^t \|{\M^s}^\top \A^s\|^2 \d s\right)^p.
\end{align}
Applying $\|\A^s\|_\op \lesssim \|\X\|_\op\|f\|_\infty \lesssim 1$, this implies
\[\E\sup_{t\in[0,T]} (Q^t)^{2p} \lesssim
\int_0^t \E \norm{\M^s}^{2p} \d s.\]
Then applying It\^o's formula to $F(\x)=\norm{\x}^{2p}$, since $\M^t$ has no drift
coefficient and
$\nabla^2 F(\x)=2p\norm{\x}^{2p-2}\Id+4p(p-1)\norm{\x}^{2p-4}\x\x^\top$,
we have
\begin{align}
\frac{\d}{\d t} \E\norm{\M^t}^{2p} &\lesssim
\E\brk{\norm{\M^t}^{2p-2}\|\A^t\|_\sF^2} +
\E\brk{\norm{\M^t}^{2p-4}\norm{{\A^t}\M^t}^2}\nonumber\\
    &\lesssim n\,\E\norm{\M^t}^{2p-2} \lesssim
n(\E\norm{\M^t}^{2p})^{\frac{p-1}{p}}
\lesssim \E\norm{\M^t}^{2p}+n^p.
\end{align}
Thus by Gr\"onwall's inequality,
$\E\norm{\M^t}^{2p} \lesssim n^p$, implying that
$\E \sup_{t\in[0,T]} (Q^t)^{2p} \lesssim n^p$.
Then Markov's inequality shows that
for any fixed and sufficiently small $\iota>0$, a.s.\ for all large $n,d$,
\begin{align}
\sup_{t\in[0,T]} |Q^t|<n^{\frac{1}{2}+\iota}=o(n).
\end{align}
Applying these bounds gives
\[\|\M^t\|^2 \lesssim \int_0^t \|\btheta_s-\bnu_\delta^s\|^2 \d s+o(n).\]
Then
\begin{equation}\label{eq:SMEIIIbound}
\norm{(\mathrm{III})}^2=\norm{\M^{\flr{t}}}^2 \lesssim \int_0^t
\norm{\btheta^s-\bnu_\delta^s}^2 \d s + o(n).
\end{equation}
Similarly, a.s.\ for all large $n,d$,
\begin{equation}\label{eq:SMEIVbound}
\norm{(\mathrm{IV})}^2 \lesssim \delta n + o(n).
\end{equation}\\

\paragraph{Bounding the drift terms (I) and (II)}

Write
\begin{align*}
\btheta^t &=\btheta^0-\int_0^t \bar\eta^s \prn{\X^\top
f(\X\btheta^s,\X\btheta^*,\beps) + g(\btheta^s)} \d s + 
\M^t,\\
\M^t&=\int_0^t \frac{\bar\eta^s}{\sqrt{\bar \kappa}} \X^\top \diag(\d\B^s)
f(\X\btheta^{s},\X\btheta^*,\beps).
\end{align*}
The same argument as above shows, on an event holding a.s.\ for all large $n,d$,
for every $t \in [0,T]$,
\begin{align*}
\norm{\M^t}^2 \lesssim \int_0^t
\norm{f(\X\btheta^s,\X\btheta^*,\beps)}^2 \d s + o(n) \lesssim n.
\end{align*}
Then by the boundedness of $f(\cdot)$ and Lipschitz continuity of $g(\cdot)$, also
$\|\btheta^t\|^2 \lesssim \int_0^t \|\btheta^s\|^2 \d s+n$, so Gr\"onwall's
inequality gives
\[\sup_{t \in [0,T]} \|\btheta^t\|^2 \lesssim n.\]
This implies, by the Lipschitz continuity of $f(\cdot)$ and $g(\cdot)$, on an event
holding a.s.\ for all large $n,d$,
\begin{equation}\label{eq:SMEIbound}
\norm{(\mathrm{I})} \lesssim \int_0^t \norm{\btheta^s-\bnu_\delta^s}
\d s, \qquad \norm{(\mathrm{II})} \lesssim \delta\sqrt{n}.
\end{equation}

\paragraph{Bounding the discretization terms (V) and (VI)}

Applying \eqref{eq:etadiscrapprox} and similar arguments as for
the preceding bounds for $\mathrm{(II)}$ and $\mathrm{(IV)}$, on an event
holding a.s.\ for all large $n,d$, we have
\begin{equation}\label{eq:SMEVbound}
    \norm{\mathrm{(V)}}^2 \lesssim \delta n,
\qquad
    \norm{\mathrm{(VI)}}^2 \lesssim \delta n+o(n).
\end{equation}

Combining \eqref{eq:SMEIIIbound}, \eqref{eq:SMEIVbound},
\eqref{eq:SMEIbound}, and \eqref{eq:SMEVbound},
\begin{align*}
    \norm{\btheta^t - \bnu_\delta^t}^2&\lesssim 
    \int_0^t \norm{\btheta^s-\bnu_\delta^s}^2 \d s + \delta n + o(n),
\end{align*}
and thus Gr\"onwall's inequality yields $\sup_{t \in [0,T]}
\norm{\btheta^t - \bnu_\delta^t}^2 \lesssim \delta n+o(n)$, implying the lemma.
\end{proof}

\begin{proof}[Proof of Theorem \ref{thm:DMFT-limit}(a)]
Let $\btheta^t=\bar\btheta^{\lfloor tn^{1-\alpha} \rfloor}$ be the time-rescaled
SGD process in \eqref{eq:SGDrescaled} (where here $\flr{\cdot}$ denotes the
usual integer part), and let $\bnu^t=\bar\btheta^{N^t}$ be
the Poissonized SGD process in \eqref{eq:SGDPoissonized}. For any $\delta>0$,
let $\bnu_\delta^t$ be the discretized approximation in \eqref{eq:deltaSGD}.
Applying an induction over discrete time increments of $\delta$ to
\eqref{eq:deltaSGD}, we see that $\{\bnu_\delta^t\}_{t \geq 0}$
is piecewise-constant with jumps at $\delta\Z_+$, taking values
\begin{equation}\label{eq:nudeltaembedding}
\bnu_\delta^t=\btheta_\delta^k \text{ for all } t \in [k\delta,(k+1)\delta)
\end{equation}
where $\btheta_\delta^k$ are the iterations of the discrete-time dynamics
\eqref{eq:discretedynamics} defined using the Poisson variables
$\z_\delta^k=\bPi\bP_\delta^k$
and $\bP_\delta^k=\int_{k\delta}^{(k+1)\delta} \d\bN^s$.

Fix any $T>0$. For each $t \in [0,T]$, set
$\tau(t)=\inf\{\tau \geq 0:N^\tau \geq \lfloor tn^{1-\alpha} \rfloor\}$,
so that
\[\btheta^t=\bnu^{\tau(t)}.\]
Note that the bound \eqref{eq:Poissontimeapprox} implies that for some
$\iota>0$, on an event holding a.s.\ for all large $n,d$,
\begin{equation}\label{eq:tauapprox}
\sup_{t \in [0,T]} |\tau(t)-t| \lesssim n^{-\iota}.
\end{equation}
Fix any $0 \leq t_1 \leq \ldots \leq t_m \leq T$, and fix any
$\delta \in (0,1)$ such that none of $t_1,\ldots,t_m$ is an integer multiple of
$\delta$. Then a.s.\ for all large $n,d$, the piecewise-constant nature of
$\{\bnu_\delta^t\}_{t \geq 0}$ and \eqref{eq:tauapprox} imply that
$\bnu_\delta^{\tau(t_i)}=\bnu_\delta^{t_i}$ for each $i=1,\ldots,m$. Then
by Lemma \ref{lem:disc-sgd}, a.s.\ for all large $n,d$,
\[\sum_{i=1}^m \frac{\|\btheta^{t_i}-\bnu_\delta^{t_i}\|^2}{d}
=\sum_{i=1}^m \frac{\|\bnu^{\tau(t_i)}-\bnu_\delta^{\tau(t_i)}\|^2}{d}
\lesssim \delta.\]
This implies that
\begin{equation}\label{eq:W2compare1}
\limsup_{n,d \to \infty} W_2\left(\frac{1}{d}\sum_{j=1}^d
\delta_{(\theta_j^{t_1},\ldots,\theta_j^{t_m},\theta_j^*)},\,
\frac{1}{d}\sum_{j=1}^d
\delta_{(\nu_{\delta,j}^{t_1},\ldots,\nu_{\delta,j}^{t_m},\theta_j^*)}\right)^2
\lesssim \delta.
\end{equation}
Let $k_1,\ldots,k_m$ be the integers for which $t_i \in
[k_i\delta,(k_i+1)\delta)$ for each $i=1,\ldots,m$. 
Then by
\eqref{eq:nudeltaembedding} and Lemma \ref{lem:finite_dim_converge},
\[\limsup_{n,d \to \infty} W_2\left(\frac{1}{d}\sum_{j=1}^d
\delta_{(\nu_{\delta,j}^{t_1},\ldots,\nu_{\delta,j}^{t_m},\theta_j^*)},
\,\Law(\theta_\delta^{k_1},\ldots,\theta_\delta^{k_m},\theta^*)\right)=0\]
where the right side denotes the joint law under the discrete DMFT system of
Section \ref{sec:discreteDMFT}. Recalling the embedding 
\eqref{eq:embeddings} of this discrete system, where
$\bar\theta_\delta^{t_i}=\theta_\delta^{k_i}$ for each $i=1,\ldots,m$, this is
equivalently written as
\begin{equation}\label{eq:W2compare2}
\limsup_{n,d \to \infty} W_2\left(\frac{1}{d}\sum_{j=1}^d
\delta_{(\nu_{\delta,j}^{t_1},\ldots,\nu_{\delta,j}^{t_m},\theta_j^*)},
\,\Law(\bar\theta_\delta^{t_1},\ldots,\bar\theta_\delta^{t_m},\theta^*)\right)
=0.
\end{equation}
Finally, by Lemma \ref{lem:dmft_discretization_error}(c),
\begin{equation}\label{eq:W2compare3}
\lim_{\delta \to 0} W_2\Big(
\Law(\bar\theta_\delta^{t_1},\ldots,\bar\theta_\delta^{t_m},\theta^*),\,
\Law(\theta^{t_1},\ldots,\theta^{t_m},\theta^*)\Big)=0
\end{equation}
where the right side denotes the joint law under the continuous DMFT system
defined via the fixed point of
Theorem \ref{thm:DMFT-limit}. Combining \eqref{eq:W2compare1},
\eqref{eq:W2compare2}, and \eqref{eq:W2compare3} and taking the limit $n,d \to
\infty$ followed by $\delta \to 0$, we obtain \eqref{eq:DMFT-limit-theta}.
The argument for \eqref{eq:DMFT-limit-xi} is the same.
\end{proof}

\begin{proof}[Proof of Theorem \ref{thm:DMFT-limit}(b)]
Let $\{\btheta^t\}_{t \geq 0}$ be the SME process \eqref{eq:SME}, and let
$\bnu_\delta^t$ be its discretized approximation \eqref{eq:deltaSME}.
Again by induction over discrete time increments of $\delta$, we see that
$\{\bnu_\delta^t\}_{t \geq 0}$
is piecewise-constant with jumps at $\delta\Z_+$, taking values
\begin{equation}\label{eq:nudeltaembedding-SME}
\bnu_\delta^t=\btheta_\delta^k \text{ for all } t \in [k\delta,(k+1)\delta)
\end{equation}
where $\btheta_\delta^k$ are the iterations of the discrete-time SME dynamics
\eqref{eq:discretedynamics} defined using the Gaussian variables
$\z_\delta^k=\int_{k\delta}^{(k+1)\delta} (\bar\kappa\1_n\d t
+\sqrt{\bar\kappa}\,\d\B^t)$. 
By Lemma \ref{lem:disc-sme}, a.s.\ for all large $n,d$,
\[\sum_{i=1}^m \frac{\|\btheta^{t_i}-\bnu_\delta^{t_i}\|^2}{d} \lesssim
\delta.\]
The proof is then completed as above, using 
Lemmas \ref{lem:finite_dim_converge} and
\ref{lem:dmft_discretization_error}(c) for the SME.
\end{proof}

%\section{Applications: Multi-Index Model}
%In this section, we apply Theorem \lw{TBD} to illustrate DMFT characterization of the training dynamic on a multi-index model using the multi-pass SGD algorithm.
%A multi-index function $f\colon \R^d\to \R$ is a function satisfies $f(x)=g(Wx)$, where $g\colon \R^p\to \R$ is a low-dimensional link function, and $W\in\R^{p\times d}$ is a row-orthogonal matrix. We are interested in the learning problems in a teacher-student setting, where the teacher model is given by a multi-index function while the student model is trained via a two-layer neural network.
%To be precise, the dataset $\mathcal{D} = \{x_i, y_i\}_{i\in[n]}$ contains $n$
%i.i.d. isotropic Gaussian vector $x_i\sim\cN(0,I_d/d)$, and the corresponding label $y_i$ generated by the ground-truth multi-index function, that is, $y_i = f^\star(x_i) = g^\star(W^\star x_i)$.
%Then we put the data $\mathcal{D}$ inside a two-layer neural network with $k$ neurons, which can be written as
%\[
%f(x) = a^\top \sigma(Wx) = \sum_{i=1}^k a_i \sigma(w_i^\top x)\, .
%\]
%Here $W\in\R^{k\times d}$, each row vector of $W$ represents the first layer weights of one neuron, while $a\in\R^k$ is the second layer weights.
%We are trying to understand the dynamics of minibatch SGD learning the first-layer weight $W$ on the empirical MSE loss given by
%\begin{align}
%    \mathcal{R}_n = \sum_{i=1}^n \prn{y_i - f(x_i)}^2 = \sum_{i=1}^n \prn{g^\star(W^\star x_i) - a^\top \sigma(Wx_i)}^2\,.
%\end{align}
%With a penalty term $\iota \|W\|_F^2/2$, the SGD algorithm is given by
%\[
%W^{k+1} - W^k = -\eta \left[ x_{I_k}  \partial_1 L( W^k x_{I_k}, W^\star x_{I_k}) + \frac{\iota}{n} W^k \right]\,.
%\]
%Since $L( W^k x_{i}, W^\star x_{i}) = \prn{g^\star(W^\star x_i) - a^\top \sigma(W^kx_i)}^2$, we can get
%\[f( W^k x_{i}, W^\star x_{i})=\partial_1 L( W^k x_{i}, W^\star x_{i}) = 2\prn{a^\top \sigma(W^k x_i) - g^\star(W^\star x_i)}\cdot \prn{a\odot \sigma'(W^k x_i)}\, . \]
%Using the DMFT system we derived from the last section with the planted signal inside, we can write it down specifically as 
%\begin{subequations}
%\begin{align}
%W^t-W^0&=-\int_0^t\prn{\Gamma_{s}W^{s}+ \eta \iota W^{s} + \alpha\eta \Rf^{0,s}(W^{[0,s]}) + \alpha\eta R_f(s,\ast)w^\ast} \d s+u^t \, , & & u^t \sim \mathsf{GP}(0, \Cf) \, , \label{eq:def-cont-theta-2} \\
%\xi^t&=-2\eta\int_0^t \Rt^{s,t}\prn{a^\top \sigma(\xi^s)- g^\star(w^\ast)}\prn{a\odot \sigma'(\xi^s)}\prn{\d s+\d Z^s}+w^t \, , & & w^t \sim \mathsf{GP}(0, C_\theta) \, , \label{eq:def-cont-xi-2}\\
%\d M^t &= \prn{a^\top \sigma(\xi^s)- g^\star(w^\ast)}\prn{a\odot \sigma'(\xi^s)}\prn{\d t + \d Z^t}\, , & & 0\leq t < \infty, \label{eq:def-mart-fxi-2}\\
%\Rt^{s,t} & = \E \left[\frac{\partial W^t}{\partial u^s}\right] \, ,& & 0 \leq s \leq t < \infty\, , \label{eq:def-cont-R-t-2}  \\
%\Gamma_t & = \alpha\eta \E\brk{\prn{a\odot \sigma'(\xi^s)}\prn{a\odot \sigma'(\xi^s)}^\top + \prn{a^\top \sigma(\xi^s)- g^\star(w^\ast)} \diag\prn{a\odot \sigma''(\xi^s)} }\, , & & 0\leq t < \infty\, ,\label{eq:def-cont-Gamma-2}\\
%\Rf^{s,t}(W^{[s,t]}) &= \E_{\xi^{[s,t]}, Z^{[s,t]}}\brk{r^{s,t}\prn{W^{[s,t]}}}\, , & &0 \leq s\leq t <\infty \, , \label{eq:def-cont-R-h-2} \\
%R_f(t,\ast) & = \E \left[\frac{\partial f(\xi^t, w^\ast)}{\partial w^\ast}\right] \, , & &0 \leq t <\infty\, , \label{eq:def-cont-R-f-ast-2} \\
%C_\theta(t,s) & = \E \left[W^t {W^s}^\top\right]\, , & & 0 \leq s , t < \infty\, ,  \label{eq:def-cont-C-t-2} \\
%\Cf^{t,s} & = \alpha \eta^2 \E\brk{M^t {M^s}^\top} \, ,  & & 0 \leq s \leq t < \infty\, . \label{eq:def-cont-C-h-2}
%\end{align}
%\end{subequations}
%where the partial derivatives
%are defined via \begin{subequations}
%\begin{align}
%\frac{\d}{\d t}\left(\frac{\partial W^{t}}{\partial u^s}\right)&=-\left(\Gamma_t+\eta \iota I_k \right)\cdot\frac{\partial W^{t}}{\partial u^s}-\alpha\eta\cdot \Rf^{s,t}\prn{\frac{\partial W^{[s,t]}}{\partial u^s}} \, , & & 0 \leq s \leq t < \infty\, , \label{eq:def-cont-deriv-t-2}  \\
%r^{s,t}(x^{[s,t]}) & = -\eta \prn{\prn{a\odot \sigma'(\xi^t)}\prn{a\odot \sigma'(\xi^t)}^\top + \prn{a^\top \sigma(\xi^t)- g^\star(w^\ast)} \diag\prn{a\odot \sigma''(\xi^t)}}\int_s^t \Rt(t,r) \nonumber\\
%&\cdot\prn{r^{s,r}(x^{[s,r]}) + \prn{\prn{a\odot \sigma'(\xi^s)}\prn{a\odot \sigma'(\xi^s)}^\top + \prn{a^\top \sigma(\xi^s)- g^\star(w^\ast)} \diag\prn{a\odot \sigma''(\xi^s)}}x^r}(\d r + \d Z^r) \, , & &0 \leq s\leq t <\infty\, , \label{eq:def-cont-resp-xi-2}\\
%\frac{\partial f(\xi^t, w^\ast)}{\partial w^\ast}&=-\prn{a\odot \sigma'(\xi^t)} \nabla g^\star(w^\ast)^\top
%-\eta\prn{\prn{a\odot \sigma'(\xi^s)}\prn{a\odot \sigma'(\xi^s)}^\top + \prn{a^\top \sigma(\xi^s)- g^\star(w^\ast)} \diag\prn{a\odot \sigma''(\xi^s)}}\nonumber\\
%&\quad\cdot\int_0^t \Rt(t,s)\frac{\partial f(\xi^s, w^\ast)}{\partial w^*}\prn{\d s+\d Z^s}
%\, . & & 0 \leq s < t < \infty\, . \label{eq:def-cont-deriv-xi-ast-2}
%\end{align}
%\end{subequations}

\appendix

\section{Burkholder-Davis-Gundy inequality}

\begin{lem}\label{lem:burkholder}
Let $\{\mathcal{F}_t\}_{t \geq 0}$ be a complete right-continuous filtration on
the probability space $(\Omega,\cF,\P)$. Let
$\{X^t\}_{t \in [0,T]}$ be a $\cF_t$-predictable process taking
values in any Euclidean space $\R^m$. If $\{z^t\}_{t \in [0,T]}$
is the $\cF_t$-adapted diffusion process \eqref{eq:gaussian-z},
then for any $p \geq 1$,
there exists a constant $C_p>0$ such that for all $t \geq 0$,
\[\E \sup_{s \in [0,t]} \left\|\int_0^s X^r\,\d z^r\right\|_2^p \leq
C_p\,\E\left(\bar\kappa \int_0^t \|X^r\|_2\,\d r\right)^p
+C_p\,\E\left(\bar\kappa \int_0^t \|X^r\|_2^2\,\d r\right)^{p/2}.\]
If $\{z^t\}_{t \in [0,T]}$ is the $\cF_t$-adapted Poisson process 
\eqref{eq:poisson-z}, then for any $p \geq 1$,
there exists a constant $C_p>0$ such that for all $t \geq 0$,
\[\E \sup_{s \in [0,t]} \left\|\int_0^s X^r\,\d z^r\right\|_2^p \leq
C_p\,\E\left(\bar\kappa \int_0^t \|X^r\|_2\,\d r\right)^p
+C_p\,\E\left(\bar\kappa \int_0^t \|X^r\|_2^2\,\d r\right)^{p/2}
+C_p\bar\kappa\,\E\int_0^t \|X^r\|_2^p\d r.\]
\end{lem}
\begin{proof}
This follows from \cite[Thm 2.11]{kunita2004stochastic} applied to
$\int_0^s X^r \d z^r
=\int_0^s X^r (\d z^r-\bar \kappa\,\d r)+\int_0^s X^r \bar\kappa\,\d r$,
where $\d z^r-\bar \kappa\,\d r$ is either $\sqrt{\bar\kappa}\,\d b^r$ for a
standard Brownian motion $\{b^t\}_{t \geq 0}$, or a compensated Poisson 
process with rate $\bar\kappa$.
\end{proof}

\section*{Acknowledgments}

We would like to thank Sinho Chewi, Alex Damian, Bruno Loureiro, and Theodor Misiakiewicz for helpful discussions. This research was supported in part by NSF DMS-2142476 and a Sloan Research Fellowship.

\bibliographystyle{alpha}
\bibliography{references}

\end{document}
